\documentclass[10pt]{article}
% Layout of a NeurIPS/ICML submission: single column, 10pt Times, 5.5in x 9in text block.
% Swapping in a venue's style file is a one-line change.
\usepackage[letterpaper,textwidth=5.5in,textheight=9in,top=1in,headsep=0pt,headheight=0pt,footskip=30pt]{geometry}
\usepackage[T1]{fontenc}
\usepackage{times}
\usepackage{microtype}
\usepackage{amsmath,amssymb,amsthm,mathtools}
\usepackage{thm-restate}
\usepackage{booktabs,array}
\usepackage{placeins}
\usepackage{graphicx}
\usepackage[font=small,labelfont=bf]{caption}
\usepackage{tikz}
\usepackage{framed}
\usepackage{needspace}
\usetikzlibrary{positioning,arrows.meta,calc,matrix,fit}
\usepackage[round,sort]{natbib}
\renewcommand{\bibfont}{\small\raggedright}
\renewcommand{\bibsection}{\section*{References}\addcontentsline{toc}{section}{References}}
\usepackage[colorlinks,linkcolor=blue!45!black,citecolor=blue!45!black,urlcolor=blue!45!black]{hyperref}
\usepackage[capitalise,noabbrev,nameinlink]{cleveref}
\usepackage{enumitem}
\setlist{nosep,leftmargin=1.4em}
\setlength{\textfloatsep}{10pt plus 2pt minus 2pt}
\renewcommand{\topfraction}{0.85}
\renewcommand{\textfraction}{0.1}

% TeX Live 2026 (LaTeX 2026-06-01) makes a theorem that shares a counter an alias of it (\newcounteralias),
% and thmtools, loaded by thm-restate, would make the same alias again, an error. Leave the alias to LaTeX;
% on older TeX Live, without \newcounteralias, nothing changes.
\makeatletter
\@ifundefined{newcounteralias}{}{%
  \renewcommand\thmt@autorefsetup{\@xa\def\csname\thmt@envname autorefname\@xa\endcsname\@xa{\thmt@thmname}}%
}
\makeatother
\newtheorem{theorem}{Theorem}[section]
\newtheorem{proposition}[theorem]{Proposition}
\newtheorem{corollary}[theorem]{Corollary}
\newtheorem{lemma}[theorem]{Lemma}
\newtheorem{conjecture}{Conjecture}
\theoremstyle{definition}
\newtheorem{definition}[theorem]{Definition}
\newtheorem*{definition*}{Definition}
\newtheorem{example}[theorem]{Example}
\newtheorem*{example*}{Example}
\theoremstyle{remark}
\newtheorem{remark}[theorem]{Remark}
\newtheorem*{remark*}{Remark}
\crefname{equation}{}{}
\crefname{conjecture}{Conjecture}{Conjectures}
\Crefname{conjecture}{Conjecture}{Conjectures}

\newcommand{\R}{\mathbb{R}}
\newcommand{\Wread}{W_{\mathrm{read}}}
\newcommand{\Wwrite}{W_{\mathrm{write}}}
\newcommand{\dm}{d_{\mathrm{model}}}
\newcommand{\dhid}{d_{\mathrm{hidden}}}
\newcommand{\GL}{\mathrm{GL}}
\newcommand{\Ort}{\mathrm{O}}
\newcommand{\CO}{\mathrm{CO}}
\newcommand{\Aff}{\mathrm{Aff}}
% The Frechet derivative, upright to set it apart from the readout D.
\newcommand{\Der}{\mathrm{D}}
\newcommand{\G}{\mathcal{G}}
\newcommand{\norm}[1]{\lVert #1\rVert}
\newcommand{\para}[1]{\smallskip\noindent\textbf{#1}\hspace{0.5em}}
% The one-sentence summary at the start of a main section.
\newcommand{\takeaway}[1]{\noindent\textit{#1}\par\smallskip}
% A link from a result in the main text to its proof in Appendix A.
% A link target at a part of a statement, and references that link name, number and part to it.
\newcommand{\ptlabel}[1]{\phantomsection\label{#1}}
\DeclareRobustCommand{\cpref}[2]{\hyperref[#1:#2]{\cref*{#1}(#2)}}
\DeclareRobustCommand{\Cpref}[2]{\hyperref[#1:#2]{\Cref*{#1}(#2)}}
\DeclareRobustCommand{\rpref}[2]{\hyperref[#1:#2]{\ref*{#1}(#2)}}
\DeclareRobustCommand{\ptlink}[2]{\hyperref[#1:#2]{(#2)}}
% No page break between a statement and the line that follows it (its "Proof in" line, or its proof).
\makeatletter
\@ifundefined{AddToHook}{}{\AddToHook{env/restatable/begin}{\@endparpenalty\@M}}
\makeatother
\newcommand{\proofin}[1]{\par\vspace{-0.4\baselineskip}\noindent\hfill{\small\itshape Proof in \cref{#1}.}\par\smallskip}
% A main-text result that is a special case of, or follows from, a result in the appendix.
\newcommand{\specialof}[2]{\par\vspace{-0.4\baselineskip}\noindent\hfill{\small\itshape A special case of \cref{#1}, proved in \cref{#2}.}\par\smallskip}
\newcommand{\followsfrom}[2]{\par\vspace{-0.4\baselineskip}\noindent\hfill{\small\itshape Follows from \cref{#1}, proved in \cref{#2}.}\par\smallskip}
% A part of a statement: a link target (if a label is given), a bold number in parentheses and an optional bold name.
\newcommand{\stpart}[3]{\ifx\relax#1\relax\else\ptlabel{#1:#2}\fi{\upshape\bfseries\boldmath(#2)\ifx\relax#3\relax\else\ #3.\fi}}
% A lemma and its proof inside another proof: indented, with a thin gray rule on the left.
\newenvironment{nested}{\par\nopagebreak\def\FrameCommand{\hspace{0.5em}{\color{black!30}\vrule width 0.6pt}\hspace{1em}}\MakeFramed{\advance\hsize-\width\FrameRestore}}{\endMakeFramed}


% Statements of the first paper, Karbevski and Mijoski [KM], are cited by their numbers there.
% The numbers are set here, in one place.
\newcommand{\KMset}[2]{\expandafter\gdef\csname KM@#1\endcsname{#2}}
\newcommand{\KMn}[1]{\csname KM@#1\endcsname}
\newcommand{\KM}[2]{[KM, #1~\KMn{#2}]}
\newcommand{\KMp}[3]{[KM, #1~\KMn{#2}(#3)]}
\newcommand{\KMeq}[1]{[KM, (\KMn{#1})]}
\KMset{thm:resdepth}{4.8}
\KMset{thm:unlock}{4.3}
\KMset{thm:oneblock}{4.7}
\KMset{thm:onelayer}{B.1}
\KMset{thm:gaugefull}{A.5}
\KMset{thm:count}{A.17}
\KMset{lem:smallwrites}{A.27}
\KMset{lem:unitid}{B.3}
\KMset{lem:translations}{A.8}
\KMset{prop:localid}{4.6}
\KMset{def:fd}{4.1}
\KMset{def:equality}{4.4}
\KMset{def:ident}{4.5}
\KMset{cor:localid}{4.9}
\KMset{conj:equality}{1}
\KMset{sec:reduce}{3}
\KMset{sec:count}{4.1}
\KMset{sec:setup}{2}
\KMset{eq:dotF}{16}
\KMset{eq:twotok}{17}

\title{\Large\bfseries\boldmath Skip Connections Unlock Exactly $\dm^2$ Parameters per Analytic Layer}

\newif\ifanon \anonfalse
\ifanon
\author{Anonymous author}
\else
\author{Marko Karbevski\\
Independent researcher\\
{\small\texttt{markokarbevski@insimplicity.tech}}}
\fi
\date{}

\begin{document}
\maketitle
\vspace{-1.5em}

\begin{abstract}
It was recently shown that skip connections unlock at least $\dm^2$ functional dimensions per layer (resp. $\approx \dm^2/2$ under pre-RMSNorm), with equality when a network without skips has no redundancy beyond its symmetries: one $\GL(\dm)$ (resp. $\Ort(\dm)$) change of basis per hidden state, and the rescalings inside its layers. We prove that this is the case at any depth, for $\tanh$, SwiGLU and GeGLU layers with or without pre-RMSNorm, and for GELU layers under pre-RMSNorm. In these networks the skips therefore unlock exactly $\dm^2$ (resp. $\dm(\dm-1)/2$) parameters per layer, and the networks without skips are locally identifiable up to their symmetries. This proves a conjecture of \citet{karbevski2026atleast} in all its cases but one. In the proof, the input of a layer is followed to infinity along a ray, where the hidden units separate by their rates, and the layers are removed one at a time.
\end{abstract}

\section{Introduction}\label{sec:intro}

Skip connections are a building block of almost every deep network trained today \citep{he2016resnet}. \citet{mijoski2026worth} show that they change which functions a network can represent, and \citet{karbevski2026atleast}, cited as [KM] below, measure the change by the \emph{functional dimension}, the number of independent directions in which the parameters move the function \citep{grigsby2022functional}. At generic parameters, that is, for almost all weights, the skips raise it by at least $\dm^2$ per block, where $\dm$ is the model dimension, and by about half as much under pre-RMSNorm. The count is exact when the network without skips has no redundancy beyond its symmetries. [KM] prove this for networks of one block without normalization, and state it as a conjecture at any depth, for $\tanh$, GELU, SwiGLU and GeGLU blocks. These four kinds of MLP block are real-analytic, and they are the analytic layers of our title. Here we prove the conjecture in all its cases but one.

In a network without skips, the function does not change when the output weights of a block are multiplied by an invertible matrix and the input weights of the next block by its inverse. Under pre-RMSNorm the matrix must be orthogonal. A SwiGLU or GeGLU block has, in addition, one rescaling per hidden unit. So every hidden state of a skipless network carries $\dm^2$ redundant parameters, or $\dm(\dm-1)/2$ under pre-RMSNorm, and a residual network keeps one change of basis for all its hidden states. A skipless network has no redundancy beyond these symmetries when every direction of its parameters that is not a direction of the symmetries changes the function. Equivalently, the network is \emph{locally identifiable} up to its symmetries: nearby parameters with the same function differ only by such changes of basis and rescalings. Whether the function of a network determines its weights is a classical question, which \citet{sussmann1992} and \citet{fefferman1994} study for $\tanh$ networks. We ask it locally, for the blocks of current networks and with normalization between them: are skipless networks locally identifiable up to their symmetries?

\subsection{Results}\label{sec:results-intro}

\para{Main result.} Consider a skipless network of $L$ MLP blocks of one kind, each with at least $\dm$ hidden units, without biases and with $P$ parameters (scope in \cref{sec:setup}). Let $\G$ be the group of the changes of basis of one hidden state, $\GL(\dm)$ without normalization and $\Ort(\dm)$ under pre-RMSNorm, and let $s$ be the number of the rescalings inside the blocks: zero for $\tanh$ and GELU blocks, and the number of hidden units of the network for SwiGLU and GeGLU blocks. We prove that at generic parameters the network has the functional dimension
\[
P - (L+1)\dim\G - s ,
\]
which is the upper bound that the symmetries give, in the following cases (\cref{thm:skiplessbound}): for $\tanh$ blocks, without normalization and under pre-RMSNorm; for GELU blocks, under pre-RMSNorm; for SwiGLU and GeGLU blocks, without normalization and under pre-RMSNorm. For $\tanh$ and GELU blocks this includes RMSNorm with the constant $\varepsilon = 0$, where one scaling per hidden state joins the symmetries. For GELU blocks without normalization we prove the same when every block has at least $2\dm$ hidden units, by another argument (\cref{prop:widegelu}).

\para{What follows.} A residual network of the same blocks has the functional dimension $P - \dim\G - s$ \KM{Theorem}{thm:resdepth}. So in these cases the skips raise the functional dimension by exactly $L\dim\G$, that is, by $\dm^2$ per block without normalization and by $\dm(\dm-1)/2$ under pre-RMSNorm, for every $L \ge 1$ (\cref{cor:exactgain}). Without normalization and for $\varepsilon > 0$, these skipless networks are locally identifiable up to their symmetries (\cref{cor:skiplesslocalid}). This proves Conjecture~\KMn{conj:equality} of [KM] in all its cases but one, GELU blocks with fewer than $2\dm$ hidden units without normalization. For $\varepsilon = 0$ we also prove the count for residual networks, attention blocks included (\cref{thm:epszero}). \Cref{tab:cases} lists the cases.

\begin{table}[t]
\centering\small
\renewcommand{\arraystretch}{1.15}\setlength{\tabcolsep}{5pt}
\newcommand{\tref}[3]{\hyperref[#2:#3]{#1~\ref*{#2}(#3)}}
\begin{tabular}{@{}llll@{}}
\toprule
\textbf{Blocks} & \textbf{No normalization} & \textbf{Pre-RMSNorm, $\varepsilon > 0$} & \textbf{Pre-RMSNorm, $\varepsilon = 0$} \\
\midrule
$\tanh$ & \tref{Thm.}{thm:skiplessbound}{i} & \tref{Thm.}{thm:skiplessbound}{i} & \tref{Thm.}{thm:skiplessbound}{i} \\
GELU & \hyperref[prop:widegelu]{Prop.~\ref*{prop:widegelu}}, if $\dhid \ge 2\dm$ & \tref{Thm.}{thm:skiplessbound}{ii} & \tref{Thm.}{thm:skiplessbound}{ii} \\
SwiGLU, GeGLU & \tref{Thm.}{thm:skiplessbound}{iii} & \tref{Thm.}{thm:skiplessbound}{iii} & open \\
\bottomrule
\end{tabular}
\caption{Skipless networks of $L \ge 1$ blocks of one kind, each with $\dhid \ge \dm$ hidden units. Each cell names the result which proves that the functional dimension at generic parameters is $P - (L+1)\dim\G - s$, so that the skips add exactly $L\dim\G$.}
\label{tab:cases}
\end{table}

\subsection{Idea of the proof}\label{sec:idea}

\para{The obstacle.} For residual networks, [KM] replace the output weights $\Wwrite^{\ell}$ of every block $\ell$ by $\tau\,\Wwrite^{\ell}$, with a small number $\tau$. To first order in $\tau$ every block then reads the embedded input, and the network is one wide block, which reduces the count to the case of one block \KM{Lemma}{lem:smallwrites}. Without skips this fails: a block with small output weights passes almost nothing on, and the network does not become one block. The blocks have to be told apart inside their composition.

\para{Three steps.} We remove the blocks one at a time, from the first to the last.
\begin{enumerate}[leftmargin=1.8em]
\item A direction of the parameters along which the function does not change gives an identity in the input of a block. Without normalization, and under pre-RMSNorm with $\varepsilon > 0$, these inputs fill an open set, so by analyticity the identity holds at every row.
\item We follow this input to infinity along a ray, as \citet{sussmann1992} does for $\tanh$ networks with one hidden layer. The hidden units of a $\tanh$ block saturate at different exponential rates. The hidden units of a GELU block and the gates of a SwiGLU or GeGLU block differ from linear or vanishing ones by small terms of different rates. A sum of terms with different rates vanishes only if every term does.
\item Term by term, the direction acts on the weights of the block as a symmetry does. What is left is the same statement for the network after this block.
\end{enumerate}
A GELU block is linear along the ray, up to its small terms. Without normalization the small terms of all blocks are then of the same order, and the block cannot be removed on its own, which is why \cref{thm:skiplessbound} leaves this case out (\cref{sec:raysgelu}). The wide blocks of \cref{prop:widegelu} come from the construction of \citet{mijoski2026worth}, which absorbs a skip into $\dm$ more hidden units. For $\varepsilon = 0$ the inputs of a block lie on a sphere, and the ray becomes a great circle at complex angles (\cref{sec:circle}).

\subsection{Related work}\label{sec:related}

\para{Identifiability.} For $\tanh$ networks, \citet{sussmann1992} and \citet{fefferman1994} prove that, under conditions on the weights, the function determines the weights up to permutations and sign flips of the hidden units. \Cref{rem:fefferman} derives one case of \cref{thm:skiplessbound} from the theorem of Fefferman. \citet{vlacic2020affine} define identifiability up to an equivalence relation on networks, and for ReLU networks \citet{stock2021embedding} and \citet{bonapellissier2022local} study identifiability, and local identifiability, up to permutations and rescalings of the hidden units. For skipless networks of attention blocks with one head and without softmax, \citet{henry2025lightning} prove that, at generic parameters, the function determines the weights up to scalings of the blocks and changes of basis inside and between them, and \citet{henry2026generic} determines the parameters with the same function for one attention block with several heads. \citet{dahlgren2026polynomial} prove that residual networks with polynomial activations, whose blocks have no output weights, are generically identifiable up to finitely many sign flips. Our statement is local and holds at generic parameters. It concerns skipless networks whose blocks meet in a hidden state of dimension $\dm$, with or without normalization between them, for $\tanh$, GELU, SwiGLU and GeGLU blocks.

\para{Dimension counts.} \citet{grigsby2022functional} bound the functional dimension of ReLU networks by their rescaling symmetries, and ReLU networks can lose more dimensions than their symmetries account for, on sets of parameters of positive measure \citep{grigsby2023hidden}. For polynomial activations, \citet{kileel2019expressive} compute the dimension of the functional variety, and \citet[Theorem~4.1]{shahverdi2026razor} prove that an MLP whose activation is a generic polynomial of large enough degree has as many dimensions of functions as weights. Such an MLP is a skipless network with $\dhid = \dm$ in which each product $\Wwrite^{\ell}\Wread^{\ell+1}$ is one weight matrix, so this is the analogue of \cref{thm:skiplessbound} without normalization. For attention, \citet{henry2026generic} gives the functional dimension of one block and conjectures it for attention-only transformers with skip connections, which [KM] prove. We give an exact count without skips, for MLP blocks in current use and with normalization.

\section{Setup}\label{sec:setup}
This section fixes the networks and recalls what we use of [KM], the paper of \citet{karbevski2026atleast}. We cite its statements with their numbers there: for example, \KM{Theorem}{thm:resdepth} is Theorem~\KMn{thm:resdepth} of that paper. The notation is that of [KM], with three differences that shorten the proofs: the map of block $\ell$ is $g_\ell$, where [KM] write $B_\ell$; the number of its hidden units is $q_\ell$, where [KM] write $\dhid$; and $N$ stands for the map that a block applies to a hidden state before it reads it, with or without normalization. The scope is MLP blocks of one kind with at least $\dm$ hidden units, no biases, and continuous inputs.

\para{Blocks.} A hidden state is a row $h \in \R^{1\times\dm}$: an MLP block acts on each token of a sequence separately, so we take inputs with one token. Block $\ell$ has two weight matrices: its input weights, the \emph{read} $\Wread^{\ell}$, and its output weights, the \emph{write} $\Wwrite^{\ell}$. For two rows of the same length, $\odot$ is their entrywise product. Let block $\ell$ have $q_\ell$ hidden units, and let $a_1, \dots, a_{q_\ell}$ be the rows of its write $\Wwrite^{\ell} \in \R^{q_\ell\times\dm}$. A $\tanh$ or GELU block has the activation $\sigma$, and its read $\Wread^{\ell} \in \R^{\dm\times q_\ell}$ has the columns $c_1, \dots, c_{q_\ell}$. A SwiGLU or GeGLU block, a \emph{gated} block, has the read $\Wread^{\ell} = [\,W_g^{\ell} \mid W_v^{\ell}\,]$, where $W_g^{\ell} \in \R^{\dm\times q_\ell}$ holds the gate columns $c_1, \dots, c_{q_\ell}$ and $W_v^{\ell} \in \R^{\dm\times q_\ell}$ the value columns $e_1, \dots, e_{q_\ell}$. Its gate $\sigma$ is GELU in a GeGLU block and SiLU, $\sigma(z) = z/(1 + e^{-z})$, in a SwiGLU block. Block $\ell$ maps its input $n$ to
\begin{align*}
g_\ell(n) \;&=\; \sigma(n\,\Wread^{\ell})\,\Wwrite^{\ell} \;=\; \sum_{j=1}^{q_\ell}\sigma(nc_j)\,a_j && \text{for a $\tanh$ or GELU block,}\\
g_\ell(n) \;&=\; \big(\sigma(n\,W_g^{\ell})\odot n\,W_v^{\ell}\big)\,\Wwrite^{\ell} \;=\; \sum_{j=1}^{q_\ell}\sigma(nc_j)\,(ne_j)\,a_j && \text{for a gated block.}
\end{align*}

\para{Networks.} A network has a linear embedding $E \in \R^{d_{\mathrm{in}}\times\dm}$, then $L$ blocks, then a linear readout $D \in \R^{\dm\times d_{\mathrm{out}}}$. Its parameters are
\[
\theta \;=\; \Big(E,\ \big\{\Wread^{\ell},\ \Wwrite^{\ell}\big\}_{\ell=1}^{L},\ D\Big) \in \R^P .
\]
Before a block or the readout reads a hidden state $h$, it applies a map $N$ to it, which the next paragraph specifies. The \emph{skipless} network maps an input $x \in \R^{1\times d_{\mathrm{in}}}$ to $F_\theta(x)$ through the hidden states
\[
h_0 = xE, \qquad h_\ell = g_\ell\big(N(h_{\ell-1})\big) \quad (1 \le \ell \le L), \qquad F_\theta(x) = N(h_L)\,D ,
\]
and the \emph{residual} network has $h_\ell = h_{\ell-1} + g_\ell\big(N(h_{\ell-1})\big)$ in the place of the second formula. Both have the same $P$ parameters. A statement holds at \emph{generic} parameters if it holds outside a set of Lebesgue measure zero.

\para{The three cases.} We write (n) for the case without normalization, (+) for pre-RMSNorm with $\varepsilon > 0$, and (0) for pre-RMSNorm with $\varepsilon = 0$. RMSNorm is the map $N_\varepsilon$ of the table, with a constant $\varepsilon \ge 0$ that is fixed and is not a parameter. The map $N$ is the identity in the case (n), and $N_\varepsilon$ in the cases (+) and (0). The table gives $N$, the set of its values, which are the inputs of the blocks, the group $\G$ of the changes of basis, and its Lie algebra $\mathfrak g$: the set of the matrices $Z$ with $\exp(\tau Z) \in \G$ for all numbers $\tau$.
\begin{center}\small
\begin{tabular}{@{}lllll@{}}
\toprule
Case & $N(h)$ & Values of $N$ & $\G$ & $\mathfrak g$ \\
\midrule
(n) & $h$ & all rows $n \in \R^{1\times\dm}$ & $\GL(\dm)$ & all matrices $Z$ \\[2pt]
(+) & $h\,\big(\norm{h}^2/\dm + \varepsilon\big)^{-1/2}$ & the ball $\norm{n}^2 < \dm$ & $\Ort(\dm)$ & antisymmetric $Z$ \\[2pt]
(0) & $\sqrt{\dm}\;h/\norm{h}$, for $h \neq 0$ & the sphere $\mathbb S$: $\norm{n}^2 = \dm$ & $\CO(\dm)$ & $Z = Y + \beta I$ \\
\bottomrule
\end{tabular}
\end{center}
In the last row $Y$ is antisymmetric and $\beta$ is a number, and $\mathbb S$ is the set of the rows $n \in \R^{1\times\dm}$ with $\norm{n}^2 = \dm$. In the case (0) a network is defined at the inputs $x$ for which no hidden state vanishes.

\para{Functional dimension.} For a finite batch $X = (x_1, \dots, x_n)$ of inputs, let $F_X(\theta)$ hold the outputs $F_\theta(x_1), \dots, F_\theta(x_n)$, and let $\Der F_X(\theta)$ be the Jacobian of $\theta \mapsto F_X(\theta)$. The \emph{functional dimension} at $\theta$ is
\[
\operatorname{fd}(\theta) \;=\; \max_X\,\operatorname{rank}\Der F_X(\theta),
\]
where the maximum is over all finite batches $X$ (\citealp{grigsby2022functional}; \KM{Definition}{def:fd}). We write $\Der_\theta F(\theta, x)\,\eta$ for the derivative of $F_\theta(x)$ along a direction $\eta \in \R^P$ of the parameters.

\para{Symmetries.} The symmetries of a skipless network are changes of basis at its hidden states (\citealp{godfrey2022symmetries}; \citealp{zhao2023symmetries}; [KM, Section~\KMn{sec:reduce}]). For $G \in \G$ let $\widehat G = G$ in the cases (n) and (+), and $\widehat G = Q$ for $G = \lambda Q$ in the case (0), with a number $\lambda > 0$ and an orthogonal $Q$. Then $N(hG) = N(h)\,\widehat G$ in all three cases. So for $G_0, \dots, G_L \in \G$ the parameters
\[
\Big(EG_0,\ \ \big\{\widehat G_{\ell-1}^{-1}\Wread^{\ell},\ \ \Wwrite^{\ell}G_\ell\big\}_{\ell=1}^{L},\ \ \widehat G_L^{-1}D\Big)
\]
give the same function, with the hidden states $h_\ell G_\ell$: one matrix of $\G$ per hidden state. For $Z \in \mathfrak g$ let $\widehat Z = Z$ in the cases (n) and (+), and $\widehat Z = Y$ for $Z = Y + \beta I$ in the case (0). For $Z_0, \dots, Z_L \in \mathfrak g$ let
\begin{equation}\label{eq:e0eta}
\eta^{Z} \;=\; \Big(EZ_0,\ \ \big\{-\widehat Z_{\ell-1}\Wread^{\ell},\ \ \Wwrite^{\ell}Z_\ell\big\}_{\ell=1}^{L},\ \ -\widehat Z_LD\Big) \in \R^P ,
\end{equation}
the derivative of these parameters at $\tau = 0$ for $G_\ell = \exp(\tau Z_\ell)$. So $\Der_\theta F(\theta, x)\,\eta^{Z} = 0$ at every input at which $F_\theta$ is defined. When $E$ and every $\Wwrite^{\ell}$ have rank $\dm$, the map $(Z_0, \dots, Z_L) \mapsto \eta^{Z}$ is injective.

A gated block has, in addition, one rescaling per hidden unit. For diagonal matrices $\Delta_1, \dots, \Delta_L$, with $\Delta_\ell$ of size $q_\ell$, the parameters
\[
\Big(EG_0,\ \ \big\{G_{\ell-1}^{-1}W_g^{\ell},\ \ G_{\ell-1}^{-1}W_v^{\ell}\,e^{-\Delta_\ell},\ \ e^{\Delta_\ell}\,\Wwrite^{\ell}G_\ell\big\}_{\ell=1}^{L},\ \ G_L^{-1}D\Big)
\]
give the same function as well: the diagonal matrix $e^{\Delta_\ell}$ multiplies the write row of every hidden unit of block $\ell$ by a positive number and divides its value column by the same number. For $Z_0, \dots, Z_L \in \mathfrak g$ let
\begin{equation}\label{eq:gteta}
\eta^{Z,\Delta} \;=\; \Big(EZ_0,\ \big\{-Z_{\ell-1}W_g^{\ell},\ -Z_{\ell-1}W_v^{\ell} - W_v^{\ell}\Delta_\ell,\ \Delta_\ell\Wwrite^{\ell} + \Wwrite^{\ell}Z_\ell\big\}_{\ell=1}^{L},\ -Z_LD\Big) \in \R^P ,
\end{equation}
the derivative of these parameters at $\tau = 0$ for $G_\ell = \exp(\tau Z_\ell)$ and with $\tau\Delta_\ell$ in the place of $\Delta_\ell$. So $\Der_\theta F(\theta, x)\,\eta^{Z,\Delta} = 0$ at every input $x$. When $E$ and every $\Wwrite^{\ell}$ have rank $\dm$ and no value column vanishes, $\eta^{Z,\Delta} = 0$ only if every $Z_\ell$ and every $\Delta_\ell$ vanishes: $EZ_0 = 0$ gives $Z_0 = 0$, then $W_v^{1}\Delta_1 = 0$ gives $\Delta_1 = 0$, then $\Wwrite^{1}Z_1 = 0$ gives $Z_1 = 0$, and so on.

Let $s = 0$ for $\tanh$ and GELU blocks and $s = q_1 + \dots + q_L$ for gated blocks. Under these conditions the directions of the symmetries form a space of dimension $(L+1)\dim\G + s$, so on every batch the Jacobian has rank at most
\begin{equation}\label{eq:e0upper}
P - (L+1)\dim\G - s .
\end{equation}
This is the upper bound of \KM{Theorem}{thm:unlock} in the cases (n) and (+).

\para{Equality condition and local identifiability.} A network satisfies the \emph{equality condition} if its functional dimension at generic parameters is the bound that its symmetries give \KM{Definition}{def:equality}: the bound \eqref{eq:e0upper} for a skipless network, and $P - \dim\G - s$ for a residual one, which keeps one change of basis for all its hidden states. A network is \emph{locally identifiable} up to its symmetries at $\theta$ if every $\theta'$ near $\theta$ with $F_{\theta'} = F_\theta$ is the image of $\theta$ under a symmetry \KM{Definition}{def:ident}. In the cases (n) and (+), a skipless or a residual network of the blocks above satisfies the equality condition if and only if it is locally identifiable up to its symmetries at generic parameters, and every symmetry that is needed is then close to the identity \KM{Proposition}{prop:localid}.

\para{The residual side.} In the cases (n) and (+), the residual network of the blocks above satisfies the equality condition: at generic parameters its functional dimension is $P - \dim\G - s$, with the same $s$ \KM{Theorem}{thm:resdepth}. That theorem also covers attention blocks, where $s$ counts the changes of basis inside the heads as well. We call these and the rescalings of gated blocks the \emph{symmetries inside the blocks}. For the case (0) we prove the same in \cref{thm:epszero}.

\section{Results}\label{sec:results}

\Cref{thm:skiplessbound} is the main result. \Cref{thm:epszero} gives the side with skips for $\varepsilon = 0$, \cref{prop:widegelu} adds one case by another argument, and the two corollaries combine them with [KM].

\begin{restatable}[Skipless networks attain their bound]{theorem}{thmSkipless}\label{thm:skiplessbound}\label{thm:epszeroskipless}
Let $d_{\mathrm{in}}, d_{\mathrm{out}} \ge \dm \ge 2$, and consider skipless networks without biases whose blocks are all of one kind, $\tanh$, GELU, SwiGLU or GeGLU, each with $q_\ell \ge \dm$ hidden units. In the following cases the functional dimension at generic parameters is the bound \eqref{eq:e0upper},
\[
P - (L+1)\dim\G - s ,
\]
so the network satisfies the equality condition.\\
\stpart{thm:skiplessbound}{i}{} $\tanh$ blocks: without normalization, and under pre-RMSNorm with $\varepsilon > 0$ and with $\varepsilon = 0$.\\
\stpart{thm:skiplessbound}{ii}{} GELU blocks: under pre-RMSNorm with $\varepsilon > 0$ and with $\varepsilon = 0$.\\
\stpart{thm:skiplessbound}{iii}{} SwiGLU blocks, and GeGLU blocks: without normalization, and under pre-RMSNorm with $\varepsilon > 0$.
\end{restatable}
\proofin{app:skiplessbound}

For GELU blocks the theorem needs no more than $q_\ell \ge \dm$ hidden units per block. \Cref{sec:proof} presents the proof, and says why part (ii) has no case without normalization.

\begin{restatable}[Residual networks under RMSNorm with $\varepsilon = 0$]{theorem}{thmEpsZero}\label{thm:epszero}
Let $d_{\mathrm{in}}, d_{\mathrm{out}} \ge \dm \ge 2$, and consider residual networks under pre-RMSNorm with $\varepsilon = 0$, without biases, in which each block, in any order and of any width, is an attention block with heads of width at most $\dm$, causal or not, on sequences of at least two tokens, or an MLP block, with all MLP blocks $\tanh$, all GELU, all SwiGLU or all GeGLU blocks. At generic parameters the functional dimension is
\[
P - \dim\CO(\dm) - s \;=\; P - \dm(\dm-1)/2 - 1 - s ,
\]
where $s$ is the dimension of the symmetries inside the blocks. So the residual network satisfies the equality condition.
\end{restatable}
\proofin{app:epszero}

The proof follows that of \KM{Theorem}{thm:resdepth}: with small writes the network is one wide block to first order, and what is new for $\varepsilon = 0$ is that this block reads only rows of a sphere (\cref{lem:sphere}).

\begin{restatable}[Wide GELU blocks without normalization]{proposition}{propWideGelu}\label{prop:widegelu}
Without normalization, let $d_{\mathrm{in}}, d_{\mathrm{out}} \ge \dm \ge 2$, and consider skipless networks of GELU blocks without biases, each with $q_\ell \ge 2\dm$ hidden units. At generic parameters the functional dimension is $P - (L+1)\dm^2$, so the network satisfies the equality condition.
\end{restatable}
\proofin{app:widegelu}

The proof starts from a residual network. Since GELU satisfies $\sigma(z) - \sigma(-z) = z$, a residual block equals a skipless block with $\dm$ more hidden units \citep{mijoski2026worth}. At small writes, the columns of the Jacobian of this skipless network span those of the Jacobian of the residual network, whose rank is known \KM{Lemma}{lem:smallwrites}, and the reads of the added hidden units give $\dm^2$ more independent columns per block. For skipless $\tanh$ networks without normalization, the equality condition also follows from the theorem of \citet{fefferman1994}, with or without biases (\cref{rem:fefferman}).

\begin{corollary}[The exact gain]\label{cor:exactgain}
In the cases of \cref{thm:skiplessbound} and of \cref{prop:widegelu}, at generic parameters and for every $L \ge 1$, the skips raise the functional dimension by exactly $L\dim\G$: by $L\,\dm^2$ without normalization, by $L\,\dm(\dm-1)/2$ under pre-RMSNorm with $\varepsilon > 0$, and by $L\,\big(\dm(\dm-1)/2 + 1\big)$ for $\varepsilon = 0$.
\end{corollary}

\begin{proof}
The skipless network has the functional dimension $P - (L+1)\dim\G - s$, by \cref{thm:skiplessbound} or \cref{prop:widegelu}. The residual network has the same parameters and the functional dimension $P - \dim\G - s$, by \KM{Theorem}{thm:resdepth} in the cases (n) and (+), and by \cref{thm:epszero} in the case (0). The difference is $L\dim\G$.
\end{proof}

For one block under pre-RMSNorm, \KMp{Theorem}{thm:oneblock}{ii} gives the gain as a lower bound, and the corollary makes it exact for the four kinds of MLP block.

\begin{corollary}[Skipless networks are locally identifiable]\label{cor:skiplesslocalid}
Without normalization or under pre-RMSNorm with $\varepsilon > 0$, in the cases of \cref{thm:skiplessbound} and of \cref{prop:widegelu}, let $\theta$ be generic parameters of the skipless network. Every $\theta'$ near $\theta$ with the same function is obtained from $\theta$ by matrices $G_0, \dots, G_L \in \G$ close to the identity, one per hidden state, as in the displays of \cref{sec:setup}, and for SwiGLU and GeGLU blocks by diagonal matrices $\Delta_1, \dots, \Delta_L$ close to zero.
\end{corollary}

This is the equality condition together with \KM{Proposition}{prop:localid}. In the residual network the matrices $G_0, \dots, G_L$ are all equal \KM{Corollary}{cor:localid}. For $\varepsilon = 0$ we state no such corollary, since \KM{Proposition}{prop:localid} is stated for $\varepsilon > 0$.

\para{Open cases.} Among the networks of \cref{tab:cases}, two cases are open: GELU blocks with $\dm \le q_\ell < 2\dm$ without normalization, which is what remains of Conjecture~\KMn{conj:equality} of [KM], and SwiGLU and GeGLU blocks with $\varepsilon = 0$. For attention blocks, the equality condition for skipless networks is proved for one block \KM{Theorem}{thm:oneblock} and is open at depth.

\section{Proof of Theorem~\ref*{thm:skiplessbound}}\label{sec:proof}

The proof removes one block at a time. \Cref{sec:claim} reduces the theorem to one step, from the network after block $k+1$ to the network after block $k$. \Cref{sec:raystanh,sec:raysgelu,sec:circle} make this step for $\tanh$ blocks, for GELU and gated blocks, and for $\varepsilon = 0$. \Cref{app:skiplessbound} has the proofs.

\subsection{One claim for every hidden state}\label{sec:claim}

For $k = 0, \dots, L$ and a row $h$, let $\mathcal B_{>k}(h)$ be the output of the network when its hidden state $h_k$ is set to $h$:
\[
\mathcal B_{>L}(h) = N(h)\,D, \qquad \mathcal B_{>k-1}(h) = \mathcal B_{>k}\big(g_k(N(h))\big), \qquad\text{so that}\quad F_\theta(x) = \mathcal B_{>0}(xE).
\]
Let $\Der\mathcal B_{>k}(h)\,v$ be its derivative at $h$ along a row $v$, and $\Der_\theta\mathcal B_{>k}(h)\,\eta$ its derivative along a direction $\eta$ of the weights on which it depends, those of the blocks $k+1, \dots, L$ and of $D$. For $Z_k, \dots, Z_L \in \mathfrak g$, the symmetries with $G_\ell = \exp(\tau Z_\ell)$ change the hidden state $h_k$ into $h_k\exp(\tau Z_k)$ and move the later weights along a direction $\eta$, and they do not change the output. To first order in $\tau$, with $Z = Z_k$,
\begin{equation}\label{eq:mainD}
\Der\mathcal B_{>k}(h)\,(hZ) + \Der_\theta\mathcal B_{>k}(h)\,\eta \;=\; 0 \qquad\text{for all rows } h .
\end{equation}
The claim is the converse.

\smallskip
\noindent\emph{Claim $(D_k)$}, for $\tanh$ and GELU blocks in the cases (n) and (+). If a $\dm\times\dm$ matrix $Z$ and a direction $\eta$ of the weights of the blocks $k+1, \dots, L$ and of $D$ satisfy \eqref{eq:mainD}, then $Z \in \mathfrak g$, and with $Z_k = Z$ there are $Z_{k+1}, \dots, Z_L \in \mathfrak g$ such that
\[
\eta \;=\; \Big(\big\{-Z_{\ell-1}\Wread^{\ell},\ \ \Wwrite^{\ell}Z_\ell\big\}_{\ell=k+1}^{L},\ \ -Z_LD\Big).
\]
\smallskip

\Cref{app:claim} states the claim in all cases: for gated blocks $\eta$ also holds the directions of the rescalings, as in \eqref{eq:gteta}, and in the case (0) only the antisymmetric part of $Z_{\ell-1}$ acts on $\Wread^{\ell}$, as in \eqref{eq:e0eta}.

The claim $(D_L)$ concerns the readout alone and is linear algebra (\cref{lem:e0readout}). The claim $(D_0)$ gives the theorem. If a direction $\eta \in \R^P$ has $\Der_\theta F(\theta, x)\,\eta = 0$ at every input $x$, then its entry at $E$ is $EZ$ for a matrix $Z$, which satisfies \eqref{eq:mainD} for $k = 0$ with the rest of $\eta$ (\cref{lem:e0atwrite}). By $(D_0)$, $\eta$ is then one of the directions $\eta^Z$ of \eqref{eq:e0eta}, and the Jacobian has the rank \eqref{eq:e0upper} on some finite batch (\cref{prop:e0givesrank}). So the proof is the step from $(D_{k+1})$ to $(D_k)$.

\para{The identity for one block.} Let block $k+1$ be a $\tanh$ or GELU block with the read $W$ and the write $V$, in the case (n), and write $\mathcal B = \mathcal B_{>k+1}$ and $\eta = (\eta_W, \eta_V, \eta')$, where $\eta_W$ and $\eta_V$ are the entries of $\eta$ at $W$ and $V$. By the chain rule, \eqref{eq:mainD} at the row $n$ reads
\begin{equation}\label{eq:mainL}
\Der\mathcal B\big(\tilde wV\big)\,\Big(\big(\tilde w'\odot n\,(ZW + \eta_W)\big)\,V + \tilde w\,\eta_V\Big) + \Der_\theta\mathcal B\big(\tilde wV\big)\,\eta' \;=\; 0 ,
\end{equation}
where $\tilde w = \sigma(nW)$ is the row of the outputs of the hidden units of the block, and $\tilde w' = \sigma'(nW)$ the row of their derivatives. In the case (+) the block reads $n = N_\varepsilon(h)$, the derivative of $N_\varepsilon$ adds one term to $n\,(ZW + \eta_W)$, and the identity holds for the rows $n$ of the open ball $\norm{n}^2 < \dm$. Its left side is real-analytic in $n$, so the identity holds at every row $n \in \R^{1\times\dm}$ (\cref{lem:sballrows}). This is what lets us follow $n$ to infinity.

\subsection{\texorpdfstring{$\tanh$}{tanh} blocks on a ray}\label{sec:raystanh}

Take $n = t\,\xi$ with a row $\xi$, and let $t \to \infty$. With $u_j = \xi c_j$, hidden unit $j$ outputs
\[
\tanh(t\,u_j) \;=\; \pm\,\frac{1 - \zeta_j(t)}{1 + \zeta_j(t)}, \qquad \zeta_j(t) = e^{-2\lvert u_j\rvert t},
\]
with the sign of $u_j$: it saturates at the exponential rate $2\lvert u_j\rvert$. When the numbers $u_j$ are linearly independent over $\mathbb Q$, the hidden units can be told apart.

\begin{restatable}[Separation on a ray]{lemma}{lemRaySep}\label{lem:raysep}
Let $u_1, \dots, u_q$ be real numbers that are linearly independent over $\mathbb Q$, and let $\mathcal T_0, \dots, \mathcal T_K$ be real-analytic maps from $\R^{1\times q}$ to $\R^{1\times d_{\mathrm{out}}}$. If
\[
\sum_{i=0}^{K} t^i\ \mathcal T_i\big(\tanh(tu_1), \dots, \tanh(tu_q)\big) \;=\; 0 \qquad\text{for all } t > 0,
\]
then $\mathcal T_i = 0$ on $\R^{1\times q}$ for every $i$.
\end{restatable}

In the proof, every $\mathcal T_i$ is expanded in powers of $\zeta_1, \dots, \zeta_q$. The monomials decay at pairwise different rates, because the $u_j$ are linearly independent over $\mathbb Q$, so the slowest term that is left must vanish. This is the argument by which \citet[Lemma~3.1]{sussmann1992} proves that $\tanh$ hidden units with different arguments are linearly independent, applied here to analytic maps of their outputs.

For $\sigma = \tanh$ we have $\sigma' = 1 - \sigma^2$, so the left side of \eqref{eq:mainL} at $n = t\,\xi$ is $t\,\mathcal T_1(\tilde w) + \mathcal T_0(\tilde w)$ at $\tilde w = \tanh(t\,\xi W)$, with
\[
\mathcal T_1(\tilde w) = \Der\mathcal B(\tilde wV)\,\Big(\big((1 - \tilde w^2)\odot\xi\,(ZW + \eta_W)\big)\,V\Big), \qquad \mathcal T_0(\tilde w) = \Der\mathcal B(\tilde wV)\,(\tilde w\,\eta_V) + \Der_\theta\mathcal B(\tilde wV)\,\eta' .
\]
By the lemma both maps vanish at every row $\tilde w \in \R^{1\times q}$: the outputs of the hidden units have become free variables. From $\mathcal T_1 = 0$ we get $\eta_W = -ZW$, since $\Der\mathcal B$ is injective at the rows of an open dense set (\cref{lem:e0tailrank}) and no row of $V$ vanishes. From $\mathcal T_0 = 0$ we get $\eta_V = VZ'$ for a matrix $Z'$ that satisfies \eqref{eq:mainD} for $k+1$ with $\eta'$ (\cref{lem:e0atwrite}), and $(D_{k+1})$ gives the rest. In the case (+) there is one more term, $t^3\,\mathcal T_3(\tilde w)$, and $\mathcal T_3 = 0$ says that $Z$ is antisymmetric (\cref{lem:sbsteptanh}). The step needs that no row of $V$ vanishes and that the columns of $W$ are linearly independent over $\mathbb Q$, which holds for almost all weights.

\subsection{GELU and gated blocks on a ray}\label{sec:raysgelu}

GELU and SiLU, the gate of a SwiGLU block, do not saturate. Both satisfy $\sigma(z) - \sigma(-z) = z$, so with $\delta(y) = y - \sigma(y)$ we have, for every $y > 0$,
\[
\sigma(y) = y - \delta(y), \qquad \sigma(-y) = -\delta(y),
\]
and $\delta$ is small: for $y \ge 2$,
\[
\tfrac16\,e^{-y^p/p} \;\le\; \delta(y) \;\le\; y\,e^{-y^p/p}, \qquad p = 2 \text{ for GELU}, \quad p = 1 \text{ for SiLU} .
\]
So on a ray a GELU hidden unit, and the gate of a hidden unit of a SwiGLU or GeGLU block, differ from a linear or a vanishing one by a small term, whose rate depends on the hidden unit through $\lvert\xi c_j\rvert$. A GELU block has the form of a GeGLU block in which the value of every hidden unit is the number $1$, and we treat the three kinds of block together (\cref{lem:gtstep}).

The identity for one block is expanded in the products of these small terms and of their derivatives, and the products are removed in the order of their rates. The term without a small factor is the identity for the block with its gates \emph{frozen}: every hidden unit that is active on the ray is linear in a GELU block and quadratic in a gated block, and the others vanish. A term with one small factor belongs to one hidden unit. For GELU and GeGLU its coefficient holds the Mills ratio
\[
\operatorname{mr}(z) \;=\; \sqrt{\frac{\pi}{2}}\;e^{z^2/2}\,\Big(1 - \operatorname{erf}\frac{z}{\sqrt2}\Big),
\]
through $\delta(y) = (2\pi)^{-1/2}\,e^{-y^2/2}\,y\,\operatorname{mr}(y)$, and the next lemma separates it from the rest of the network.

\begin{restatable}[The Mills ratio at infinity]{lemma}{lemMills}\label{lem:e0nomero}
Let $z(t) = \omega\,t + \omega'/t$ with $\omega, \omega' \in \mathbb C$ and $\operatorname{Re}\omega > 0$, and let $\Theta_1$ and $\Theta_2$ be meromorphic at infinity, with values in $\mathbb C^{1\times d_{\mathrm{out}}}$. If $\Theta_1(t)\,\operatorname{mr}(z(t)) + \Theta_2(t)$ tends to $0$ faster than every power of $1/t$ as $t \to +\infty$, then $\Theta_1 = 0$ and $\Theta_2 = 0$.
\end{restatable}

For SwiGLU a parity argument takes the place of the lemma. The terms with two or more small factors follow by induction. At the end the direction does not move what the gates read, and the identity holds with free variables in the place of the outputs of the hidden units, as for $\tanh$.

\para{The later blocks.} The coefficients of the expansion come from the network after the block. Under RMSNorm the later blocks read a normalized row, which stays bounded along the ray, and the coefficients are meromorphic at infinity (\cref{lem:sbcoeff}). Without normalization the inputs of the later blocks grow along the ray, and their hidden units add small terms of their own. A SwiGLU or GeGLU block is quadratic along the ray up to its small terms, so the inputs of the later blocks grow like the square of its own input. Their small terms are then of a smaller order, and the network after the block can be replaced by a polynomial map (\cref{lem:gtfar}). A GELU block is linear along the ray, the small terms of all blocks are of the same order, and the block cannot be removed on its own. This is why \cpref{thm:skiplessbound}{ii} has no case without normalization.

\para{The parameters.} For $\tanh$ blocks the step holds at almost every $\theta$. For GELU and gated blocks we prove the claim at one explicit choice of parameters $\theta^\dagger$: $E$ and $D$ are the identity padded with zeros, and in every block
\[
W_g^{\ell} = \big[\,I \;\big|\; {-\mu_1\mathbf 1} \;\big|\; \cdots \;\big|\; {-\mu_{q-\dm}\mathbf 1}\,\big], \qquad \Wwrite^{\ell} = \begin{bmatrix} I \\ \mathbf 1^\top \\[-3pt] \vdots \\ \mathbf 1^\top \end{bmatrix}, \qquad W_v^{\ell} = \big[\,\mathbf 1 \;\big|\; \cdots \;\big|\; \mathbf 1\,\big] ,
\]
where $\mathbf 1$ is the column of $\dm$ ones, $W_g^{\ell} = \Wread^{\ell}$ in a GELU block, and only a gated block has $W_v^{\ell}$. On rows with positive entries the first $\dm$ hidden units are active and the others are not, so the frozen block maps $n$ to $n$ in a GELU block and to $(n\mathbf 1)\,n$ in a gated block, and every hidden state stays in the positive orthant (\cref{lem:gtstar}). The numbers $\mu_j > 0$ are chosen so that, on a suitable ray, the rates of the hidden units are linearly independent over $\mathbb Q$. At $\theta^\dagger$ the Jacobian then has the rank \eqref{eq:e0upper} on some finite batch. The network is analytic in $\theta$, so the Jacobian has this rank on the same batch at generic parameters.

\subsection{The case \texorpdfstring{$\varepsilon = 0$}{of no regularization}: a great circle}\label{sec:circle}

For $\varepsilon = 0$ the inputs of a block lie on the sphere $\mathbb S$, and the identity for one block holds only there. A ray leaves the sphere, so we follow a great circle and continue it to complex angles. For orthonormal rows $\xi, \xi'$ let $\nu = \xi + \mathrm i\xi'$ and
\[
n(t) \;=\; \frac{\sqrt{\dm}}{2}\,\Big(t\,\nu + \frac{1}{t}\,\bar\nu\Big), \qquad t \in \mathbb C\setminus\{0\}.
\]
Then $n(t)\,n(t)^\top = \dm$ for all $t$, without complex conjugation. For $\lvert t\rvert = 1$ the row $n(t)$ is real and runs through the great circle of $\xi$ and $\xi'$, and for real $t \to \infty$ it is a complex row that grows like $t$. The network after the block extends to a holomorphic map on a complex cone around the real rows (\cref{lem:e0cone}), so the identity continues along the path (\cref{lem:e0cont}). There the hidden units of a $\tanh$ block separate as on a ray (\cref{lem:e0sep}), and those of a GELU block through the Mills ratio: on the path a hidden unit reads $z(t) = \omega\,t + \omega'/t$, the form of \cref{lem:e0nomero}. We prove the claim at parameters $\theta^*$ that are chosen block by block, from the last to the first. The theorem follows by analyticity on the set of the parameters at which no hidden state vanishes on a batch, which is connected (\cref{lem:e0domain}).

\section{Conclusion}\label{sec:conclusion}

In the cases of \cref{tab:cases}, a network without skip connections has no redundancy beyond one change of basis per hidden state and the rescalings inside its blocks. So the lower bound of [KM] on what the skips add is an equality there: $\dm^2$ functional dimensions per block without normalization, and $\dm(\dm-1)/2$ under pre-RMSNorm. Without normalization and for $\varepsilon > 0$, these skipless networks are locally identifiable up to their symmetries, as the residual networks are.

\para{Open questions.} The equality condition is open for skipless networks of GELU blocks with fewer than $2\dm$ hidden units without normalization, of SwiGLU and GeGLU blocks with $\varepsilon = 0$, and of attention blocks at depth. Whether the networks of this paper are identifiable, and not only locally, is open as well. Networks with biases are not treated here, apart from \cref{rem:fefferman}.

\section*{Use of AI}
\addcontentsline{toc}{section}{Use of AI}
The conception of the study is the author's own. The author obtained the results and proofs working with AI assistants, through continuous conversation; the assistants were used to find, verify and simplify proofs.

\bibliographystyle{plainnat}
\bibliography{refs}


\appendix
\section{Proof of Theorem~\ref*{thm:skiplessbound}}\label{app:skiplessbound}\label{app:rmsskipless}\label{app:epszeroskipless}

We recall the theorem.

\thmSkipless*

The three cases (n), (+) and (0), the blocks, the symmetries with their directions $\eta^{Z}$ and $\eta^{Z,\Delta}$, and the bound \eqref{eq:e0upper} are those of \cref{sec:setup}. \Cref{app:claim} states Claim~$(D_k)$ and reduces the theorem to the step from $(D_{k+1})$ to $(D_k)$. \Cref{app:raystanh,app:raysgelu} make this step on rays, in the cases (n) and (+), and \cref{app:circle} on a great circle, in the case (0). \Cref{app:proofmain} puts the parts together. The gate $\sigma$ of a GeGLU block is GELU, and that of a SwiGLU block is SiLU. Gated blocks are treated in the cases (n) and (+) only.

\subsection{The claim, and the reduction to one step}\label{app:claim}

\para{Notation.} For rows $u, v$, real or complex, $\langle u, v\rangle = uv^\top$, without complex conjugation. For $k = 0, \dots, L$ and a row $h$, with $h \neq 0$ in the case (0), let $\mathcal B_{>k}(h)$ be the output of the network when its hidden state $h_k$ is set to $h$:
\[
\mathcal B_{>L}(h) = N(h)\,D, \qquad \mathcal B_{>k-1}(h) = \mathcal B_{>k}\big(g_k(N(h))\big), \qquad\text{so that}\quad F_\theta(x) = \mathcal B_{>0}(xE).
\]
$\mathcal B_{>k}$ depends on the weights of the blocks $k+1, \dots, L$ and on $D$. We write $\Der\mathcal B_{>k}(h)\,v$ for its derivative at $h$ along a row $v$, and $\Der_\theta\mathcal B_{>k}(h)\,\eta$ for its derivative at fixed $h$ along a direction $\eta$ of these weights. In the cases (+) and (0), for a $\dm\times\dm$ matrix $Z$ and $n = N(h)$,
\begin{equation}\label{eq:sbDN}
\Der N(h)\,(hZ) \;=\; nZ - \frac{\langle nZ, n\rangle}{\dm}\,n .
\end{equation}
In the case (0), $\mathcal B_{>k}(\lambda h) = \mathcal B_{>k}(h)$ for $\lambda > 0$, so
\begin{equation}\label{eq:e0radial}
\Der\mathcal B_{>k}(h)\,h = 0 .
\end{equation}
We call $\theta$ \emph{regular} if
\begin{itemize}[leftmargin=2.9em]
\item[(R1)] $E$, $D$ and every $\Wwrite^{\ell}$ have rank $\dm$;
\item[(R2)] in the case (0), for every $\ell$ and every $n \in \mathbb S$: $g_\ell(n) \neq 0$;
\item[(R3)] for every $\ell$: in the cases (n) and (+), $\Der g_\ell(n)$ is invertible at some row $n$; in the case (0), there is an $n \in \mathbb S$ such that no row $v \neq 0$ with $\langle v, n\rangle = 0$ has $\Der g_\ell(n)\,v \in \R\,g_\ell(n)$.
\end{itemize}
In the cases (n) and (+), every $\mathcal B_{>k}$ is real-analytic on all rows $h$. In the case (0), by (R2), it is defined and real-analytic on all rows $h \neq 0$.

\begin{lemma}[The derivative of $\mathcal B_{>k}$]\label{lem:e0tailrank}
Let $\theta$ be regular. For every $k$ there is an open dense set $\mathcal U_k$ of rows such that, for $h \in \mathcal U_k$, $\Der\mathcal B_{>k}(h)\,v = 0$ only if $v = 0$, in the cases (n) and (+), and only if $v \in \R\,h$, in the case (0). In the case (0), $\mathcal U_k$ consists of rows $h \neq 0$ and is closed under multiplication by positive numbers.
\end{lemma}

\begin{proof}
\emph{The cases (n) and (+).} $N$ is the identity, or a diffeomorphism of $\R^{1\times\dm}$ onto the ball of the rows $n$ with $\norm{n}^2 < \dm$, with the inverse $n \mapsto \sqrt{\varepsilon}\;n\,\big(1 - \norm{n}^2/\dm\big)^{-1/2}$. So $\Der N(h)$ is invertible at every $h$. As $D$ is injective on rows, $\Der\mathcal B_{>L}(h)\,v = \big(\Der N(h)\,v\big)D$ vanishes only for $v = 0$: $\mathcal U_L$ is the set of all rows. Let $k < L$. The determinant of $\Der g_{k+1}(n)$ is real-analytic in $n$ and does not vanish identically, by (R3). So $g_{k+1}$ is a local diffeomorphism on an open dense set of rows $n$, and $g_{k+1}\circ N$ on an open dense set of rows $h$. Let $\mathcal U_k$ be the set of the $h$ of this set with $g_{k+1}(N(h)) \in \mathcal U_{k+1}$. It is open and dense, and for $h \in \mathcal U_k$,
\[
\Der\mathcal B_{>k}(h)\,v \;=\; \Der\mathcal B_{>k+1}\big(g_{k+1}(N(h))\big)\,\Big(\Der\big(g_{k+1}\circ N\big)(h)\,v\Big)
\]
vanishes only for $v = 0$.

\emph{The case (0).} For a row $h \neq 0$, $\Der N_\varepsilon(h)\,v = \frac{\sqrt{\dm}}{\norm{h}}\big(v - \frac{\langle v, h\rangle}{\norm{h}^2}\,h\big)$, which vanishes exactly for $v \in \R\,h$. For $k = L$, $\Der\mathcal B_{>L}(h)\,v = \big(\Der N_\varepsilon(h)\,v\big)D$ and $D$ is injective on rows, so $\mathcal U_L$ is the set of all rows $h \neq 0$. Let $k < L$ and $g = g_{k+1}$. For $n \in \mathbb S$ and a row $v$,
\[
\Der\mathcal B_{>k}(n)\,v \;=\; \Der\mathcal B_{>k+1}\big(g(n)\big)\,\big(\Der g(n)\,v^\perp\big), \qquad v^\perp = v - \frac{\langle v, n\rangle}{\dm}\,n ,
\]
where $v^\perp$ is orthogonal to $n$ and vanishes exactly for $v \in \R\,n$. For $n \in \mathbb S$ consider the linear map $v \mapsto \Der g(n)\,v^\perp + \frac{\langle v, n\rangle}{\dm}\,g(n)$ of $\R^{1\times\dm}$. By (R2) it is injective exactly when $n$ is as in (R3), and its determinant is real-analytic in $n$. By (R3) the determinant does not vanish identically on the connected set $\mathbb S$, so the $n$ as in (R3) form an open dense subset of $\mathbb S$. On this subset the map $n \mapsto N_\varepsilon(g(n))$ from $\mathbb S$ to $\mathbb S$ has an injective derivative, so it is a local diffeomorphism, and the $n$ of the subset with $g(n) \in \mathcal U_{k+1}$ form an open dense subset of $\mathbb S$ again. For such $n$, $\Der\mathcal B_{>k}(n)\,v = 0$ gives $\Der g(n)\,v^\perp \in \R\,g(n)$, so $v^\perp = 0$ and $v \in \R\,n$. Let $\mathcal U_k$ be the set of the positive multiples of these $n$.
\end{proof}

\begin{lemma}[A direction at a write]\label{lem:e0atwrite}
Let $\theta$ be regular, $0 \le k \le L$, and let $A$ be a matrix with $\dm$ columns and rank $\dm$. Let $\eta_A$ be a matrix of the size of $A$ and $\eta$ a direction of the weights of the blocks $k+1, \dots, L$ and of $D$. If
\[
\Der\mathcal B_{>k}(wA)\,(w\,\eta_A) + \Der_\theta\mathcal B_{>k}(wA)\,\eta \;=\; 0 \qquad\text{for all rows } w, \text{ with } wA \neq 0 \text{ in the case (0)},
\]
then $\eta_A = AZ$ for a $\dm\times\dm$ matrix $Z$, and
\begin{equation}\label{eq:e0D}
\Der\mathcal B_{>k}(h)\,(hZ) + \Der_\theta\mathcal B_{>k}(h)\,\eta \;=\; 0 \qquad\text{for all rows } h, \text{ with } h \neq 0 \text{ in the case (0)}.
\end{equation}
\end{lemma}

\begin{proof}
Let $w_1A = 0$. Subtracting the hypothesis at $w$ and at $w + w_1$ gives $\Der\mathcal B_{>k}(wA)\,(w_1\eta_A) = 0$. For $wA \in \mathcal U_k$, \cref{lem:e0tailrank} gives $w_1\eta_A = 0$ in the cases (n) and (+), and $w_1\eta_A \in \R\,wA$ in the case (0). In the case (0), as $\mathcal U_k$ is open and $\dm \ge 2$, $wA$ takes two linearly independent values in $\mathcal U_k$, so $w_1\eta_A = 0$ as well. Let $A^+$ be a matrix with $A^+A = I$. For every row $w$, $(w - wAA^+)\,A = 0$, so $w\,\eta_A = wAA^+\eta_A$: $\eta_A = AZ$ with $Z = A^+\eta_A$. The hypothesis is now \eqref{eq:e0D} at $h = wA$, which takes every value.
\end{proof}

\smallskip
\noindent\emph{Claim $(D_k)$}, for $k = 0, \dots, L$. Let $Z$ be a $\dm\times\dm$ matrix and $\eta$ a direction of the weights of the blocks $k+1, \dots, L$ and of $D$. If \eqref{eq:e0D} holds, then $Z \in \mathfrak g$, and with $Z_k = Z$ there are $Z_{k+1}, \dots, Z_L \in \mathfrak g$ such that, for $\tanh$ and GELU blocks,
\begin{equation}\label{eq:sbDconcl}
\eta \;=\; \Big(\big\{-\widehat Z_{\ell-1}\Wread^{\ell},\ \ \Wwrite^{\ell}Z_\ell\big\}_{\ell=k+1}^{L},\ \ -\widehat Z_LD\Big).
\end{equation}
In the case (0), with $Z_\ell = Y_\ell + \beta_\ell I$ for antisymmetric matrices $Y_k, \dots, Y_L$ and numbers $\beta_k, \dots, \beta_L$, this reads
\begin{equation}\label{eq:e0Dconcl}
Z = Y_k + \beta_k I, \qquad \eta = \Big(\big\{-Y_{\ell-1}\Wread^{\ell},\ \ \Wwrite^{\ell}(Y_\ell + \beta_\ell I)\big\}_{\ell=k+1}^{L},\ \ -Y_LD\Big).
\end{equation}
For gated blocks there are, in addition, diagonal matrices $\Delta_{k+1}, \dots, \Delta_L$, with $\Delta_\ell$ of size $q_\ell$, such that
\begin{equation}\label{eq:gtDconcl}
\eta \;=\; \Big(\big\{-Z_{\ell-1}W_g^{\ell},\ \ -Z_{\ell-1}W_v^{\ell} - W_v^{\ell}\Delta_\ell,\ \ \Delta_\ell\Wwrite^{\ell} + \Wwrite^{\ell}Z_\ell\big\}_{\ell=k+1}^{L},\ \ -Z_LD\Big).
\end{equation}

\smallskip
Conversely, every $Z = Z_k \in \mathfrak g$ and $\eta$ of the form \eqref{eq:sbDconcl}, or \eqref{eq:gtDconcl} for gated blocks, satisfy \eqref{eq:e0D}: it is the derivative at $\tau = 0$ of the identity
\[
\mathcal B_{>k}\big(h\exp(\tau Z_k)\big) \;=\; \mathcal B_{>k}(h),
\]
in which the left side is taken at the weights moved by the symmetries with $G_\ell = \exp(\tau Z_\ell)$ for $\ell \ge k$, and for gated blocks with $\tau\Delta_\ell$ in the place of $\Delta_\ell$.

\begin{proposition}[The claim gives the rank]\label{prop:e0givesrank}
Let $\theta$ be regular, with no vanishing value column if the blocks are gated, and let $(D_0)$ hold at $\theta$. Then the directions $\eta \in \R^P$ with $\Der_\theta F(\theta, x)\,\eta = 0$ for all $x$, with $xE \neq 0$ in the case (0), are the directions $\eta^{Z}$ of \eqref{eq:e0eta}, or $\eta^{Z,\Delta}$ of \eqref{eq:gteta} for gated blocks, and the Jacobian of the network has the rank \eqref{eq:e0upper} on some finite batch of such inputs.
\end{proposition}

\begin{proof}
Let $\eta = (\eta_E, \eta')$ be such a direction, with $\eta_E$ its entry at $E$. As $F_\theta(x) = \mathcal B_{>0}(xE)$,
\[
\Der_\theta F(\theta, x)\,\eta \;=\; \Der\mathcal B_{>0}(xE)\,(x\,\eta_E) + \Der_\theta\mathcal B_{>0}(xE)\,\eta' .
\]
By \cref{lem:e0atwrite} with $A = E$, $\eta_E = EZ$ and \eqref{eq:e0D} holds for $k = 0$. By $(D_0)$, $\eta = \eta^{Z}$ for some $Z_0, \dots, Z_L \in \mathfrak g$ with $Z_0 = Z$, or $\eta = \eta^{Z,\Delta}$ for gated blocks. The kernels of the Jacobian on finite batches are subspaces of $\R^P$ that can only shrink when the batch grows, so on some finite batch the kernel is the space of the directions of the symmetries, which has dimension $(L+1)\dim\G + s$ by the hypotheses on $\theta$.
\end{proof}

\begin{lemma}[The readout]\label{lem:e0readout}
If $D$ has rank $\dm$, then $(D_L)$ holds.
\end{lemma}

\begin{proof}
\emph{The case (n).} Here \eqref{eq:e0D} is $hZD + h\,\eta_D = 0$ for all rows $h$, so $\eta_D = -ZD$, and $\mathfrak g$ contains every $Z$.

\emph{The case (+).} By \eqref{eq:sbDN}, \eqref{eq:e0D} is $\big(nZ - \frac{\langle nZ, n\rangle}{\dm}\,n\big)D + n\,\eta_D = 0$ for all $n$ in the ball, hence as a polynomial in $n$. Its part of degree $1$ gives $\eta_D = -ZD$. Its part of degree $3$ gives $\langle nZ, n\rangle\,nD = 0$, and $nD \neq 0$ for $n \neq 0$, so $\langle nZ, n\rangle = 0$ for all $n$: $Z$ is antisymmetric.

\emph{The case (0).} Here \eqref{eq:e0D} is $\big(\Der N_\varepsilon(h)\,(hZ)\big)D + N_\varepsilon(h)\,\eta_D = 0$. Multiplying by $\norm{h}^3/\sqrt{\dm}$ gives
\[
\norm{h}^2\;h\,(ZD + \eta_D) \;=\; \langle hZ, h\rangle\;hD \qquad\text{for all rows } h .
\]
So $h\,(ZD + \eta_D) = \mu(h)\,hD$ with a number $\mu(h)$, for every $h \neq 0$. As $h \mapsto hD$ is injective and $\dm \ge 2$, comparing $h$, $h'$ and $h + h'$ for linearly independent $h, h'$ shows that $\mu$ is a constant $\beta$, so $ZD + \eta_D = \beta D$. Then $\langle hZ, h\rangle = \beta\norm{h}^2$ for all $h$, so the symmetric part of $Z$ is $\beta I$: $Z = Y_L + \beta I$ with $Y_L$ antisymmetric, and $\eta_D = \beta D - ZD = -Y_LD$.
\end{proof}

What remains is the step from $(D_{k+1})$ to $(D_k)$. For $\tanh$ blocks in the cases (n) and (+) we prove it at almost all parameters (\cref{lem:sbsteptanh}). For GELU blocks in the case (+) and for gated blocks we prove it at explicit parameters (\cref{lem:gtstep,lem:gtstar}), and in the case (0) at suitable parameters (\cref{lem:e0steptanh,lem:e0stepgelu}).

\para{The identity for one block.} Fix $k < L$, let block $k+1$ be a $\tanh$ or GELU block, and write $W = \Wread^{k+1}$, with the columns $c_1, \dots, c_q$, $V = \Wwrite^{k+1}$, with the rows $a_1, \dots, a_q$, and $\mathcal B = \mathcal B_{>k+1}$. Let $Z$ and $\eta = (\eta_W, \eta_V, \eta')$ be as in $(D_k)$, where $\eta_W$ and $\eta_V$ are the entries of $\eta$ at $W$ and $V$ and $\eta'$ is the rest. For rows $n$, $\tilde w$ and $\tilde w'$ let
\begin{align}
\dot z(n) &= \begin{cases} nZ\,W + n\,\eta_W & \text{in the case (n),}\\[3pt] \Big(nZ - \dfrac{\langle nZ, n\rangle}{\dm}\,n\Big)\,W + n\,\eta_W & \text{in the cases (+) and (0),}\end{cases} \label{eq:e0zdot}\\[2pt]
\mathcal L(\tilde w, \tilde w'; n) &= \Der\mathcal B(\tilde wV)\,\Big(\big(\tilde w'\odot\dot z(n)\big)\,V + \tilde w\,\eta_V\Big) + \Der_\theta\mathcal B(\tilde wV)\,\eta' . \label{eq:e0L}
\end{align}
By the chain rule and \eqref{eq:sbDN}, the left side of \eqref{eq:e0D} at $h$ is $\mathcal L\big(\sigma(nW), \sigma'(nW); n\big)$ at $n = N(h)$: $\tilde w = \sigma(nW)$ is the row of the outputs of the hidden units of block $k+1$, $\tilde w' = \sigma'(nW)$ the row of their derivatives, and $\dot z(n)$ the derivative of the row $nW$ that they read. So \eqref{eq:e0D} for $k$ is equivalent to
\begin{equation}\label{eq:e0star}
\mathcal L\big(\sigma(nW),\ \sigma'(nW);\ n\big) \;=\; 0 \qquad\text{for all values } n \text{ of } N,
\end{equation}
that is, for all rows $n$ in the case (n), for the rows $n$ with $\norm{n}^2 < \dm$ in the case (+), and for $n \in \mathbb S$ in the case (0).

\begin{lemma}[The identity at every real row]\label{lem:sballrows}
In the cases (n) and (+), \eqref{eq:e0D} for $k$ holds if and only if
\begin{equation}\label{eq:sbstarall}
\mathcal L\big(\sigma(nW),\ \sigma'(nW);\ n\big) \;=\; 0 \qquad\text{for every row } n \in \R^{1\times\dm} .
\end{equation}
\end{lemma}

\begin{proof}
In the case (n) this is \eqref{eq:e0star}. In the case (+), the left side of \eqref{eq:sbstarall} is real-analytic in $n$ on all of $\R^{1\times\dm}$, because $\mathcal B$ is real-analytic on all rows. So it vanishes on the ball if and only if it vanishes everywhere.
\end{proof}

In the case (0) the identity holds only on the sphere $\mathbb S$, which contains no ray. This is where the cases part: we follow rays in the cases (n) and (+), and a great circle at complex angles in the case (0).

\subsection{Rays: \texorpdfstring{$\tanh$}{tanh} blocks}\label{app:raystanh}

\para{Rays: no normalization, and $\varepsilon > 0$.} In the cases (n) and (+) we take $n = t\,\xi$ in \eqref{eq:sbstarall}, with a row $\xi$, and let $t \to \infty$. The next lemma separates the hidden units of a $\tanh$ block.

\lemRaySep*

\begin{proof}
No $u_j$ is zero. Let $s_j$ be the sign of $u_j$ and $\zeta_j(t) = e^{-2\lvert u_j\rvert t} \in (0, 1)$, so that $\tanh(tu_j) = s_j\,(1 - \zeta_j(t))/(1 + \zeta_j(t))$. The maps $\zeta \mapsto \mathcal T_i\big(s_1\frac{1 - \zeta_1}{1 + \zeta_1}, \dots, s_q\frac{1 - \zeta_q}{1 + \zeta_q}\big)$ are real-analytic near $\zeta = 0$, so for some $\rho \in (0, 1)$ they are sums of power series $\sum_\alpha b_{i,\alpha}\,\zeta^\alpha$, over the multi-indices $\alpha = (\alpha_1, \dots, \alpha_q)$ of nonnegative integers, that converge absolutely for $\lvert\zeta_j\rvert \le \rho$. Choose $t_0 \ge 1$ with $\zeta_j(t_0) \le \rho$ for all $j$. For $t \ge t_0$ the hypothesis reads
\[
\sum_\alpha p_\alpha(t)\,e^{-\mu_\alpha t} = 0, \qquad p_\alpha(t) = \sum_{i=0}^{K} t^i\,b_{i,\alpha}, \qquad \mu_\alpha = 2\sum_{j=1}^{q}\alpha_j\,\lvert u_j\rvert .
\]
The numbers $\mu_\alpha$ are distinct for distinct $\alpha$, because the $u_j$ are linearly independent over $\mathbb Q$, and only finitely many of them lie below any bound, because every $\lvert u_j\rvert > 0$. Suppose that some $p_\alpha$ is not zero, and among these let $\alpha^*$ have the smallest $\mu_\alpha$, say $\mu^*$. Let $\gamma > 0$ be the distance from $\mu^*$ to the next larger $\mu_\alpha$. For $t \ge t_0$,
\[
\lvert p_{\alpha^*}(t)\rvert \;\le\; \sum_{\mu_\alpha > \mu^*}\lvert p_\alpha(t)\rvert\,e^{-(\mu_\alpha - \mu^*)t}
\;\le\; t^K\,e^{-\gamma(t - t_0)}\,e^{\mu^*t_0}\,\sum_{i=0}^{K}\sum_\alpha\lvert b_{i,\alpha}\rvert\,\zeta(t_0)^\alpha ,
\]
which tends to $0$ as $t \to \infty$. A polynomial that tends to $0$ is zero, a contradiction. So every $b_{i,\alpha} = 0$, and every $\mathcal T_i$ vanishes on the image of a neighborhood of $\zeta = 0$, which is a neighborhood of the row $(s_1, \dots, s_q)$. A real-analytic map on $\R^{1\times q}$ that vanishes on an open set vanishes everywhere.
\end{proof}

Consider, for a block $\ell$ with $W = \Wread^{\ell}$ and $V = \Wwrite^{\ell}$:
\begin{itemize}[leftmargin=2.9em]
\item[(T1)] no row of $V$ vanishes;
\item[(T2)] the columns of $W$ are linearly independent over $\mathbb Q$.
\end{itemize}

\begin{lemma}[The step for $\tanh$ on rays]\label{lem:sbsteptanh}
Let $\sigma = \tanh$, in the case (n) or (+), let $\theta$ be regular, and let block $k+1$ satisfy (T1) and (T2). Then $(D_{k+1})$ implies $(D_k)$.
\end{lemma}

\begin{proof}
Let $Z$ and $\eta$ satisfy \eqref{eq:e0D} for $k$, so that \eqref{eq:sbstarall} holds, with the notation above. Take $n = t\,\xi$ with a row $\xi$ and $t > 0$. Then $\tilde w = \tanh(t\,\xi W)$ and $\tilde w' = 1 - \tilde w^2$, the row with the entries $1 - \tilde w_j^2$. By \eqref{eq:e0zdot},
\[
\dot z(t\,\xi) \;=\; t\;\xi\,(ZW + \eta_W) \;-\; t^3\,\frac{\langle\xi Z, \xi\rangle}{\dm}\;\xi W
\]
in the case (+), and without the term in $t^3$ in the case (n). So \eqref{eq:sbstarall} at $n = t\,\xi$ reads
\[
t^3\,\mathcal T_3(\tilde w) + t\,\mathcal T_1(\tilde w) + \mathcal T_0(\tilde w) = 0 \qquad\text{at } \tilde w = \tanh(t\,\xi W),
\]
with the maps
\begin{align*}
\mathcal T_3(\tilde w) &= -\frac{\langle\xi Z, \xi\rangle}{\dm}\;\Der\mathcal B(\tilde wV)\,\Big(\big((1 - \tilde w^2)\odot\xi W\big)\,V\Big),\\
\mathcal T_1(\tilde w) &= \Der\mathcal B(\tilde wV)\,\Big(\big((1 - \tilde w^2)\odot\xi\,(ZW + \eta_W)\big)\,V\Big),\\
\mathcal T_0(\tilde w) &= \Der\mathcal B(\tilde wV)\,(\tilde w\,\eta_V) + \Der_\theta\mathcal B(\tilde wV)\,\eta' ,
\end{align*}
where $\mathcal T_3 = 0$ in the case (n). They are defined and real-analytic for every row $\tilde w \in \R^{1\times q}$. By (T2), for integers $\kappa_1, \dots, \kappa_q$ that are not all zero, $\sum_j\kappa_j\,\xi c_j = \xi\sum_j\kappa_jc_j$ vanishes only for $\xi$ in a hyperplane. So for $\xi$ outside countably many hyperplanes the numbers $\xi c_1, \dots, \xi c_q$ are linearly independent over $\mathbb Q$, and \cref{lem:raysep} gives $\mathcal T_3 = \mathcal T_1 = \mathcal T_0 = 0$ on $\R^{1\times q}$.

We use the following fact: if a row $b \in \R^{1\times q}$ has $\Der\mathcal B(\tilde wV)\big(((1 - \tilde w^2)\odot b)\,V\big) = 0$ for all rows $\tilde w$, then $b = 0$. Indeed, by \cref{lem:e0tailrank} the row $\sum_j b_j\,(1 - \tilde w_j^2)\,a_j$ vanishes for all $\tilde w$ with $\tilde wV \in \mathcal U_{k+1}$. These $\tilde w$ form an open dense set, because $V$ has rank $\dm$, so the row vanishes as a polynomial in $\tilde w$. Its coefficient of $\tilde w_j^2$ is $-b_ja_j$, and $a_j \neq 0$ by (T1).

\emph{(a) $\mathcal T_3 = 0$ gives $Z$.} In the case (n) every matrix $Z$ lies in $\mathfrak g$. In the case (+), $\xi W \neq 0$, so the fact gives $\langle\xi Z, \xi\rangle = 0$ for a dense set of rows $\xi$, hence for all: $Z$ is antisymmetric. In both cases $Z \in \mathfrak g$ and $\widehat Z = Z$.

\emph{(b) $\mathcal T_1 = 0$ gives $\eta_W$.} The fact gives $\xi\,(ZW + \eta_W) = 0$ for a dense set of rows $\xi$, so $\eta_W = -ZW$.

\emph{(c) $\mathcal T_0 = 0$ gives the rest.} $\Der\mathcal B(\tilde wV)\,(\tilde w\,\eta_V) + \Der_\theta\mathcal B(\tilde wV)\,\eta' = 0$ for all rows $\tilde w$. By \cref{lem:e0atwrite} with $A = V$, $\eta_V = VZ'$ and \eqref{eq:e0D} holds for $k+1$ with $Z'$ and $\eta'$. By $(D_{k+1})$, $Z' \in \mathfrak g$ and $\eta'$ is of the form \eqref{eq:sbDconcl} with $Z_{k+1} = Z'$. With (a) and (b) this is \eqref{eq:sbDconcl} for $k$.
\end{proof}

\subsection{Rays: GELU and gated blocks}\label{app:raysgelu}

\para{GELU and the Mills ratio.} GELU is $\sigma(z) = \frac z2\big(1 + \operatorname{erf}(z/\sqrt2)\big)$, which is entire. For $z \in \mathbb C$ let
\[
\varphi(z) = \frac{e^{-z^2/2}}{\sqrt{2\pi}}, \qquad \operatorname{mr}(z) \;=\; \sqrt{\frac{\pi}{2}}\;e^{z^2/2}\,\Big(1 - \operatorname{erf}\frac{z}{\sqrt2}\Big),
\]
the normal density and the Mills ratio.

\begin{lemma}[The Mills ratio]\label{lem:e0mills}
\stpart{lem:e0mills}{i}{} For all $z \in \mathbb C$: $\sigma(z) = z - \varphi(z)\,z\,\operatorname{mr}(z)$ and $\sigma'(z) = 1 - \varphi(z)\,\big(\operatorname{mr}(z) - z\big)$.\\
\stpart{lem:e0mills}{ii}{} For $\operatorname{Re}z > 0$ and every integer $m \ge 0$, with $(-1)!! = 1$,
\[
\Big\lvert\,\operatorname{mr}(z) - \sum_{l=0}^{m-1}\frac{(-1)^l\,(2l-1)!!}{z^{2l+1}}\,\Big\rvert \;\le\; \frac{(2m-1)!!}{(\operatorname{Re}z)^{2m+1}} .
\]
\end{lemma}

\begin{proof}
(i) follows from $\varphi(z)\,\operatorname{mr}(z) = \frac12\big(1 - \operatorname{erf}(z/\sqrt2)\big)$ and $\sigma'(z) = \frac12\big(1 + \operatorname{erf}(z/\sqrt2)\big) + z\,\varphi(z)$. For (ii): for real $z$, $\operatorname{mr}(z) = e^{z^2/2}\int_z^\infty e^{-y^2/2}\,dy$, and the substitution $y = z + u$ gives $\operatorname{mr}(z) = \int_0^\infty e^{-zu - u^2/2}\,du$. Both sides are entire, so this holds for all $z$. For a number $y \ge 0$, $\big\lvert e^{-y} - \sum_{l<m}\frac{(-y)^l}{l!}\big\rvert \le \frac{y^m}{m!}$. With $y = u^2/2$, and with $\int_0^\infty e^{-zu}u^{2l}\,du = \frac{(2l)!}{z^{2l+1}}$ and $\frac{(2l)!}{2^l\,l!} = (2l-1)!!$, this gives (ii).
\end{proof}

We say that a map $\Theta$ on an interval $(t_0, \infty)$ is \emph{meromorphic at infinity} if $\Theta(t) = t^{l}\,\Theta_0(1/t)$ for an integer $l$ and a map $\Theta_0$ that is holomorphic near $0$. Such a map vanishes identically if it tends to $0$ faster than every power of $1/t$.

\lemMills*

\begin{proof}
Suppose that a component of $\Theta_1$ is not zero, and let $\Theta_1, \Theta_2$ stand for this component of both. Then $\operatorname{mr}(z(t)) = -\Theta_2(t)/\Theta_1(t) + O(t^{-m})$ for every $m$. For large $t$, $t = \frac{z}{2\omega}\big(1 + \sqrt{1 - 4\omega\omega'/z^2}\big)$ with $z = z(t)$, so $1/t$ is a convergent power series in $1/z$ without constant term. Hence $-\Theta_2/\Theta_1$ is, for large $t$, a convergent series in powers of $1/z$ with finitely many positive powers of $z$, and this series is an asymptotic expansion of $\operatorname{mr}(z)$ along the curve $z = z(t)$. On this curve $\operatorname{Re}z$ is at least a positive constant times $\lvert z\rvert$ for large $t$, so by \cpref{lem:e0mills}{ii} the series $\sum_l(-1)^l(2l-1)!!\,z^{-2l-1}$ is one as well. Asymptotic expansions are unique, so the two series have the same coefficients. But the numbers $(2l-1)!!$ grow faster than every geometric sequence, so the second series converges for no $z$, a contradiction. Hence $\Theta_1 = 0$, then $\Theta_2$ tends to $0$ faster than every power, and $\Theta_2 = 0$.
\end{proof}

\para{A gate on a ray.} GELU and SiLU both satisfy $\sigma(z) - \sigma(-z) = z$, and so $\sigma'(z) + \sigma'(-z) = 1$. Let $\delta(y) = y - \sigma(y)$, with the derivative $\delta'(y) = 1 - \sigma'(y)$. For every real number $y > 0$,
\begin{equation}\label{eq:sbunit}
\sigma(y) = y - \delta(y), \qquad \sigma(-y) = -\delta(y), \qquad \sigma'(y) = 1 - \delta'(y), \qquad \sigma'(-y) = \delta'(y) .
\end{equation}
By \cpref{lem:e0mills}{i}, for GELU $\delta(y) = \varphi(y)\,y\,\operatorname{mr}(y)$ and $\delta'(y) = \varphi(y)\,\big(\operatorname{mr}(y) - y\big)$, and for SiLU
\[
\delta(y) = \frac{y}{1 + e^{y}}, \qquad \delta'(y) = \frac{1 + e^{y} - y\,e^{y}}{(1 + e^{y})^2} .
\]
Let $p = 2$ for GELU and $p = 1$ for SiLU. For every real number $y \ge 2$,
\begin{equation}\label{eq:gtdelta}
\tfrac16\,e^{-y^p/p} \;\le\; \delta(y) \;\le\; y\,e^{-y^p/p}, \qquad \lvert\delta'(y)\rvert \;\le\; 4\,y\,\delta(y) .
\end{equation}
For GELU, \cpref{lem:e0mills}{ii} gives $\operatorname{mr}(y) \le 1/y$ and $\lvert\operatorname{mr}(y) - 1/y\rvert \le 1/y^3$, so $\frac34 \le y\operatorname{mr}(y) \le 1$ and $\lvert\operatorname{mr}(y) - y\rvert \le 2y$. For SiLU, $\frac12\,y\,e^{-y} \le \delta(y) \le y\,e^{-y}$ and $\lvert\delta'(y)\rvert \le (2 + y)\,e^{-y}$. So on a ray a GELU hidden unit, and the gate of a hidden unit of a gated block, differ from a linear or a vanishing one by the small term $\delta$.

\begin{lemma}[The rest of the network on a ray]\label{lem:sbcoeff}
In the case (+), let $k \le L$, let $h^\circ \neq 0$ be a row, let $v_1, \dots, v_i$ be rows, and let $\eta$ be a direction of the weights of the blocks $k+1, \dots, L$ and of $D$. With $\Der^i$ the derivative of order $i$ in $h$, the maps
\[
t \;\mapsto\; t^{\,i}\;\Der^{i}\mathcal B_{>k}(t\,h^\circ)\,(v_1, \dots, v_i) \qquad\text{and}\qquad t \;\mapsto\; t^{\,i}\;\Der^{i}\big(\Der_\theta\mathcal B_{>k}(\cdot)\,\eta\big)(t\,h^\circ)\,(v_1, \dots, v_i)
\]
are real-analytic functions of $1/t$ near $1/t = 0$, so they are meromorphic at infinity. Moreover, there are numbers $\varrho, t_0 > 0$ such that the maps $u \mapsto \mathcal B_{>k}\big(t\,(h^\circ + u)\big)$ and $u \mapsto \Der_\theta\mathcal B_{>k}\big(t\,(h^\circ + u)\big)\,\eta$ extend to holomorphic maps on the rows $u \in \mathbb C^{1\times\dm}$ with $\lvert u_j\rvert \le \varrho$ for all $j$, which are bounded uniformly in $t \ge t_0$.
\end{lemma}

\begin{proof}
$\mathcal B_{>k} = \mathcal B^\sharp\circ N_\varepsilon$, with $\mathcal B^\sharp(n) = nD$ if $k = L$ and $\mathcal B^\sharp = \mathcal B_{>k+1}\circ g_{k+1}$ otherwise, which is real-analytic on all rows. For $t > 0$ and a row $u$,
\[
N_\varepsilon\big(t\,(h^\circ + u)\big) \;=\; \frac{h^\circ + u}{\sqrt{\norm{h^\circ + u}^2/\dm + \varepsilon\,t^{-2}}} ,
\]
and the right side is a real-analytic function of $(1/t, u)$ near $(0, 0)$, also at $1/t = 0$, because $h^\circ \neq 0$. So $\mathcal B_{>k}\big(t\,(h^\circ + u)\big)$ is a real-analytic function of $(1/t, u)$ near $(0, 0)$. It extends to a holomorphic and bounded function on a complex neighborhood of $(0, 0)$, which gives $\varrho$ and $t_0$. A derivative of order $i$ in $v = t\,u$ at $v = 0$ is $t^{-i}$ times the derivative of order $i$ in $u$ at $u = 0$. The weights enter analytically as well, which gives the same for $\Der_\theta\mathcal B_{>k}(\cdot)\,\eta$.
\end{proof}

\para{GELU and gated blocks in one form.} The rest of the treatment of rays is for GELU, SwiGLU and GeGLU blocks. A GELU block has the form of a GeGLU block in which the value $ne_j$ of every hidden unit is replaced by the number $1$. So in a GELU block let $W_g^{\ell} = \Wread^{\ell}$, and call $c_1, \dots, c_{q_\ell}$ its gate columns. For a row $n$ let $\iota^{\ell}(n) \in \R^{1\times q_\ell}$ be the row of the values of its hidden units:
\begin{gather*}
\iota^{\ell}(n) = n\,W_v^{\ell} \ \text{ in a gated block}, \qquad\qquad \iota^{\ell}(n) = (1, \dots, 1) \ \text{ in a GELU block},\\
g_\ell(n) \;=\; \big(\sigma(n\,W_g^{\ell})\odot\iota^{\ell}(n)\big)\,\Wwrite^{\ell} \ \text{ in both.}
\end{gather*}
Let $m = 2$ for gated blocks and $m = 1$ for GELU blocks, so that $\iota^{\ell}(\lambda n) = \lambda^{m-1}\,\iota^{\ell}(n)$ for every number $\lambda > 0$. The exponent $p$ of \eqref{eq:gtdelta} is $2$ for GELU and GeGLU blocks and $1$ for SwiGLU blocks.

\para{Cones and good parameters.} A \emph{cone} is a set of rows that contains $\lambda h$ for each of its rows $h$ and every number $\lambda > 0$. Let $\mathcal C_0, \dots, \mathcal C_L$ be open convex cones in $\R^{1\times\dm}$ that do not contain $0$, one for each hidden state. We call $\theta$ \emph{good} for these cones if (R1) holds and every block $\ell$, with its gate columns $c_j$, its write rows $a_j$ and the entries $\iota_j(n)$ of $\iota^{\ell}(n)$, satisfies:
\begin{itemize}[leftmargin=2.9em]
\item[(K1)] $nc_j \neq 0$ for every $n \in \mathcal C_{\ell-1}$ and every $j$. So every gate has one sign on $\mathcal C_{\ell-1}$. Let $\mathcal A_\ell$ be the set of the hidden units $j$ with $nc_j > 0$ on $\mathcal C_{\ell-1}$, the \emph{active} ones, and let
\[
g^\circ_\ell(n) \;=\; \sum_{j\in\mathcal A_\ell}(nc_j)\,\iota_j(n)\,a_j
\]
be the block with its gates \emph{frozen}. It is a linear map for a GELU block and a quadratic map for a gated block: $g^\circ_\ell(\lambda n) = \lambda^m\,g^\circ_\ell(n)$.
\item[(K2)] $g^\circ_\ell(n) \in \mathcal C_\ell$ for every $n \in \mathcal C_{\ell-1}$, and $\Der g^\circ_\ell(n)$ is invertible at some row $n$.
\item[(K3)] No $a_j$ vanishes, and in a gated block no value column $e_j$ vanishes. For GELU and GeGLU blocks the matrices $c_1c_1^\top, \dots, c_{q_\ell}c_{q_\ell}^\top$ are linearly independent over $\mathbb Q$, and for SwiGLU blocks the columns $c_1, \dots, c_{q_\ell}$ are.
\end{itemize}
The determinant of $\Der g^\circ_\ell(n)$ is a polynomial in $n$, so by (K2) $g^\circ_\ell$ is a local diffeomorphism on an open dense set of rows. A good $\theta$ is regular. Indeed, take $\xi \in \mathcal C_{\ell-1}$ with $\Der g^\circ_\ell(\xi)$ invertible. For a row $v$, by \eqref{eq:sbunit} and \eqref{eq:gtdelta}, in a gated block
\begin{align*}
\frac1t\,\Der g_\ell(t\,\xi)\,v \;&=\; \sum_j\Big[\sigma'(t\,\xi c_j)\,(vc_j)\,(\xi e_j) + \frac{\sigma(t\,\xi c_j)}{t}\,(ve_j)\Big]\,a_j \\
&\longrightarrow\; \sum_{j\in\mathcal A_\ell}\Big[(vc_j)\,(\xi e_j) + (\xi c_j)\,(ve_j)\Big]\,a_j \;=\; \Der g^\circ_\ell(\xi)\,v \qquad\text{as } t \to \infty ,
\end{align*}
and in a GELU block $\Der g_\ell(t\,\xi)\,v = \sum_j\sigma'(t\,\xi c_j)\,(vc_j)\,a_j$ tends to $\sum_{j\in\mathcal A_\ell}(vc_j)\,a_j = \Der g^\circ_\ell(\xi)\,v$. So $\Der g_\ell(t\,\xi)$ is invertible for large $t$, which is (R3).

\para{The frozen network.} In the case (n), for gated blocks and a good $\theta$, let $\mathcal B^\circ_{>k}$ be the network after $h_k$ with the gates of all its blocks frozen:
\[
\mathcal B^\circ_{>L}(h) = hD, \qquad \mathcal B^\circ_{>k-1}(h) = \mathcal B^\circ_{>k}\big(g^\circ_k(h)\big).
\]
It is a polynomial map of $h$ and of the weights. \Cref{lem:e0tailrank,lem:e0atwrite} hold for $\mathcal B^\circ_{>k}$ in the place of $\mathcal B_{>k}$, with an open dense set $\mathcal U^\circ_k$ of rows and with the same proofs, in which $g^\circ_{k+1}$ takes the place of $g_{k+1}$.

\para{The identity for one GELU or gated block.} Fix $k < L$ and write $W_g = W_g^{k+1}$, with the columns $c_1, \dots, c_q$, $V = \Wwrite^{k+1}$, with the rows $a_1, \dots, a_q$, $\iota = \iota^{k+1}$ and $\mathcal B = \mathcal B_{>k+1}$, and in a gated block $W_v = W_v^{k+1}$, with the columns $e_1, \dots, e_q$. Let $Z$ and $\eta = (\eta_g, \eta_v, \eta_V, \eta')$ be as in $(D_k)$, where $\eta_g$, $\eta_v$ and $\eta_V$ are the entries of $\eta$ at $W_g$, $W_v$ and $V$ and $\eta'$ is the rest; in a GELU block there is no $\eta_v$. For a row $n$ let $\tilde n = nZ$ in the case (n) and $\tilde n = nZ - \frac{\langle nZ, n\rangle}{\dm}\,n$ in the case (+). For rows $n$ and $\tilde w, \tilde w' \in \R^{1\times q}$ let
\begin{align}
\dot z^g(n) &= \tilde n\,W_g + n\,\eta_g, \qquad\qquad \dot z^v(n) = \begin{cases} \tilde n\,W_v + n\,\eta_v & \text{in a gated block,}\\ 0 & \text{in a GELU block,}\end{cases} \label{eq:gtzdot}\\[3pt]
\mathcal L(\tilde w, \tilde w'; n) &= \Der\mathcal B\big(\hat wV\big)\,\Big(\big(\tilde w'\odot\dot z^g(n)\odot\iota(n) + \tilde w\odot\dot z^v(n)\big)\,V + \hat w\,\eta_V\Big) + \Der_\theta\mathcal B\big(\hat wV\big)\,\eta' , \label{eq:gtL}
\end{align}
where $\hat w = \tilde w\odot\iota(n)$, and $\dot z^g_j(n)$ and $\dot z^v_j(n)$ are the entries of the rows \eqref{eq:gtzdot}. Here $\tilde w$ stands for the row $\sigma(nW_g)$ of the gates, $\tilde w'$ for the row $\sigma'(nW_g)$ of their derivatives, and $\hat w$ for the row of the outputs of the hidden units, and $\dot z^g(n)$ and $\dot z^v(n)$ are the derivatives of the rows $nW_g$ and $\iota(n)$. In a GELU block $\hat w = \tilde w$, and these are the maps \eqref{eq:e0zdot} and \eqref{eq:e0L}. By the chain rule and \eqref{eq:sbDN}, the left side of \eqref{eq:e0D} at $h$ is $\mathcal L\big(\sigma(nW_g), \sigma'(nW_g); n\big)$ at $n = N(h)$. As in \cref{lem:sballrows}, \eqref{eq:e0D} for $k$ therefore holds if and only if
\begin{equation}\label{eq:gtstarall}
\mathcal L\big(\sigma(nW_g),\ \sigma'(nW_g);\ n\big) \;=\; 0 \qquad\text{for every row } n \in \R^{1\times\dm} .
\end{equation}
We use two properties of $\mathcal L$. The row $\tilde w'$ enters only through the terms $\tilde w'_j\,\dot z^g_j(n)\,\iota_j(n)\,a_j$, so
\begin{equation}\label{eq:gtLprime}
\frac{\partial\mathcal L}{\partial\tilde w'_j}(\tilde w, \tilde w'; n) \;=\; \dot z^g_j(n)\;\iota_j(n)\;\Der\mathcal B\big(\hat wV\big)\,a_j .
\end{equation}
And in a gated block $\dot z^g$ and $\dot z^v$ are odd in $n$, while $\hat w$ does not change when $\tilde w$ and $n$ both change their sign, so
\begin{equation}\label{eq:gtparity}
\mathcal L(-\tilde w, \tilde w'; -n) \;=\; \mathcal L(\tilde w, \tilde w'; n) \qquad\text{in a gated block.}
\end{equation}

\begin{lemma}[The network after $h_k$ far out in its cone]\label{lem:gtfar}
In the case (n), for gated blocks, let $\theta$ be good, $k \le L$, $h^\circ \in \mathcal C_k$, and let $\eta$ be a direction of the weights of the blocks $k+1, \dots, L$ and of $D$. There are numbers $\varrho, \kappa, t_0 > 0$ such that, for $t \ge t_0$, the maps $u \mapsto \mathcal B_{>k}\big(t\,(h^\circ + u)\big)$ and $u \mapsto \Der_\theta\mathcal B_{>k}\big(t\,(h^\circ + u)\big)\,\eta$ extend to holomorphic maps on the rows $u \in \mathbb C^{1\times\dm}$ with $\lvert u_j\rvert \le \varrho$ for all $j$, and for these $u$, at $h = t\,(h^\circ + u)$,
\[
\norm{\mathcal B_{>k}(h) - \mathcal B^\circ_{>k}(h)} \;\le\; e^{-\kappa t^p}, \qquad \norm{\Der_\theta\mathcal B_{>k}(h)\,\eta - \Der_\theta\mathcal B^\circ_{>k}(h)\,\eta} \;\le\; e^{-\kappa t^p} .
\]
\end{lemma}

\begin{proof}
\emph{One gate.} Let $z \in \mathbb C$ with $\operatorname{Re}z \ge 2$ and $\lvert\operatorname{Im}z\rvert \le \frac12\operatorname{Re}z$, so that $\operatorname{Re}(z^2) \ge \frac34(\operatorname{Re}z)^2$ and $\lvert z\rvert \le 2\operatorname{Re}z$. For GELU, by \cref{lem:e0mills}, $\sigma(z) - z = -\varphi(z)\,z\operatorname{mr}(z)$ and $\lvert z\operatorname{mr}(z)\rvert \le \lvert z\rvert/\operatorname{Re}z \le 2$, so $\lvert\sigma(z) - z\rvert \le e^{-\frac38(\operatorname{Re}z)^2}$. For SiLU, $\sigma(z) - z = -z/(1 + e^{z})$ and $\lvert 1 + e^{z}\rvert \ge \frac12\,e^{\operatorname{Re}z}$, so $\lvert\sigma(z) - z\rvert \le 4\operatorname{Re}z\;e^{-\operatorname{Re}z}$. Both gates are holomorphic at $z$ and at $-z$, and $\sigma(-z) = \sigma(z) - z$. So there is a number $\kappa_1 > 0$ such that, for large $\operatorname{Re}z$,
\begin{equation}\label{eq:gtgate}
\lvert\sigma(z) - z\rvert \;\le\; e^{-\kappa_1(\operatorname{Re}z)^p} \qquad\text{and}\qquad \lvert\sigma(-z)\rvert \;\le\; e^{-\kappa_1(\operatorname{Re}z)^p} .
\end{equation}

\emph{The first bound, by induction on $L - k$.} We prove it at the weights $\theta + \tau\eta$ as well, for all complex numbers $\tau$ with $\lvert\tau\rvert \le \tau_0$ and a small $\tau_0 > 0$, where the hidden units whose gates are frozen are those of $\theta$. For $k = L$, $\mathcal B_{>L}(h) = hD = \mathcal B^\circ_{>L}(h)$. Let the bound hold for some $k \ge 1$, and let $h^\circ \in \mathcal C_{k-1}$. By (K1), $h^\circ c_j \neq 0$ for the gate columns $c_j$ of block $k$ at $\theta$. Let $\varrho$ and $\tau_0$ be small, and let $h = t\,(h^\circ + u)$. For every $j$, with $c_j$ taken at the weights $\theta + \tau\eta$, the number $z = hc_j$ or $z = -hc_j$ with the sign of $h^\circ c_j$ satisfies $\operatorname{Re}z \ge \frac t2\,\lvert h^\circ c_j\rvert$ and $\lvert\operatorname{Im}z\rvert \le \frac12\operatorname{Re}z$. By \eqref{eq:gtgate}, $g_k$ is holomorphic at these $h$, and
\[
\norm{g_k(h) - g^\circ_k(h)} \;\le\; e^{-\kappa_2t^p}, \qquad g^\circ_k(h) \;=\; t^2\,\big(g^\circ_k(h^\circ) + u'\big),
\]
with a number $\kappa_2 > 0$, where $g_k$ and $g^\circ_k$ are taken at the weights $\theta + \tau\eta$ and $g^\circ_k(h^\circ)$ at $\theta$, and where the entries of the row $u'$ are as small as we like when $\varrho$ and $\tau_0$ are small. So $g_k(h)$ is $t^2\,\big(g^\circ_k(h^\circ) + u''\big)$ with a row $u''$ of the same kind. By (K2), $g^\circ_k(h^\circ) \in \mathcal C_k$. So the bound for $k$, at the row $g^\circ_k(h^\circ)$ and with $t^2$ in the place of $t$, gives $\norm{\mathcal B_{>k}(g_k(h)) - \mathcal B^\circ_{>k}(g_k(h))} \le e^{-\kappa t^{2p}}$. The derivative of the polynomial map $\mathcal B^\circ_{>k}$ is at most a fixed power of $t$ at the rows $t^2\,(g^\circ_k(h^\circ) + u')$, so $\norm{\mathcal B^\circ_{>k}(g_k(h)) - \mathcal B^\circ_{>k}(g^\circ_k(h))}$ is at most a fixed power of $t$ times $e^{-\kappa_2t^p}$. As $\mathcal B_{>k-1} = \mathcal B_{>k}\circ g_k$ and $\mathcal B^\circ_{>k-1} = \mathcal B^\circ_{>k}\circ g^\circ_k$, this gives the bound for $k - 1$, with a smaller $\kappa$.

\emph{The second bound.} At $h = t\,(h^\circ + u)$ the difference of $\mathcal B_{>k}(h)$ and $\mathcal B^\circ_{>k}(h)$ at the weights $\theta + \tau\eta$ is holomorphic in $\tau$ for $\lvert\tau\rvert \le \tau_0$ and at most $e^{-\kappa t^p}$. By Cauchy's estimate its derivative at $\tau = 0$, which is $\Der_\theta\mathcal B_{>k}(h)\,\eta - \Der_\theta\mathcal B^\circ_{>k}(h)\,\eta$, is at most $e^{-\kappa t^p}/\tau_0$.
\end{proof}

\begin{lemma}[The step for GELU and gated blocks on rays]\label{lem:gtstep}
Let the blocks be GELU blocks in the case (+), or gated blocks in the case (n) or (+). Let $\theta$ be good for the cones $\mathcal C_0, \dots, \mathcal C_L$, and let $k < L$. Then $(D_{k+1})$ implies $(D_k)$.
\end{lemma}

\begin{proof}
Let $Z$ and $\eta$ satisfy \eqref{eq:e0D} for $k$, so that \eqref{eq:gtstarall} holds, with the notation above. Let $\mathcal A = \mathcal A_{k+1}$ and $g^\circ = g^\circ_{k+1}$. Let $s_j$ be the sign of $nc_j$ on $\mathcal C_k$ and $\bar s = (s_1, \dots, s_q)$, let $I_{\mathcal A}$ be the diagonal matrix with the entries $1$ at the hidden units of $\mathcal A$ and $0$ elsewhere, and $1_{\mathcal A}$ the row with these entries. Let
\[
\mathcal B^\infty = \mathcal B_{>k+1} \ \text{ in the case (+)}, \qquad \mathcal B^\infty = \mathcal B^\circ_{>k+1} \ \text{ in the case (n)},
\]
and let $\mathcal L^\infty$ be \eqref{eq:gtL} with $\mathcal B^\infty$ in the place of $\mathcal B$. It is real-analytic in $(\tilde w, \tilde w', n)$ on all rows, a polynomial in the case (n), and it satisfies \eqref{eq:gtLprime}, and \eqref{eq:gtparity} in a gated block, with $\mathcal B^\infty$. \Cref{lem:e0tailrank,lem:e0atwrite} hold for $\mathcal B^\infty$, and we write $\mathcal U^\infty$ for its open dense set $\mathcal U_{k+1}$ or $\mathcal U^\circ_{k+1}$.

\emph{The ray.} Let $\xi \in \mathcal C_k$ and $n = t\,\xi$ with $t > 0$. Let
\[
r_j = \lvert\xi c_j\rvert^p, \qquad \delta_j(t) = \delta\big(t\,\lvert\xi c_j\rvert\big), \qquad \delta'_j(t) = \delta'\big(t\,\lvert\xi c_j\rvert\big),
\]
and let $\delta(t)$ and $\delta'(t)$ be the rows with these entries. By \eqref{eq:gtdelta}, with a constant $\kappa$, for large $t$
\begin{equation}\label{eq:gtsizes}
\tfrac16\,e^{-r_jt^p/p} \;\le\; \delta_j(t) \;\le\; \kappa\,t\;e^{-r_jt^p/p}, \qquad \lvert\delta'_j(t)\rvert \;\le\; \kappa\,t\;\delta_j(t) .
\end{equation}
By \eqref{eq:sbunit}, $\sigma(nW_g) = n\,W_gI_{\mathcal A} - \delta(t)$ and $\sigma'(nW_g) = 1_{\mathcal A} - \bar s\odot\delta'(t)$. So \eqref{eq:gtstarall} at $n = t\,\xi$ is
\begin{equation}\label{eq:gtray}
\mathcal L\big(t\,\xi W_gI_{\mathcal A} - \delta(t),\ \ 1_{\mathcal A} - \bar s\odot\delta'(t);\ \ t\,\xi\big) \;=\; 0 \qquad\text{for all } t > 0 .
\end{equation}
Without the small terms the hidden units of the block give the row $t\,\xi W_gI_{\mathcal A}\odot\iota(t\,\xi)$, and the block gives $g^\circ(t\,\xi) = t^m\,h^\circ$ with $h^\circ = g^\circ(\xi)$, which lies in $\mathcal C_{k+1}$ by (K2), so $h^\circ \neq 0$. With a row $\delta$ in the place of $\delta(t)$ the block gives $t^m\,(h^\circ + u)$ with $u = -\frac1t\,\big(\delta\odot\iota(\xi)\big)\,V$.

\emph{The expansion.} For a multi-index $\alpha = (\alpha_1, \dots, \alpha_q)$ of nonnegative integers let $\delta^\alpha = \prod_j\delta_j^{\alpha_j}$, $\lvert\alpha\rvert = \sum_j\alpha_j$, $\alpha! = \prod_j\alpha_j!$ and $r_\alpha = \sum_j\alpha_jr_j$, and let $1_j$ be the multi-index with the entry $1$ at $j$ and $0$ elsewhere. Let
\[
\mathcal L_\alpha(t) = \frac{1}{\alpha!}\,\frac{\partial^{\lvert\alpha\rvert}\mathcal L^\infty}{\partial\tilde w^\alpha}, \qquad \mathcal L'_{\alpha,i}(t) = \frac{1}{\alpha!}\,\frac{\partial^{\lvert\alpha\rvert}}{\partial\tilde w^\alpha}\,\frac{\partial\mathcal L^\infty}{\partial\tilde w'_i}, \qquad\text{both at } \big(t\,\xi W_gI_{\mathcal A},\ 1_{\mathcal A};\ t\,\xi\big).
\]
We write $\mathcal L_j = \mathcal L_{1_j}$ and $\mathcal L'_j = \mathcal L'_{0,j}$. As $\mathcal L^\infty$ is real-analytic in $\tilde w$ and affine in $\tilde w'$,
\begin{equation}\label{eq:gtseries}
\mathcal L^\infty\big(t\,\xi W_gI_{\mathcal A} - \delta,\ \ 1_{\mathcal A} - \bar s\odot\delta';\ \ t\,\xi\big) \;=\; \sum_\alpha(-1)^{\lvert\alpha\rvert}\,\delta^\alpha\,\Big[\mathcal L_\alpha(t) - \sum_{i=1}^{q}s_i\,\mathcal L'_{\alpha,i}(t)\,\delta'_i\Big]
\end{equation}
for small rows $\delta$ and all rows $\delta'$. By \eqref{eq:gtsizes}, at $\delta = \delta(t)$ and $\delta' = \delta'(t)$ the term of $\delta^\alpha$ is of the order of $e^{-r_\alpha t^p/p}$, up to powers of $t$, and the term of $\delta^\alpha\delta'_i$ of the order of $e^{-(r_\alpha + r_i)t^p/p}$. We call $r_\alpha$ and $r_\alpha + r_i$ the \emph{exponents} of these terms. We use three facts.
\begin{itemize}[leftmargin=2.2em]
\item[(a)] Every $\mathcal L_\alpha$ and $\mathcal L'_{\alpha,i}$ is a real-analytic function of $t \in \R$ that is meromorphic at infinity. In the case (n) it is a polynomial in $t$. In the case (+) it is built from derivatives of $\mathcal B$ and of $\Der_\theta\mathcal B(\cdot)\,\eta'$ at $t^m\,h^\circ$ and from polynomials in $t$, so this is \cref{lem:sbcoeff} with $t^m$ in the place of $t$.
\item[(b)] At $\delta = \delta(t)$ and $\delta' = \delta'(t)$, the left side of \eqref{eq:gtseries} vanishes in the case (+), by \eqref{eq:gtray}. In the case (n), where the blocks are gated and $m = 2$, it is at most $e^{-\kappa' t^{2p}}$ for large $t$, with a number $\kappa' > 0$. Indeed, by \cref{lem:gtfar} for $k+1$, at the row $h^\circ$ and with $t^2$ in the place of $t$, the maps $\mathcal B - \mathcal B^\infty$ and $\Der_\theta\mathcal B(\cdot)\,\eta' - \Der_\theta\mathcal B^\infty(\cdot)\,\eta'$ are holomorphic and at most $e^{-2\kappa' t^{2p}}$ at the rows $t^2\,(h^\circ + u)$ with $\lvert u_j\rvert \le \varrho$, and by Cauchy's estimate so is the derivative of the first, for large $t$, where $\lvert u_j\rvert \le \varrho/2$. At $\delta = \delta(t)$ the block gives such a row, with $u \to 0$, and the rows to which these derivatives are applied in \eqref{eq:gtL} are at most a fixed power of $t$. So $\mathcal L$ and $\mathcal L^\infty$ differ at the point of \eqref{eq:gtray} by at most a fixed power of $t$ times $e^{-2\kappa' t^{2p}}$, and $\mathcal L$ vanishes there.
\item[(c)] For large $t$ the series \eqref{eq:gtseries} converges absolutely at $\delta(t)$, $\delta'(t)$, and the sum of any set of its terms is at most a fixed power of $t$ times $e^{-rt^p/p}$, where $r$ is the smallest exponent of a term of the set. In the case (n) the series is a finite sum with coefficients that are polynomials in $t$, and this follows from \eqref{eq:gtsizes}. In the case (+), by \cref{lem:sbcoeff} with $t^m$ in the place of $t$, the maps $u \mapsto \mathcal B\big(t^m(h^\circ + u)\big)$ and $u \mapsto \Der_\theta\mathcal B\big(t^m(h^\circ + u)\big)\,\eta'$ are holomorphic and bounded on the rows $u$ with $\lvert u_j\rvert \le \varrho$, uniformly in $t \ge t_0$, and by Cauchy's estimate so is the derivative of the first where $\lvert u_j\rvert \le \varrho/2$. With a row $\delta$ in the place of $\delta(t)$ the block gives $t^m\,(h^\circ + u)$, and there is a number $\varrho' > 0$ such that $\lvert u_j\rvert \le \varrho/2$ for the rows $\delta \in \mathbb C^{1\times q}$ with $\lvert\delta_j\rvert \le \varrho'\,t$ for all $j$. So the left side of \eqref{eq:gtseries}, which is affine in $\delta'$, is holomorphic in these rows $\delta$, and there its value at $\delta' = 0$ and its coefficients of the entries of $\delta'$ are at most a fixed power of $t$. By Cauchy's estimates, $\mathcal L_\alpha(t)$ and $\mathcal L'_{\alpha,i}(t)$ are at most a fixed power of $t$ times $(\varrho'\,t)^{-\lvert\alpha\rvert}$. With \eqref{eq:gtsizes}, a term with the exponent $r'$ is then at most a fixed power of $t$ times $(\kappa/\varrho')^{\lvert\alpha\rvert}\,e^{-r't^p/p}$. As $r' \ge \lvert\alpha\rvert\,\min_jr_j$, the terms with large $\lvert\alpha\rvert$ are bounded by a geometric series.
\end{itemize}

\emph{A generic ray.} We now take $\xi \in \mathcal C_k$ such that
\begin{itemize}[leftmargin=2.2em]
\item[(i)] the numbers $r_1, \dots, r_q$ are linearly independent over $\mathbb Q$;
\item[(ii)] $\iota_j(\xi) \neq 0$ for every $j$;
\item[(iii)] $g^\circ(\xi) \in \mathcal U^\infty$.
\end{itemize}
These $\xi$ are dense in $\mathcal C_k$. For (i): if integers $\kappa_1, \dots, \kappa_q$ are not all zero, $\sum_j\kappa_jr_j$ is $\xi\big(\sum_j\kappa_jc_jc_j^\top\big)\xi^\top$ for GELU and GeGLU blocks and $\xi\sum_j\kappa_js_jc_j$ for SwiGLU blocks, and by (K3) it vanishes only for $\xi$ on a quadric, or in a hyperplane. For (ii): in a gated block $\iota_j(\xi) = \xi e_j$, and $e_j \neq 0$ by (K3). For (iii): $g^\circ$ is a local diffeomorphism on an open dense set of rows, and $\mathcal U^\infty$ is open and dense. By (i), different multi-indices have different exponents $r_\alpha$, and only finitely many exponents lie below any bound.

\emph{The induction.} We show that, for every multi-index $\alpha$,
\begin{equation}\label{eq:gtclaim}
\mathcal L_\alpha(t) = 0 \quad\text{and}\quad \mathcal L'_{\alpha - 1_i,\,i}(t) = 0 \ \text{ for every } i \text{ with } \alpha_i \ge 1, \qquad\text{for all } t \in \R,
\end{equation}
by induction on the exponent $r_\alpha$. These are the coefficients of the terms of \eqref{eq:gtseries} with the exponent $r_\alpha$. Let \eqref{eq:gtclaim} hold for all multi-indices with a smaller exponent, and let $r^+ > r_\alpha$ be the next exponent. In \eqref{eq:gtseries} at $\delta(t)$, $\delta'(t)$, the terms with a smaller exponent vanish, and by (c) the terms with a larger exponent add up to at most a fixed power of $t$ times $e^{-r^+t^p/p}$. By (b), so do the terms with the exponent $r_\alpha$, which are
\[
(-1)^{\lvert\alpha\rvert}\,\delta(t)^\alpha\,\Big[\mathcal L_\alpha(t) + \sum_{i:\,\alpha_i\ge1}s_i\,\mathcal L'_{\alpha - 1_i,\,i}(t)\;\frac{\delta'_i(t)}{\delta_i(t)}\Big] .
\]
By \eqref{eq:gtsizes}, $\delta(t)^\alpha \ge 6^{-\lvert\alpha\rvert}\,e^{-r_\alpha t^p/p}$. So
\begin{equation}\label{eq:gtlevel}
\mathcal L_\alpha(t) \;+\; \sum_{i:\,\alpha_i\ge1}s_i\,\mathcal L'_{\alpha - 1_i,\,i}(t)\;\frac{\delta'_i(t)}{\delta_i(t)} \;\longrightarrow\; 0 \qquad\text{faster than every power of } 1/t .
\end{equation}

\emph{No small factor, $\alpha = 0$.} \eqref{eq:gtlevel} is $\mathcal L_0(t) \to 0$, and by (a) $\mathcal L_0 = 0$.

\emph{One small factor, $\alpha = 1_j$.} With $y = t\,\lvert\xi c_j\rvert$, \eqref{eq:gtlevel} is $\mathcal L_j + s_j\,\mathcal L'_j\,\delta'(y)/\delta(y) \to 0$.

For GELU and GeGLU blocks, $\delta'(y)/\delta(y) = \big(\operatorname{mr}(y) - y\big)/\big(y\operatorname{mr}(y)\big)$ and $y\operatorname{mr}(y) \le 1$, so
\[
\big(y\,\mathcal L_j + s_j\,\mathcal L'_j\big)\,\operatorname{mr}(y) \;-\; s_j\,y\,\mathcal L'_j \;\longrightarrow\; 0 \qquad\text{faster than every power of } 1/t .
\]
By (a), the maps $y\,\mathcal L_j + s_j\mathcal L'_j$ and $y\,\mathcal L'_j$ are meromorphic at infinity. By \cref{lem:e0nomero} with $z(t) = \lvert\xi c_j\rvert\,t$ both vanish, so $\mathcal L'_j = 0$ and $\mathcal L_j = 0$.

For SwiGLU blocks, $y\,\delta'(y)/\delta(y) = 1 - y + \delta(y)$, and $\delta(y)$ tends to $0$ faster than every power of $1/t$, so
\[
f(t) \;=\; \lvert\xi c_j\rvert\,t\;\mathcal L_j(t) + s_j\,\big(1 - \lvert\xi c_j\rvert\,t\big)\,\mathcal L'_j(t)
\]
tends to $0$ faster than every power of $1/t$ as $t \to +\infty$. By (a), $f$ is meromorphic at infinity, so $f(t) = 0$ for large $t$, and then for all $t \in \R$, because $f$ is real-analytic on $\R$. By \eqref{eq:gtparity}, $\mathcal L_j$ is odd in $t$ and $\mathcal L'_j$ is even. So $0 = f(t) - f(-t) = -2\,s_j\,\lvert\xi c_j\rvert\,t\;\mathcal L'_j(t)$: $\mathcal L'_j = 0$, and then $\mathcal L_j = 0$.

For all three kinds of block, by \eqref{eq:gtLprime}, $\mathcal L'_j(t) = \dot z^g_j(t\,\xi)\;\iota_j(t\,\xi)\;\Der\mathcal B^\infty(t^m\,h^\circ)\,a_j$. By (iii), (K3) and \cref{lem:e0tailrank}, $\Der\mathcal B^\infty(h^\circ)\,a_j \neq 0$, so the last factor, which is real-analytic in $t$, vanishes only at isolated $t$. By (ii), $\iota_j(t\,\xi) = t^{m-1}\,\iota_j(\xi) \neq 0$ for $t > 0$. So
\begin{equation}\label{eq:gtzg}
\dot z^g_j(t\,\xi) = 0 \qquad\text{for all } t > 0 .
\end{equation}

\emph{Two or more small factors, $\lvert\alpha\rvert \ge 2$.} Let $\alpha_i \ge 1$. Then $r_i < r_\alpha$, so \eqref{eq:gtzg} holds for the hidden unit $i$, and by \eqref{eq:gtLprime} the derivative of $\mathcal L^\infty(\tilde w, \tilde w'; t\,\xi)$ in $\tilde w'_i$ vanishes for all rows $\tilde w$. So $\mathcal L'_{\alpha - 1_i,\,i} = 0$. Then \eqref{eq:gtlevel} is $\mathcal L_\alpha(t) \to 0$, and by (a) $\mathcal L_\alpha = 0$. This ends the induction.

\emph{What the induction gives.} For this $\xi$ and every $t > 0$, all coefficients of the expansion of $(\tilde w, \tilde w') \mapsto \mathcal L^\infty(\tilde w, \tilde w'; t\,\xi)$ at $(t\,\xi W_gI_{\mathcal A}, 1_{\mathcal A})$ vanish, so this real-analytic map vanishes. The rows $t\,\xi$ of this kind are dense in $\mathcal C_k$. So, with \eqref{eq:gtzg}, for all $n \in \mathcal C_k$ and all rows $\tilde w \in \R^{1\times q}$, with $\hat w = \tilde w\odot\iota(n)$,
\begin{equation}\label{eq:gtall}
\dot z^g(n) = 0, \qquad \Der\mathcal B^\infty(\hat wV)\,\Big(\big(\tilde w\odot\dot z^v(n)\big)\,V + \hat w\,\eta_V\Big) + \Der_\theta\mathcal B^\infty(\hat wV)\,\eta' = 0 .
\end{equation}

\emph{Step 1: the gates.} By \eqref{eq:gtzdot}, $\dot z^g(n)$ is a polynomial in $n$ that vanishes on the open set $\mathcal C_k$, so it vanishes for all $n$. In the case (n) this is $n\,(ZW_g + \eta_g) = 0$, so $\eta_g = -ZW_g$, and $\mathfrak g$ contains every $Z$. In the case (+) its part of degree $3$ is $\langle nZ, n\rangle\,nW_g = 0$, so $Z$ is antisymmetric, as $W_g \neq 0$, and its part of degree $1$ gives $\eta_g = -ZW_g$. In both cases $Z \in \mathfrak g$, and in a gated block $\dot z^v(n) = n\,(ZW_v + \eta_v)$.

\emph{Step 2: the values and the write.} In a gated block, fix $n \in \mathcal C_k$ with $ne_j \neq 0$ for every $j$, and let $\Delta(n)$ be the diagonal matrix with the entries $-\dot z^v_j(n)/(ne_j)$, so that $\tilde w\odot\dot z^v(n) = -\hat w\,\Delta(n)$. In a GELU block, fix $n \in \mathcal C_k$ and let $\Delta(n) = 0$. The row $\hat w$ takes every value in $\R^{1\times q}$, and \eqref{eq:gtall} reads
\[
\Der\mathcal B^\infty(\hat wV)\,\Big(\hat w\,\big(\eta_V - \Delta(n)\,V\big)\Big) + \Der_\theta\mathcal B^\infty(\hat wV)\,\eta' = 0 \qquad\text{for all rows } \hat w .
\]
By \cref{lem:e0atwrite} for $\mathcal B^\infty$, with $A = V$, there is a matrix $Z'(n)$ with
\begin{equation}\label{eq:gtZprime}
\eta_V - \Delta(n)\,V = V\,Z'(n), \qquad \Der\mathcal B^\infty(h)\,\big(hZ'(n)\big) + \Der_\theta\mathcal B^\infty(h)\,\eta' = 0 \quad\text{for all rows } h .
\end{equation}
For two such rows $n$ and $n_1$, subtracting gives $\Der\mathcal B^\infty(h)\,\big(h\,(Z'(n) - Z'(n_1))\big) = 0$ for all $h$, so $Z'(n) = Z'(n_1)$ by \cref{lem:e0tailrank}. Let $Z'$ be this matrix. Then $\Delta(n)\,V = \eta_V - VZ'$ does not depend on $n$, and no row of $V$ vanishes, so $\Delta(n)$ is one diagonal matrix $\Delta_{k+1}$ for all these $n$, with $\Delta_{k+1} = 0$ in a GELU block. So $\eta_V = \Delta_{k+1}V + VZ'$. In a gated block, $\dot z^v(n) = -n\,W_v\Delta_{k+1}$ on an open set of rows $n$, that is, $n\,(ZW_v + \eta_v + W_v\Delta_{k+1}) = 0$: $\eta_v = -ZW_v - W_v\Delta_{k+1}$.

\emph{Step 3: the rest.} In the case (+), $\mathcal B^\infty = \mathcal B_{>k+1}$, and the second part of \eqref{eq:gtZprime} is \eqref{eq:e0D} for $k+1$, with $Z'$ and $\eta'$. By $(D_{k+1})$, $Z' \in \mathfrak g$ and $\eta'$ is of the form \eqref{eq:gtDconcl}, or \eqref{eq:sbDconcl} for GELU blocks, with $Z_{k+1} = Z'$. With Steps~1 and~2 this is \eqref{eq:gtDconcl}, or \eqref{eq:sbDconcl}, for $k$.

In the case (n), where the blocks are gated, \eqref{eq:gtZprime} is an identity for the frozen network. The frozen network does not change along some directions along which the network does, for example along every direction of the weights of a hidden unit that is not active, so \eqref{eq:gtZprime} does not determine $\eta'$. Let $\eta^\flat$ be the direction \eqref{eq:gtDconcl} for $k$ with $Z_k = Z$, with $\Delta_{k+1}$ as in Step~2, with $Z_{k+1} = Z'$, and with $Z_\ell = 0$ and $\Delta_\ell = 0$ for all later $\ell$; here $\mathfrak g$ contains every matrix. By the converse of the claim, $Z$ and $\eta^\flat$ satisfy \eqref{eq:e0D} for $k$, and so do $Z$ and $\eta$. So the matrix $0$ and the direction $\eta - \eta^\flat$ satisfy it. By Steps~1 and~2, $\eta - \eta^\flat$ vanishes at the weights of block $k+1$. Let $\eta''$ be its entries at the later weights. Then \eqref{eq:e0D} for $k$ reads
\[
\Der_\theta\mathcal B_{>k+1}\big(g_{k+1}(h)\big)\,\eta'' = 0 \qquad\text{for all rows } h .
\]
By (R3) the values of $g_{k+1}$ contain an open set, and the left side is real-analytic, so $\Der_\theta\mathcal B_{>k+1}(h)\,\eta'' = 0$ for all rows $h$. This is \eqref{eq:e0D} for $k+1$ with the matrix $0$. By $(D_{k+1})$, $\eta''$ is of the form \eqref{eq:gtDconcl} for $k+1$ with $Z_{k+1} = 0$. So $\eta = \eta^\flat + \eta''$ is of the form \eqref{eq:gtDconcl} for $k$.
\end{proof}

\para{The parameters $\theta^\dagger$.} Let $\mathbf 1$ be the column of $\dm$ ones. Choose numbers $\mu_1, \mu_2, \dots > 0$ such that $1, \mu_1, \mu_2, \dots$ are linearly independent over $\mathbb Q$, for SwiGLU blocks, or such that $\mu_1^2, \mu_2^2, \dots$ are, for GELU and GeGLU blocks. Almost every choice of finitely many numbers is of this kind. Let $\theta^\dagger$ be the parameters with
\[
E = \begin{bmatrix} I \\ 0 \end{bmatrix} \in \R^{d_{\mathrm{in}}\times\dm}, \qquad D = \begin{bmatrix} I & 0 \end{bmatrix} \in \R^{\dm\times d_{\mathrm{out}}},
\]
and, in every block $\ell$, with $q = q_\ell$ hidden units,
\[
W_g^{\ell} = \big[\,I \;\big|\; {-\mu_1\mathbf 1} \;\big|\; \cdots \;\big|\; {-\mu_{q-\dm}\mathbf 1}\,\big], \qquad \Wwrite^{\ell} = \begin{bmatrix} I \\ \mathbf 1^\top \\[-3pt] \vdots \\ \mathbf 1^\top \end{bmatrix}, \qquad W_v^{\ell} = \big[\,\mathbf 1 \;\big|\; \cdots \;\big|\; \mathbf 1\,\big] ,
\]
where only a gated block has $W_v^{\ell}$. So the gate column $c_j$ is column $j$ of $I$ for $j \le \dm$ and $c_j = -\mu_{j-\dm}\,\mathbf 1$ for $j > \dm$, the write row $a_j$ is row $j$ of $I$ for $j \le \dm$ and $a_j = \mathbf 1^\top$ for $j > \dm$, and in a gated block every value column is $e_j = \mathbf 1$. Let $\R^{1\times\dm}_{>0}$ be the positive orthant, the set of the rows $n$ with $n_i > 0$ for every $i$. It is an open convex cone that does not contain $0$.

\begin{lemma}[$\theta^\dagger$ is good]\label{lem:gtstar}
$\theta^\dagger$ is good for the cones $\mathcal C_0 = \dots = \mathcal C_L = \R^{1\times\dm}_{>0}$. For every block, $\mathcal A_\ell = \{1, \dots, \dm\}$, and $g^\circ_\ell(n) = n$ for a GELU block and $g^\circ_\ell(n) = (n\mathbf 1)\,n$ for a gated block, where $n\mathbf 1 = n_1 + \dots + n_{\dm}$.
\end{lemma}

\begin{proof}
(R1): $E$, $D$ and every $\Wwrite^{\ell}$ contain the identity matrix. (K1): for $n \in \R^{1\times\dm}_{>0}$, $nc_j = n_j > 0$ for $j \le \dm$ and $nc_j = -\mu_{j-\dm}\,n\mathbf 1 < 0$ for $j > \dm$. So $\mathcal A_\ell = \{1, \dots, \dm\}$, and $g^\circ_\ell(n) = \sum_{j\le\dm}n_j\,\iota_j(n)\,a_j$ is $n$ for a GELU block and $(n\mathbf 1)\,n$ for a gated block, where $\iota_j(n) = n\mathbf 1$. (K2): $g^\circ_\ell(n)$ is a positive multiple of $n$. Its derivative is the identity for a GELU block, and for a gated block the map $v \mapsto (n\mathbf 1)\,v + (v\mathbf 1)\,n$, with the determinant $2\,(n\mathbf 1)^{\dm} \neq 0$. (K3): no $a_j$ and no $e_j$ vanishes. For SwiGLU blocks, $\sum_j\kappa_jc_j = 0$ with integers $\kappa_j$ says that $\kappa_j = \sum_{i>\dm}\kappa_i\,\mu_{i-\dm}$ for every $j \le \dm$, so every $\kappa_j$ vanishes, by the choice of the $\mu_i$. For GELU and GeGLU blocks, the entries of $\sum_j\kappa_jc_jc_j^\top = 0$ off the diagonal, which exist because $\dm \ge 2$, give $\sum_{i>\dm}\kappa_i\,\mu_{i-\dm}^2 = 0$, so $\kappa_i = 0$ for $i > \dm$, and then the diagonal gives $\kappa_j = 0$ for $j \le \dm$.
\end{proof}

\subsection{The great circle: \texorpdfstring{$\varepsilon = 0$}{no regularization}}\label{app:circle}

\para{The great circle: $\varepsilon = 0$.} We now treat the case (0), for $\tanh$ and GELU blocks. The identity \eqref{eq:e0star} then holds only on the sphere $\mathbb S$. We continue it along a great circle of $\mathbb S$ to complex angles, on which the inputs of the hidden units go to infinity.

\para{RMSNorm at complex rows.} For $h \in \mathbb C^{1\times\dm}$ with $\operatorname{Re}\langle h, h\rangle > 0$ let
\[
N_\varepsilon(h) \;=\; \frac{\sqrt{\dm}\;h}{\sqrt{\langle h, h\rangle}},
\]
with the principal square root. It is holomorphic, $N_\varepsilon(\lambda h) = N_\varepsilon(h)$ for $\lambda > 0$, and
\begin{equation}\label{eq:e0DN}
\langle n, n\rangle = \dm \quad\Longrightarrow\quad N_\varepsilon(n) = n, \qquad \Der N_\varepsilon(n)\,v = v - \frac{\langle v, n\rangle}{\dm}\,n .
\end{equation}
For $\varkappa \in (0, 1]$ let
\[
\mathbb C_\varkappa \;=\; \big\{u + \mathrm iu' \ :\ u, u' \in \R^{1\times\dm},\ \norm{u'} < \varkappa\norm{u}\big\},
\]
an open cone in $\mathbb C^{1\times\dm}$ that contains every real row $h \neq 0$. On $\mathbb C_1$, $\operatorname{Re}\langle h, h\rangle = \norm{u}^2 - \norm{u'}^2 > 0$.

\begin{lemma}[$\mathcal B_{>k}$ on a complex cone]\label{lem:e0cone}
Let $\theta$ satisfy (R1) and (R2). For every $k$ there is a $\varkappa_k \in (0, 1]$ such that $\mathcal B_{>k}$, with its derivatives $\Der\mathcal B_{>k}$ and $\Der_\theta\mathcal B_{>k}$, extends to a holomorphic map on $\mathbb C_{\varkappa_k}$ with $\mathcal B_{>k}(\lambda h) = \mathcal B_{>k}(h)$ for $\lambda > 0$. The number $\varkappa_k$ depends only on the weights of the blocks $k+1, \dots, L$ and on $D$.
\end{lemma}

\begin{proof}
$\mathcal B_{>L}(h) = N_\varepsilon(h)\,D$ is holomorphic on $\mathbb C_1$: $\varkappa_L = 1$. Let $k < L$. The map $g_{k+1}$ is holomorphic near $\mathbb S$ and maps the compact set $\mathbb S$ into the real rows $h \neq 0$, by (R2), hence into the open set $\mathbb C_{\varkappa_{k+1}}$. So there is a $\rho > 0$ such that $n \mapsto \mathcal B_{>k+1}(g_{k+1}(n))$ is holomorphic on the set of the $n \in \mathbb C^{1\times\dm}$ at distance less than $\rho$ from $\mathbb S$. On the compact set of the $u + \mathrm iu'$ with $\norm{u} = 1$ and $\norm{u'} \le \frac12$, $N_\varepsilon$ is uniformly continuous, and $N_\varepsilon(u) \in \mathbb S$. So there is a $\varkappa_k$ such that $N_\varepsilon$ maps $\mathbb C_{\varkappa_k}$ into the points at distance less than $\rho$ from $\mathbb S$, and $\mathcal B_{>k} = \big(\mathcal B_{>k+1}\circ g_{k+1}\big)\circ N_\varepsilon$ is holomorphic on $\mathbb C_{\varkappa_k}$. The weights enter holomorphically as well.
\end{proof}

\para{The path.} For orthonormal rows $\xi, \xi' \in \R^{1\times\dm}$ let $\nu = \xi + \mathrm i\xi'$, $\bar\nu = \xi - \mathrm i\xi'$ and
\[
n(t) \;=\; \frac{\sqrt{\dm}}{2}\,\Big(t\,\nu + \frac{1}{t}\,\bar\nu\Big), \qquad t \in \mathbb C\setminus\{0\}.
\]
As $\langle\nu, \nu\rangle = \langle\bar\nu, \bar\nu\rangle = 0$ and $\langle\nu, \bar\nu\rangle = 2$, we have $\langle n(t), n(t)\rangle = \dm$ for all $t$. For $\lvert t\rvert = 1$ the row $n(t)$ is real and runs through the great circle of $\xi$ and $\xi'$ in $\mathbb S$, with $n(1) = \sqrt{\dm}\;\xi$. For real $t \ge 1$,
\begin{equation}\label{eq:e0path}
n(t) \;=\; \sqrt{\dm}\,\Big(\frac{t + t^{-1}}{2}\;\xi + \mathrm i\,\frac{t - t^{-1}}{2}\;\xi'\Big).
\end{equation}
In the notation of the proof of \cref{lem:sphere}, this is the great circle of $\xi$ and $\xi'$ at the angle $\vartheta = \mathrm i\log t$.

\para{The identity on the path.} With the notation of \eqref{eq:e0zdot} and \eqref{eq:e0L}, let $\varkappa = \varkappa_{k+1}$. The same formulas define $\dot z(n) \in \mathbb C^{1\times q}$ for $n \in \mathbb C^{1\times\dm}$, and $\mathcal L(\tilde w, \tilde w'; n)$ for rows $\tilde w, \tilde w' \in \mathbb C^{1\times q}$ with $\tilde wV \in \mathbb C_\varkappa$. $\mathcal L$ is holomorphic in $\tilde w$, affine in $\tilde w'$ and polynomial in $n$.

\begin{lemma}[Continuation]\label{lem:e0cont}
In the case (0), let \eqref{eq:e0star} hold, and let $\xi, \xi'$ be orthonormal. Suppose that for every real $t \ge 1$ the function $\sigma$ is holomorphic at every $n(t)\,c_j$ and $\sigma(n(t)W)\,V \in \mathbb C_\varkappa$. Then
\[
\mathcal L\big(\sigma(n(t)W),\ \sigma'(n(t)W);\ n(t)\big) \;=\; 0 \qquad\text{for all real } t \ge 1 .
\]
\end{lemma}

\begin{proof}
The left side is a holomorphic function of $t$ on an open set that contains $[1, \infty)$, hence on a connected open neighborhood of $[1, \infty)$. This neighborhood contains an arc of the unit circle around $t = 1$, on which $n(t) \in \mathbb S$ and the function vanishes by \eqref{eq:e0star}. So it vanishes on the whole neighborhood.
\end{proof}

\para{The step for $\tanh$ on the path.} Let $\sigma = \tanh$. Consider, for a block $\ell$ with $W = \Wread^{\ell}$, its columns $c_j$, $V = \Wwrite^{\ell}$ and $\varkappa_\ell$ of \cref{lem:e0cone}:
\begin{itemize}[leftmargin=2.9em]
\item[(S1)] any two rows of $V$ are linearly independent;
\item[(S2)] the columns of $W$ are linearly independent over $\mathbb Q$;
\item[(S3)] there are a unit row $\xi_0$ and a number $\rho \in (0, 1]$ such that every $\xi_0c_j \neq 0$ and, with $s_j$ the sign of $\xi_0c_j$ and $\bar s = (s_1, \dots, s_q)$: the rows $\tilde wV$ with $\lvert\tilde w_j - s_j\rvert \le \rho$ for all $j$ lie in $\mathbb C_{\varkappa_\ell}$, and $e^{-2\sqrt{\dm}\,\lvert\xi_0c_j\rvert} < \rho/3$ for every $j$.
\end{itemize}
(S3) asks for large reads: on the row $\sqrt{\dm}\,\xi_0$ every hidden unit of the block is close to saturation. It then stays so on the whole path \eqref{eq:e0path}, and the outputs of the block stay in a small set, on which the rest of the network is holomorphic.

\begin{lemma}[Separation on the path]\label{lem:e0sep}
Let $u_1, \dots, u_q > 0$ be linearly independent over $\mathbb Q$, let $u'_1, \dots, u'_q$ be real, and for $t \ge 1$ let
\[
\zeta_j(t) \;=\; \exp\Big(-\big(t + t^{-1}\big)\,u_j - \mathrm i\,\big(t - t^{-1}\big)\,u'_j\Big).
\]
Let $\rho_0 > \max_j e^{-2u_j}$, and let $\mathcal T_{-K}, \dots, \mathcal T_K$ be holomorphic maps from the polydisc $\{\zeta \in \mathbb C^q : \lvert\zeta_j\rvert < \rho_0\}$ to $\mathbb C^{1\times d_{\mathrm{out}}}$. If
\[
\sum_{i=-K}^{K} t^i\ \mathcal T_i\big(\zeta_1(t), \dots, \zeta_q(t)\big) \;=\; 0 \qquad\text{for all } t \ge 1,
\]
then $\mathcal T_i = 0$ for every $i$.
\end{lemma}

\begin{proof}
This is the proof of \cref{lem:raysep}, with a phase and with negative powers. On the polydisc $\mathcal T_i(\zeta) = \sum_\alpha b_{i,\alpha}\,\zeta^\alpha$, over the multi-indices $\alpha$ of nonnegative integers, and the series converges absolutely at $\lvert\zeta_j\rvert = e^{-2u_j}$. Let $\mu_\alpha = \sum_j\alpha_ju_j$ and $p_\alpha(t) = \sum_i t^i\,b_{i,\alpha}$. Then $\lvert\zeta(t)^\alpha\rvert = e^{-(t + t^{-1})\mu_\alpha}$, and the hypothesis reads $\sum_\alpha p_\alpha(t)\,\zeta(t)^\alpha = 0$. The numbers $\mu_\alpha$ are pairwise different, and finitely many lie below any bound. Suppose that some $p_\alpha$ is not zero, let $\alpha^*$ be the one among these with the smallest $\mu_\alpha$, say $\mu^*$, and let $\gamma > 0$ be the distance from $\mu^*$ to the next larger $\mu_\alpha$. For $t \ge 1$, as $t + t^{-1} \ge 2$,
\[
\lvert p_{\alpha^*}(t)\rvert \;\le\; \sum_{\mu_\alpha > \mu^*}\lvert p_\alpha(t)\rvert\,e^{-(t + t^{-1})(\mu_\alpha - \mu^*)} \;\le\; t^K\,e^{-(t + t^{-1} - 2)\gamma}\,e^{2\mu^*}\sum_{i}\sum_\alpha\lvert b_{i,\alpha}\rvert\,e^{-2\mu_\alpha},
\]
which tends to $0$ faster than every power of $1/t$. But $p_{\alpha^*}(t) = \sum_{i=-K}^{K}t^i\,b_{i,\alpha^*}$ is not zero, so $\lvert p_{\alpha^*}(t)\rvert$ is at least a positive constant times $t^{-K}$ for large $t$, a contradiction.
\end{proof}

\begin{lemma}[A change of the reads moves the output]\label{lem:e0wedge}
In the case (0), let $\theta$ be regular with (S1) for block $k+1$, let $V = \Wwrite^{k+1}$ have the rows $a_1, \dots, a_q$, and let $b \in \mathbb C^{1\times q}$. If
\[
\Der\mathcal B_{>k+1}(\tilde wV)\,\Big(\big((1 - \tilde w^2)\odot b\big)\,V\Big) \;=\; 0
\]
for all $\tilde w$ in a nonempty open subset of $\R^{1\times q}$, where $1 - \tilde w^2$ is the row with the entries $1 - \tilde w_j^2$, then $b = 0$.
\end{lemma}

\begin{proof}
At real $\tilde w$ the derivative is a real linear map, so the real and the imaginary part of $b$ satisfy the hypothesis: let $b$ be real. By \cref{lem:e0tailrank}, for $\tilde w$ in an open dense subset of the open set, the row $\sum_j b_j(1 - \tilde w_j^2)\,a_j$ is a multiple of the row $\sum_l\tilde w_l\,a_l$. So all $2\times2$ minors of these two rows vanish on an open set, hence as polynomials in $\tilde w$. Their part of degree $1$ in $\tilde w$ shows that $bV$ is a multiple of every $a_l$, so $bV = 0$ by (S1). The coefficient of $\tilde w_j^2\,\tilde w_l$ with $l \neq j$ shows that $b_ja_j$ is a multiple of $a_l$, so $b_j = 0$ by (S1).
\end{proof}

\begin{lemma}[The step for $\tanh$ on the path]\label{lem:e0steptanh}
Let $\sigma = \tanh$, in the case (0), let $\theta$ be regular, and let block $k+1$ satisfy (S1) to (S3). Then $(D_{k+1})$ implies $(D_k)$.
\end{lemma}

\begin{proof}
Let $Z$ and $\eta$ satisfy \eqref{eq:e0D} for $k$, so that \eqref{eq:e0star} holds, with the notation above. By (S3) there is an open set $\mathcal P$ of unit rows around $\xi_0$ such that $\xi c_j$ has the sign $s_j$ and $e^{-2\sqrt{\dm}\,\lvert\xi c_j\rvert} < \rho/3$ for every $\xi \in \mathcal P$ and every $j$. By (S2), for integers $\kappa_j$ that are not all zero, $\sum_j\kappa_j\,\xi c_j = \xi\sum_j\kappa_jc_j$ vanishes only for $\xi$ in a hyperplane. Let $\mathcal P' \subset \mathcal P$ be the dense set of the $\xi$ for which $\xi c_1, \dots, \xi c_q$ are linearly independent over $\mathbb Q$.

\emph{The path.} Fix $\xi \in \mathcal P'$ and a unit row $\xi'$ orthogonal to $\xi$, and let $n(t)$ be the path \eqref{eq:e0path}. Let
\[
\zeta_j(t) \;=\; e^{-2s_j\,n(t)\,c_j} \;=\; \exp\Big(-\big(t + t^{-1}\big)\sqrt{\dm}\,\lvert\xi c_j\rvert - \mathrm i\,\big(t - t^{-1}\big)\sqrt{\dm}\,s_j\,\xi'c_j\Big),
\]
so that $\lvert\zeta_j(t)\rvert < \rho/3$ for $t \ge 1$, and
\[
\tanh\big(n(t)\,c_j\big) = s_j\,\frac{1 - \zeta_j(t)}{1 + \zeta_j(t)}, \qquad \tanh'\big(n(t)\,c_j\big) = 1 - \tanh^2\big(n(t)\,c_j\big).
\]
For $\zeta$ in the polydisc $\mathbb D = \{\lvert\zeta_j\rvert < \rho/3\}$ let $\tilde w(\zeta)$ be the row with the entries $s_j\frac{1 - \zeta_j}{1 + \zeta_j}$. As $\lvert\tilde w_j(\zeta) - s_j\rvert = \frac{2\lvert\zeta_j\rvert}{\lvert1 + \zeta_j\rvert} \le 3\lvert\zeta_j\rvert < \rho$, (S3) gives $\tilde w(\zeta)\,V \in \mathbb C_\varkappa$. The function $\tanh$ is holomorphic at $n(t)\,c_j$, as $\zeta_j(t) \neq -1$. So \cref{lem:e0cont} applies:
\begin{equation}\label{eq:e0tanhpath}
\mathcal L\big(\tilde w(\zeta(t)),\ 1 - \tilde w(\zeta(t))^2;\ n(t)\big) \;=\; 0 \qquad\text{for all } t \ge 1 .
\end{equation}

\emph{The powers of $t$.} By \eqref{eq:e0zdot}, $\dot z(n(t)) = t^3\dot z^{(3)} + t\,\dot z^{(1)} + t^{-1}\dot z^{(-1)} + t^{-3}\dot z^{(-3)}$ with rows $\dot z^{(i)} \in \mathbb C^{1\times q}$ that do not depend on $t$, and
\begin{equation}\label{eq:e0powers}
\begin{aligned}
\dot z^{(3)} &= -\frac{\sqrt{\dm}}{8}\,\langle\nu Z, \nu\rangle\;\nu W, \\
\dot z^{(1)} &= \frac{\sqrt{\dm}}{2}\,\Big(\nu\,(ZW + \eta_W) - \frac14\Big(\big(\langle\nu Z, \bar\nu\rangle + \langle\bar\nu Z, \nu\rangle\big)\,\nu + \langle\nu Z, \nu\rangle\,\bar\nu\Big)W\Big).
\end{aligned}
\end{equation}
$\mathcal L$ is affine in $\dot z$, so \eqref{eq:e0tanhpath} reads $\sum_i t^i\,\mathcal T_i(\zeta(t)) = 0$ with the maps
\begin{align*}
\mathcal T_i(\zeta) &= \Der\mathcal B\big(\tilde w(\zeta)V\big)\,\Big(\big((1 - \tilde w(\zeta)^2)\odot\dot z^{(i)}\big)\,V\Big) \qquad (i = \pm1, \pm3),\\
\mathcal T_0(\zeta) &= \Der\mathcal B\big(\tilde w(\zeta)V\big)\,\big(\tilde w(\zeta)\,\eta_V\big) + \Der_\theta\mathcal B\big(\tilde w(\zeta)V\big)\,\eta' ,
\end{align*}
which are holomorphic on $\mathbb D$. By \cref{lem:e0sep}, with $u_j = \sqrt{\dm}\,\lvert\xi c_j\rvert$ and $\rho_0 = \rho/3$, every $\mathcal T_i$ vanishes on $\mathbb D$. For real $\zeta \in \mathbb D$ the rows $\tilde w(\zeta)$ fill a neighborhood of $\bar s$ in $\R^{1\times q}$.

\emph{(a) $\mathcal T_3 = 0$ gives $Z$.} Here $\dot z^{(3)}$ is the number $-\frac{\sqrt{\dm}}{8}\langle\nu Z, \nu\rangle$ times the row $\nu W$, which is not zero. By \cref{lem:e0wedge}, $\langle\nu Z, \nu\rangle = 0$. Its imaginary part is $\xi\,(Z + Z^\top)\,\xi'^\top = 0$. This holds for every unit row $\xi'$ orthogonal to $\xi$, so $\xi\,(Z + Z^\top)$ is a multiple of $\xi$, for every $\xi \in \mathcal P'$ and then, by continuity, for every $\xi$ in the open set $\mathcal P$. A matrix for which every row of an open set is an eigenvector is a multiple of the identity. So $Z + Z^\top = 2\beta_kI$, and $Z = Y_k + \beta_kI$ with $Y_k$ antisymmetric.

\emph{(b) $\mathcal T_1 = 0$ gives $\eta_W$.} Now $\langle\nu Z, \nu\rangle = 0$ and $\langle\nu Z, \bar\nu\rangle + \langle\bar\nu Z, \nu\rangle = 2\beta_k\langle\nu, \bar\nu\rangle = 4\beta_k$, so by \eqref{eq:e0powers} $\dot z^{(1)} = \frac{\sqrt{\dm}}{2}\,\nu\,(Y_kW + \eta_W)$. By \cref{lem:e0wedge}, $\dot z^{(1)} = 0$, so $\xi\,(Y_kW + \eta_W) = 0$ for every $\xi \in \mathcal P'$, and $\eta_W = -Y_kW$.

\emph{(c) $\mathcal T_0 = 0$ gives the rest.} $\Der\mathcal B(\tilde wV)\,(\tilde w\,\eta_V) + \Der_\theta\mathcal B(\tilde wV)\,\eta' = 0$ for real $\tilde w$ near $\bar s$. The left side is real-analytic on the rows $\tilde w \in \R^{1\times q}$ with $\tilde wV \neq 0$. These form a connected set, the complement of a subspace of codimension $\dm \ge 2$, so the identity holds on all of it. By \cref{lem:e0atwrite} with $A = V$, $\eta_V = VZ'$ and \eqref{eq:e0D} holds for $k+1$ with $Z'$ and $\eta'$. By $(D_{k+1})$, $Z' = Y_{k+1} + \beta_{k+1}I$ and $\eta'$ is as in \eqref{eq:e0Dconcl}. With (a) and (b) this is \eqref{eq:e0Dconcl} for $k$.
\end{proof}

\para{The step for GELU on the path.} Consider, for a block $\ell$ with $W = \Wread^{\ell}$, its columns $c_j$, $V = \Wwrite^{\ell}$, its rows $a_j$, and $\varkappa_\ell$ of \cref{lem:e0cone}:
\begin{itemize}[leftmargin=2.9em]
\item[(P1)] the columns of $W$ are pairwise different, and the matrix $WV$ is invertible;
\item[(P2)] there is a unit row $\xi_0$ such that, for every unit row $\xi'$ orthogonal to $\xi_0$,
\begin{gather*}
\lvert\xi'c_j\rvert < \xi_0c_j \ \text{ for every } j, \qquad\qquad
\norm{\xi'\,WV} < \frac{\varkappa_\ell}{2}\,\norm{\xi_0\,WV}, \\
\sum_{j}e^{-\dm(\xi_0c_j)^2/2}\,\norm{a_j} \;<\; \frac{\varkappa_\ell\sqrt{\dm}}{8}\,\norm{\xi_0\,WV} .
\end{gather*}
\end{itemize}
(P2) asks for reads that are large and close to a matrix of rank one: every column of $W$ is close to a positive multiple of $\xi_0^\top$. Then every hidden unit is close to linear on the whole path, the output of the block is close to $n(t)\,WV$, and its imaginary part stays small against its real part, so that the rest of the network is holomorphic along the path (\cref{lem:e0cone}).

\begin{lemma}[The step for GELU on the path]\label{lem:e0stepgelu}
Let $\sigma$ be GELU, in the case (0), let $\theta$ be regular, and let block $k+1$ satisfy (P1) and (P2). Then $(D_{k+1})$ implies $(D_k)$.
\end{lemma}

\begin{proof}
Let $Z$ and $\eta$ satisfy \eqref{eq:e0D} for $k$, so that \eqref{eq:e0star} holds, with the notation above. Both sides of every inequality of (P2), the left sides as suprema over $\xi'$, are continuous in $\xi_0$. So there is an open set $\mathcal P$ of unit rows around $\xi_0$ such that (P2) holds for every $\xi \in \mathcal P$ in the place of $\xi_0$.

\emph{The path.} Fix $\xi \in \mathcal P$ and a unit row $\xi'$ orthogonal to $\xi$, let $n(t)$ be the path \eqref{eq:e0path}, and
\[
z_j(t) = n(t)\,c_j, \qquad r_j = (\xi c_j)^2 - (\xi'c_j)^2 > 0, \qquad \delta_j(t) = z_j - \sigma(z_j), \qquad \delta'_j(t) = 1 - \sigma'(z_j).
\]
Let $z(t), \delta(t), \delta'(t)$ be the rows with these entries. For $t \ge 1$, by \eqref{eq:e0path} and (P2),
\begin{equation}\label{eq:e0zfacts}
\begin{gathered}
\operatorname{Re}z_j \ge \sqrt{\dm}\;\xi c_j > 0, \qquad \lvert\operatorname{Im}z_j\rvert \le \operatorname{Re}z_j, \\
\operatorname{Re}\big(z_j^2\big) = \dm\,(\xi c_j)^2 + \frac{\dm}{4}\,\big(t - t^{-1}\big)^2\,r_j .
\end{gathered}
\end{equation}
By \cref{lem:e0mills}, $\delta_j = \varphi(z_j)\,z_j\,\operatorname{mr}(z_j)$ and $\delta'_j = \varphi(z_j)\,\big(\operatorname{mr}(z_j) - z_j\big)$, with $\lvert\operatorname{mr}(z_j)\rvert \le 1/\operatorname{Re}z_j$ and $\lvert z_j\rvert \le \sqrt2\,\operatorname{Re}z_j$. So, with a constant $\kappa$,
\begin{equation}\label{eq:e0deltabound}
\lvert\delta_j(t)\rvert \;\le\; e^{-\dm(\xi c_j)^2/2}\;e^{-\frac{\dm}{8}(t - t^{-1})^2r_j}, \qquad \lvert\delta'_j(t)\rvert \;\le\; \kappa\,t\;e^{-\frac{\dm}{8}(t - t^{-1})^2r_j} .
\end{equation}
The real and the imaginary part of $n(t)\,WV$ are $\sqrt{\dm}\,\frac{t + t^{-1}}{2}\,\xi WV$ and $\sqrt{\dm}\,\frac{t - t^{-1}}{2}\,\xi'WV$. By (P2), the second is shorter than $\frac{\varkappa}{2}$ times the first, and by \eqref{eq:e0deltabound} $\norm{\delta(t)\,V} \le \sum_j\lvert\delta_j\rvert\,\norm{a_j}$ is at most $\frac{\varkappa}{8}$ times the norm of the first. Let $\varrho = \varkappa\sqrt{\dm}\,\norm{\xi WV}\big/\big(16\sum_j\norm{a_j}\big)$. For a row $\delta \in \mathbb C^{1\times q}$,
\begin{equation}\label{eq:e0incone}
\lvert\delta_j\rvert \le \varrho\,t \ \text{ for all } j, \ \text{ or } \delta = \delta(t) \qquad\Longrightarrow\qquad \big(z(t) - \delta\big)\,V = n(t)\,WV - \delta V \ \in \mathbb C_\varkappa,
\end{equation}
because then $\norm{\delta V}$ is at most $\frac{\varkappa}{8}$ times the norm of the real part of $n(t)\,WV$, and $\frac{\varkappa/2 + \varkappa/8}{1 - \varkappa/8} < \varkappa$. As $\sigma(z) = z - \delta(t)$ and $\sigma'(z) = 1 - \delta'(t)$ on the path, \cref{lem:e0cont} applies:
\begin{equation}\label{eq:e0gelupath}
\mathcal L\big(z(t) - \delta(t),\ 1 - \delta'(t);\ n(t)\big) \;=\; 0 \qquad\text{for all } t \ge 1 .
\end{equation}

\emph{The expansion.} For fixed $t$, by \eqref{eq:e0incone} the map $(\delta, \delta') \mapsto \mathcal L\big(z(t) - \delta, 1 - \delta'; n(t)\big)$ is holomorphic in the rows $\delta$ with $\lvert\delta_j\rvert \le \varrho\,t$ and affine in the row $\delta'$. We expand it in powers of $\delta$ and $\delta'$. Its value at $0$ and its first derivatives are
\[
\mathcal L_0(t) = \mathcal L\big(z(t), 1; n(t)\big), \qquad \mathcal L_j(t) = \frac{\partial\mathcal L}{\partial\tilde w_j}\big(z(t), 1; n(t)\big), \qquad \mathcal L'_j(t) = \frac{\partial\mathcal L}{\partial\tilde w'_j}\big(z(t), 1; n(t)\big),
\]
so that \eqref{eq:e0gelupath} is
\begin{equation}\label{eq:e0expansion}
\mathcal L_0(t) \;-\; \sum_j\Big[\mathcal L_j(t)\,\delta_j(t) + \mathcal L'_j(t)\,\delta'_j(t)\Big] \;+\; \big(\text{terms with two or more factors}\big) \;=\; 0 ,
\end{equation}
where the factors are entries of $\delta(t)$ and $\delta'(t)$. We use two facts about the coefficients. First, every coefficient is meromorphic at infinity. Indeed, it is built from derivatives of $\mathcal B$ at $z(t)\,V = n(t)\,WV$ and from polynomials in $n(t)$, and
\[
n(t)\,WV \;=\; \sqrt{\dm}\;\frac{t + t^{-1}}{2}\;\Big(\xi\,WV + \mathrm i\,\frac{1 - t^{-2}}{1 + t^{-2}}\;\xi'WV\Big),
\]
where the row in the bracket is holomorphic in $1/t$ near $0$, with values in $\mathbb C_{\varkappa/2}$, and a derivative of order $i$ of $\mathcal B$ is homogeneous of degree $-i$. Second, on the rows $u + \mathrm iu' \in \mathbb C_{3\varkappa/4}$ with $\norm{u} = 1$, $\Der\mathcal B$ and $\Der_\theta\mathcal B$ are bounded. So by Cauchy's estimates on the polydisc of \eqref{eq:e0incone}, the coefficient of a product of $i$ factors $\delta_j$ and of at most one factor $\delta'_j$ is at most a fixed power of $t$ times $(\varrho\,t)^{-i}$. With \eqref{eq:e0deltabound}, the series \eqref{eq:e0expansion} therefore converges absolutely for large $t$, and the sum of any set of its terms is at most a power of $t$ times $e^{-\frac{\dm}{8}(t - t^{-1})^2r}$, where $r$ is the smallest sum of the numbers $r_j$ of the factors of a term of the set.

\emph{Step 1: the term without $\delta$ gives $\eta'$.} By \eqref{eq:e0expansion}, $\mathcal L_0(t)$ tends to $0$ faster than every power of $1/t$. It is meromorphic at infinity, so $\mathcal L_0(t) = 0$ for large $t$, and then for all $t \ge 1$, because it is holomorphic near $[1, \infty)$. At $t = 1$, where $n(t)$ is the row $n = \sqrt{\dm}\;\xi$, this reads, by \eqref{eq:e0L}, \eqref{eq:e0zdot} and \eqref{eq:e0radial},
\[
0 \;=\; \mathcal L\big(nW, 1; n\big) \;=\; \Der\mathcal B(n\,WV)\,\Big(n\,\big(Z\,WV + \eta_WV + W\eta_V\big)\Big) + \Der_\theta\mathcal B(n\,WV)\,\eta' .
\]
The right side is real-analytic and homogeneous of degree $0$ in the row $n \neq 0$, and it vanishes for $\xi$ in the open set $\mathcal P$, hence for all $n \neq 0$. With $h = n\,WV$, which takes every value by (P1), this is \eqref{eq:e0D} for $k+1$ with
\begin{equation}\label{eq:e0Zprime}
Z' = (WV)^{-1}\,\big(Z\,WV + \eta_WV + W\eta_V\big)
\end{equation}
and $\eta'$. By $(D_{k+1})$, $Z' = Y_{k+1} + \beta' I$ and $\eta'$ is as in \eqref{eq:e0Dconcl}, for antisymmetric $Y_{k+1}, \dots, Y_L$ and numbers $\beta', \beta_{k+2}, \dots, \beta_L$.

\emph{Step 2: the slowest hidden unit gives $Z$.} On an open dense set of pairs $(\xi, \xi')$ the numbers $r_1, \dots, r_q$ are pairwise different: $r_i = r_j$ on an open set of pairs would make $c_ic_i^\top - c_jc_j^\top$ a multiple of the identity, which is impossible for $c_i \neq \pm c_j$, and $c_i \neq c_j$ by (P1), $c_i \neq -c_j$ by (P2). Take a hidden unit $j$ and a nonempty open set of pairs, with $\xi \in \mathcal P$, on which $r_j < r_i$ for all $i \neq j$, and let $\gamma = \min\big(r_j, \min_{i\neq j}(r_i - r_j)\big)$. By Step~1, all terms of \eqref{eq:e0expansion} other than the two of unit $j$ are, together, at most a power of $t$ times $e^{-\frac{\dm}{8}(t - t^{-1})^2(r_j + \gamma)}$. Dividing by $\varphi(z_j)$, whose absolute value is a constant times $e^{-\frac{\dm}{8}(t - t^{-1})^2r_j}$ by \eqref{eq:e0zfacts},
\[
\big(z_j\,\mathcal L_j + \mathcal L'_j\big)\,\operatorname{mr}(z_j) \;-\; z_j\,\mathcal L'_j \;\longrightarrow\; 0 \qquad\text{faster than every power of } 1/t .
\]
Here $z_j(t) = \frac{\sqrt{\dm}}{2}\big(t\,\nu c_j + \frac1t\,\bar\nu c_j\big)$ with $\operatorname{Re}(\nu c_j) = \xi c_j > 0$, and $z_j\mathcal L_j + \mathcal L'_j$ and $z_j\mathcal L'_j$ are meromorphic at infinity. By \cref{lem:e0nomero} both vanish, so
\begin{equation}\label{eq:e0unitj}
\mathcal L'_j(t) = 0 \quad\text{and}\quad \mathcal L_j(t) = 0 \qquad\text{for large } t, \text{ hence for all } t \ge 1 .
\end{equation}
By \eqref{eq:e0L}, with $\dot z_j$ the entry $j$ of $\dot z$,
\[
\mathcal L'_j(t) \;=\; \dot z_j\big(n(t)\big)\;\Der\mathcal B\big(n(t)\,WV\big)\,a_j .
\]
At $t = 1$ this is $\dot z_j(n)\,\Der\mathcal B(n\,WV)\,a_j = 0$ at $n = \sqrt{\dm}\,\xi$, for $\xi$ in an open set. By \cref{lem:e0tailrank}, $\Der\mathcal B(n\,WV)\,a_j \neq 0$ on an open dense set of $n$. So $\dot z_j(n) = 0$ on an open subset of $\mathbb S$, that is, by \eqref{eq:e0zdot},
\[
\norm{n}^2\;n\,\big(Zc_j + \eta_{c_j}\big) \;=\; \langle nZ, n\rangle\;nc_j
\]
on an open cone of rows $n$, hence for all $n$, where $\eta_{c_j}$ is column $j$ of $\eta_W$. For $n$ orthogonal to $c_j$ this gives $n\,(Zc_j + \eta_{c_j}) = 0$, so $Zc_j + \eta_{c_j} = \beta_kc_j$ for a number $\beta_k$, and then $\langle nZ, n\rangle = \beta_k\norm{n}^2$ for all $n$. So $Z = Y_k + \beta_kI$ with $Y_k$ antisymmetric.

\emph{Step 3: subtracting a symmetry.} Let $\beta_{k+1} = \beta' - \beta_k$. The pair $Y_k + \beta_kI$ and $\big(-Y_kW,\ V(Y_{k+1} + \beta_{k+1}I),\ \eta'\big)$ is of the form \eqref{eq:e0Dconcl}, so it satisfies \eqref{eq:e0D} for $k$. Subtracting it from $Z$ and $\eta$ leaves the pair $0$ and $(\eta^\circ_W, \eta^\circ_V, 0)$ with
\[
\eta^\circ_W = \eta_W + Y_kW, \qquad \eta^\circ_V = \eta_V - V\,(Y_{k+1} + \beta_{k+1}I),
\]
which satisfies \eqref{eq:e0D} for $k$ as well, and by \eqref{eq:e0Zprime} and Step~1
\begin{equation}\label{eq:e0balance}
\eta^\circ_WV + W\eta^\circ_V \;=\; WV\,(Y_{k+1} + \beta'I) - \beta_k\,WV - WV\,(Y_{k+1} + \beta_{k+1}I) \;=\; 0 .
\end{equation}
For this pair, \eqref{eq:e0L} is $\mathcal L^\circ(\tilde w, \tilde w'; n) = \Der\mathcal B(\tilde wV)\,\big((\tilde w'\odot n\eta^\circ_W)\,V + \tilde w\,\eta^\circ_V\big)$. It remains to show $\eta^\circ_W = 0$ and $\eta^\circ_V = 0$.

\emph{Step 4: the hidden units in the order of their rates.} Take a nonempty open set of pairs $(\xi, \xi')$, with $\xi \in \mathcal P$, on which $r_{j_1} < r_{j_2} < \dots < r_{j_q}$ for one order $j_1, \dots, j_q$ of the hidden units. We show by induction on $i$ that column $j_i$ of $\eta^\circ_W$ and row $j_i$ of $\eta^\circ_V$ vanish. Let this hold for the units $j_1, \dots, j_{i-1}$, and let $j = j_i$. If the rows $\delta$ and $\delta'$ vanish outside the entries $j_1, \dots, j_{i-1}$, then $(1 - \delta')\odot n\eta^\circ_W = n\eta^\circ_W$ and $(z - \delta)\,\eta^\circ_V = nW\eta^\circ_V$, so by \eqref{eq:e0balance}
\[
\mathcal L^\circ\big(z - \delta,\ 1 - \delta';\ n\big) \;=\; \Der\mathcal B\big((z - \delta)V\big)\,\Big(n\,\big(\eta^\circ_WV + W\eta^\circ_V\big)\Big) \;=\; 0 .
\]
So in the expansion \eqref{eq:e0expansion} of $\mathcal L^\circ$ every term has a factor $\delta_l(t)$ or $\delta'_l(t)$ of a unit $l$ other than $j_1, \dots, j_{i-1}$. Other than the two terms of unit $j$, such a term has a factor of a unit $j_{i'}$ with $i' > i$, or a factor of unit $j$ and a second factor. With $\gamma = \min(r_{j_1}, r_{j_{i+1}} - r_j)$, or $\gamma = r_{j_1}$ if $i = q$, we get \eqref{eq:e0unitj} as in Step~2: $\mathcal L^{\circ\prime}_j(t) = 0$ and $\mathcal L^\circ_j(t) = 0$ for $t \ge 1$. Here, by \eqref{eq:e0balance},
\[
\mathcal L^{\circ\prime}_j(t) = \big(n\,\eta^\circ_{c_j}\big)\;\Der\mathcal B(n\,WV)\,a_j, \qquad \mathcal L^\circ_j(t) = \Der\mathcal B(n\,WV)\,\eta^\circ_{a_j} \qquad\text{at } n = n(t),
\]
with $\eta^\circ_{c_j}$ column $j$ of $\eta^\circ_W$ and $\eta^\circ_{a_j}$ row $j$ of $\eta^\circ_V$. At $t = 1$, for $\xi$ in an open set and $n = \sqrt{\dm}\,\xi$: by \cref{lem:e0tailrank}, $n\,\eta^\circ_{c_j} = 0$, so $\eta^\circ_{c_j} = 0$, and $\eta^\circ_{a_j}$ is a multiple of $n\,WV$ for an open set of rows $n$, so $\eta^\circ_{a_j} = 0$.

After $i = q$: $\eta_W = -Y_kW$ and $\eta_V = V\,(Y_{k+1} + \beta_{k+1}I)$. With Step~1 for $\eta'$ and Step~2 for $Z$ this is \eqref{eq:e0Dconcl} for $k$.
\end{proof}

\para{Parameters that satisfy all conditions, in the case (0).} We choose the blocks from the last to the first. Let $D$ and the blocks $\ell+1, \dots, L$ be chosen, regular and with (S1) to (S3) for $\tanh$, or (P1) and (P2) for GELU. This fixes $\varkappa_\ell$ (\cref{lem:e0cone}). We choose $(W, V) = (\Wread^{\ell}, \Wwrite^{\ell})$.

\emph{Conditions that fail on a set of measure zero of pairs $(W, V)$.} (R1), (S1), (S2) and (P1) are of this kind; (S2) fails on countably many proper subspaces. (R2): the triples $(n, W, V)$ with $n \in \mathbb S$, $\sigma(nW) \neq 0$ and $\sigma(nW)\,V = 0$ form a manifold of dimension one less than that of the pairs $(W, V)$, so their projection to the pairs has measure zero; and $\sigma(nW) = 0$ only if $nW = 0$, which requires $W$ of rank less than $\dm$. (R3): the determinant in the proof of \cref{lem:e0tailrank} is real-analytic in $(n, W, V)$, so (R3) fails on a set of measure zero once it holds for one pair. It holds for $W = [\,I \mid 0\,]$ and $V$ with the rows of $I$ and zero rows, at any $n \in \mathbb S$ with positive entries $n_i$: there $g_\ell(n) = \sigma(n)$, and $\Der g_\ell(n)\,v = \mu\,g_\ell(n)$ gives $v_i = \mu\,\sigma(n_i)/\sigma'(n_i)$, so $0 = \langle v, n\rangle = \mu\sum_i n_i\,\sigma(n_i)/\sigma'(n_i)$ with a positive sum, hence $\mu = 0$ and $v = 0$.

\emph{Conditions that hold on a nonempty open set of pairs.} (S3): take any $V$ and $W^\circ$, with the columns $c^\circ_j$, and a unit row $\xi_0$ with $\xi_0c^\circ_j \neq 0$ and $\bar sV \neq 0$ for the row $\bar s$ of the signs. As $\bar sV$ is a real nonzero row, there is a $\rho$ as in (S3). Then $W = \lambda W^\circ$ satisfies (S3) for large $\lambda > 0$, with strict inequalities and a compact set of rows $\tilde wV$ inside the open cone, so (S3) holds near $(W, V)$ as well. (P2): take a row $\alpha$ with positive entries, any $V$ with $\alpha V \neq 0$, and $W^\circ = \xi_0^\top\alpha$. For $W^\circ$ the left sides of the first two inequalities of (P2) vanish and $\xi_0c^\circ_j = \alpha_j > 0$, so these two inequalities hold for all $W = \lambda\,(W^\circ + W')$ with $\lambda > 0$ and $W'$ small. For fixed small $W'$, the third holds for large $\lambda$: its left side decreases like $e^{-\kappa\lambda^2}$ with a constant $\kappa > 0$, and its right side grows like $\lambda$. All three are strict, so (P2) holds on a nonempty open set of pairs.

A nonempty open set minus a set of measure zero is not empty. So there are parameters $\theta^*$ that are regular and whose blocks all satisfy (S1) to (S3) for $\tanh$, or (P1) and (P2) for GELU.

\begin{lemma}[The domain]\label{lem:e0domain}
In the case (0), let $X$ be a finite batch of nonzero inputs, and let $\mathcal O_X \subset \R^P$ be the set of parameters at which the hidden states $h_0, \dots, h_L$ of the skipless network are nonzero for every $x \in X$. Then $\mathcal O_X$ is open and connected, its complement has measure zero, and the outputs $F_X$ on the batch are real-analytic on $\mathcal O_X$.
\end{lemma}

\begin{proof}
This is \cref{lem:domain} for skipless networks. Let $\mathcal O^0$ be the set of parameters with $xE \neq 0$ for all $x \in X$, and for $\ell \ge 1$ let $\mathcal O^\ell$ be the set of parameters in $\mathcal O^{\ell-1}$ with $h_\ell \neq 0$ for all $x \in X$. The complement of $\mathcal O^0$ is a finite union of subspaces of codimension $\dm \ge 2$. Let $\ell \ge 1$ and $x \in X$. On $\mathcal O^{\ell-1}$, $h_\ell = \tilde w_\ell\,\Wwrite^{\ell}$ with $\tilde w_\ell = \sigma\big(N_\varepsilon(h_{\ell-1})\,\Wread^{\ell}\big)$, and $h_{\ell-1}$ does not depend on the weights of block $\ell$. For $\tanh$ and GELU, $\sigma(z) = 0$ only at $z = 0$. So $\tilde w_\ell = 0$ exactly where $N_\varepsilon(h_{\ell-1})\,\Wread^{\ell} = 0$, a closed submanifold of $\mathcal O^{\ell-1}$ of codimension at least $\dm \ge 2$, because the derivative with respect to $\Wread^{\ell}$ is onto. Where $\tilde w_\ell \neq 0$, the derivative of $h_\ell$ with respect to $\Wwrite^{\ell}$ is onto $\R^{1\times\dm}$, so there the parameters with $h_\ell = 0$ form a closed submanifold of codimension $\dm \ge 2$. Removing finitely many closed submanifolds of codimension at least two, one after the other, from a connected open set leaves a connected open set and removes a set of measure zero.
\end{proof}

\subsection{The parts together}\label{app:proofmain}

\begin{proof}[Proof of \cref{thm:skiplessbound}]
\emph{$\tanh$ blocks in the cases (n) and (+).} Let $\theta$ satisfy (R1), let every block satisfy (T1) and (T2), and let every $\Wread^{\ell}\Wwrite^{\ell}$ be invertible. These conditions fail only on a set of measure zero; (T2) fails on countably many proper linear subspaces. Such $\theta$ are regular: $\Der g_\ell(0)$ is the map $v \mapsto v\,\Wread^{\ell}\Wwrite^{\ell}$. By \cref{lem:e0readout} and \cref{lem:sbsteptanh}, Claim $(D_k)$ holds at $\theta$ for $k = L$, then for $k = L-1$, and so on down to $k = 0$. By \cref{prop:e0givesrank}, the Jacobian at $\theta$ has the rank \eqref{eq:e0upper} on some finite batch, and by (R1) it has at most this rank on every batch. So the functional dimension is \eqref{eq:e0upper} at almost every $\theta$.

\emph{GELU blocks in the case (+), and gated blocks in the cases (n) and (+).} By \cref{lem:gtstar}, $\theta^\dagger$ is good, so it is regular, and no value column vanishes. $(D_L)$ holds at $\theta^\dagger$ (\cref{lem:e0readout}), and by \cref{lem:gtstep}, applied for $k = L-1$, then for $k = L-2$, and so on down to $k = 0$, so does $(D_0)$. By \cref{prop:e0givesrank}, the Jacobian at $\theta^\dagger$ has the rank \eqref{eq:e0upper} on some finite batch $X$. In the cases (n) and (+) the outputs on $X$ are real-analytic in $\theta \in \R^P$, and so are the minors of the Jacobian. A minor of the size \eqref{eq:e0upper} does not vanish at $\theta^\dagger$, so it vanishes only on a set of measure zero. At generic parameters the Jacobian on $X$ therefore has at least the rank \eqref{eq:e0upper}, and it has at most this rank on every batch, because $E$ and every $\Wwrite^{\ell}$ have rank $\dm$ and, in gated blocks, no value column vanishes. So the functional dimension is \eqref{eq:e0upper}.

\emph{The case (0).} By \cref{lem:e0readout} and \cref{lem:e0steptanh} or \cref{lem:e0stepgelu}, Claim $(D_k)$ holds at $\theta^*$ for $k = L$, then for $k = L-1$, and so on down to $k = 0$. By \cref{prop:e0givesrank} there is a finite batch $X^*$, with $xE \neq 0$ for all its inputs, on which the Jacobian at $\theta^*$ has the rank \eqref{eq:e0upper}. By (R2), $\theta^* \in \mathcal O_{X^*}$. The sum of the squares of the minors of this size of the Jacobian of $F_{X^*}$ is real-analytic on the connected open set $\mathcal O_{X^*}$ and does not vanish at $\theta^*$, so its zero set has measure zero, and so has the complement of $\mathcal O_{X^*}$ (\cref{lem:e0domain}). Hence the rank on $X^*$ is at least \eqref{eq:e0upper} at almost every parameter. Where $E$ and every $\Wwrite^{\ell}$ have rank $\dm$, it is at most \eqref{eq:e0upper} on every batch. So at generic parameters the functional dimension is \eqref{eq:e0upper}.
\end{proof}

\section{Residual networks under RMSNorm with \texorpdfstring{$\varepsilon = 0$}{no regularization}}\label{app:epszero}

In this appendix a block may be an attention block. As in [KM, Section~\KMn{sec:setup}], a hidden state is then a matrix $h \in \R^{T\times\dm}$ with one row per token, block $\ell$ maps it to $f_\ell(h\,\Wread^{\ell})\,\Wwrite^{\ell}$ with a fixed map $f_\ell$, and $N_\varepsilon$ acts on every row. We recall the theorem.

\thmEpsZero*

The proof has two parts. \Cref{app:sphere} determines the directions of the weights of one block that do not move its output on the sphere. \Cref{app:smallwrites0} brings the network to one such block, by the argument of small writes of [KM], and proves the theorem.

\subsection{A block whose inputs lie on a sphere}\label{app:sphere}

With $\varepsilon = 0$ a block reads only rows of the sphere
\[
\mathbb S \;=\; \big\{\,n \in \R^{1\times\dm} : \norm{n}^2 = \dm\,\big\}.
\]
The lemma of this subsection determines the directions of the weights of a block along which its output, at inputs on the sphere, stays fixed or moves along its input. \Cref{lem:smallwrites0} uses it at any depth.

\para{The block.} Let the block $g$ hold any number of attention heads, of width $d_h \le \dm$, causal or not, on sequences of $T \ge 2$ tokens, and any number of hidden units, which are all $\tanh$, all GELU, all SwiGLU or all GeGLU units. Head $i$ has the projections $W_{Q,i}$, $W_{K,i}$, $W_{V,i}$ and $W_{O,i}$, and, as in the proof of \KMp{Theorem}{thm:onelayer}{i} with $E$ and $D$ replaced by the identity,
\[
B_i = W_{Q,i}W_{K,i}^\top/\sqrt{d_h}, \qquad C_i = W_{V,i}W_{O,i}, \qquad S = \textstyle\sum_i C_i .
\]
Hidden unit $j$ has the column $c_j$ of the read and the row $a_j$ of the write, and a gated unit also the value column $e_j$ of the read. At a row $n$ the unit computes $\sigma(nc_j)\,a_j$, and a gated unit $\sigma(nc_j)\,(ne_j)\,a_j$, where $\sigma$ is SiLU for SwiGLU and GELU for GeGLU. For a direction $\eta$ of the weights of $g$, let $\psi_{c_j}$, $\psi_{e_j}$ and $\psi_{a_j}$ be its entries at $c_j$, $e_j$ and $a_j$, and let
\[
\psi_{B_i} = \big(\eta_{Q,i}W_{K,i}^\top + W_{Q,i}\eta_{K,i}^\top\big)/\sqrt{d_h}, \qquad \psi_{C_i} = \eta_{V,i}W_{O,i} + W_{V,i}\eta_{O,i}, \qquad \psi_S = \textstyle\sum_i \psi_{C_i}
\]
be the derivatives of $B_i$, $C_i$ and $S$ along $\eta$. For a sequence $n = (n_1, \dots, n_T)$ of rows, $\Der_\theta g(n)\,\eta$ is the derivative of the output $g(n)$ along $\eta$.

\begin{lemma}[Block on the sphere]\label{lem:sphere}
Let $\dm \ge 2$. For generic weights of the block $g$, let a direction $\eta$ of its weights satisfy, for every sequence $n = (n_1, \dots, n_T)$ of rows of $\mathbb S$ and every token $t$,
\begin{equation}\label{eq:radial}
\big(\Der_\theta g(n)\,\eta\big)_t \;=\; \mu_t(n)\,n_t \qquad\text{for a number } \mu_t(n).
\end{equation}
Then $\eta$ is a direction of the symmetries inside the block:\\
for every hidden unit, $\psi_{c_j} = 0$;\\
for $\tanh$ and GELU units, $\psi_{a_j} = 0$;\\
for SwiGLU and GeGLU units, $\psi_{e_j} = -\alpha_je_j$ and $\psi_{a_j} = \alpha_ja_j$ for some number $\alpha_j$;\\
for every head, $\psi_{B_i} = 0$ and $\psi_{C_i} = 0$, that is, $\eta_{Q,i} = W_{Q,i}\Gamma_i$, $\eta_{K,i} = -W_{K,i}\Gamma_i^\top$, $\eta_{V,i} = W_{V,i}\Gamma'_i$ and $\eta_{O,i} = -\Gamma'_iW_{O,i}$ for some $d_h\times d_h$ matrices $\Gamma_i$ and $\Gamma'_i$.\\
In particular $\Der_\theta g(n)\,\eta = 0$ at every input.
\end{lemma}

\begin{proof}
\emph{Notation.} For rows $n, v \in \mathbb C^{1\times\dm}$ write $\langle v, n\rangle = vn^\top$, without complex conjugation, and
\[
\operatorname{pr}_n(v) \;=\; \dm\,v - \langle v, n\rangle\,n .
\]
This is linear in $v$ and polynomial in $n$. If $\langle n, n\rangle = \dm$, then $\operatorname{pr}_n(v) = 0$ exactly when $v$ is a multiple of $n$, so \eqref{eq:radial} says that $\operatorname{pr}_{n_t}$ of its left side vanishes. For orthonormal rows $\xi, \xi' \in \R^{1\times\dm}$, the \emph{great circle} of $\xi$ and $\xi'$ is
\[
n(\vartheta) \;=\; \sqrt{\dm}\,\big(\cos\vartheta\;\xi + \sin\vartheta\;\xi'\big), \qquad \vartheta \in \mathbb C .
\]
It has $\langle n(\vartheta), n(\vartheta)\rangle = \dm$ for all $\vartheta \in \mathbb C$, and $n(\vartheta) \in \mathbb S$ for real $\vartheta$. We use two facts. (a)~A meromorphic function on $\mathbb C$ that vanishes on $\R$ vanishes identically. (b)~For a polynomial map $\Theta$ from $\mathbb C^{1\times\dm}$ to itself, $\vartheta \mapsto \Theta(n(\vartheta))$ is a trigonometric polynomial: for some integer $l_0$, $e^{\mathrm il_0\vartheta}\,\Theta(n(\vartheta))$ is a polynomial in $e^{\mathrm i\vartheta}$. So if it vanishes at infinitely many $\vartheta$ with pairwise different $e^{\mathrm i\vartheta}$, it vanishes for all $\vartheta$.

\emph{Generic weights.} The columns $c_j$ are nonzero and pairwise non-collinear, every $a_j$ and $e_j$ is nonzero, and $a_je_j \neq 0$ if $\dm = 2$. Every $B_i$ and $C_i$ is nonzero, and for two heads $i \neq i'$ the rows $uB_i$ and $uB_{i'}$ are linearly independent for some row $u$. Every $W_{Q,i}$, $W_{K,i}$ and $W_{V,i}$ has full column rank, and every $W_{O,i}$ has full row rank.

\emph{Step A: equal tokens.} Let all tokens equal $n \in \mathbb S$. Head $i$ then gives every token the output $nC_i$, whatever its scores are, so \eqref{eq:radial} reads $\operatorname{pr}_n(J(n)) = 0$ for all $n \in \mathbb S$, where $J(n)$ is the sum of $n\psi_S$ and of the terms of the hidden units. As in \KMeq{eq:dotF}, these terms are $\sigma(nc_j)\,\psi_{a_j} + \sigma'(nc_j)\,(n\psi_{c_j})\,a_j$ for a $\tanh$ or GELU unit, and $\sigma(nc_j)\,nW_j + \sigma'(nc_j)\,(n\psi_{c_j})\,(ne_j)\,a_j$ for a gated unit, with the matrix $W_j = \psi_{e_j}a_j + e_j\psi_{a_j} \in \R^{\dm\times\dm}$. So
\begin{equation}\label{eq:sphereA}
\operatorname{pr}_n\big(J(n)\big) \;=\; \operatorname{pr}_n(n\psi_S) + \sum_j\Big[\sigma(nc_j)\,\Theta_j(n) + \sigma'(nc_j)\,\Theta'_j(n)\Big],
\end{equation}
with the polynomial maps
\begin{align*}
&\tanh,\ \text{GELU:} & \Theta_j(n) &= \operatorname{pr}_n(\psi_{a_j}), & \Theta'_j(n) &= (n\psi_{c_j})\;\operatorname{pr}_n(a_j),\\
&\text{SwiGLU},\ \text{GeGLU:} & \Theta_j(n) &= \operatorname{pr}_n(nW_j), & \Theta'_j(n) &= (n\psi_{c_j})\,(ne_j)\;\operatorname{pr}_n(a_j).
\end{align*}
We show that $\Theta'_j$ and $\Theta_j$ vanish on $\mathbb S$, after saying what this gives.

\emph{What $\Theta'_j = 0$ and $\Theta_j = 0$ on $\mathbb S$ give.} Let $\Theta'_j$ vanish on $\mathbb S$. For $n \in \mathbb S$ that is not a multiple of $a_j$, $\operatorname{pr}_n(a_j) \neq 0$, so $n\psi_{c_j} = 0$, or $(n\psi_{c_j})(ne_j) = 0$ for a gated unit. As $e_j \neq 0$, in both cases $n\psi_{c_j} = 0$ on a dense subset of $\mathbb S$, so $\psi_{c_j} = 0$, and $\Theta'_j = 0$ as a polynomial map. Let $\Theta_j$ vanish on $\mathbb S$ as well. For a $\tanh$ or GELU unit, $\psi_{a_j}$ is then a multiple of every $n \in \mathbb S$, so $\psi_{a_j} = 0$. For a gated unit, every $n \in \mathbb S$ is a left eigenvector of $W_j$, so $W_j = \lambda_jI$ for a number $\lambda_j$. Let $v \neq 0$ be a row with $ve_j = 0$. Then $\lambda_jv = vW_j = (v\psi_{e_j})\,a_j$. If $\lambda_j \neq 0$, every such $v$ is a multiple of $a_j$, which is impossible for $\dm \ge 3$ and gives $a_je_j = 0$ for $\dm = 2$. So $W_j = 0$, that is, $\psi_{e_j}a_j = -e_j\psi_{a_j}$. Multiplying on the right by $a_j^\top$ gives $\psi_{e_j} = -\alpha_je_j$ with $\alpha_j = \psi_{a_j}a_j^\top/\norm{a_j}^2$, and then $e_j\,(\psi_{a_j} - \alpha_ja_j) = 0$, so $\psi_{a_j} = \alpha_ja_j$. In both cases the terms of unit $j$ in $J(n)$ vanish at every $n$.

\emph{$\tanh$ and SwiGLU units: poles.} Here $\sigma$ is $\tanh$ or SiLU. Its poles are the points $z_m = \mathrm i\,y_m$, $m \in \mathbb Z$, with $y_m = \pi(m + \tfrac12)$ for $\tanh$ and $y_m = \pi(2m+1)$ for SiLU, and as $\zeta \to 0$
\[
\sigma(z_m + \zeta) = \frac{\varrho_m}{\zeta} + O(1), \qquad \sigma'(z_m + \zeta) = -\frac{\varrho_m}{\zeta^2} + O(1),
\]
with $\varrho_m = 1$ for $\tanh$ and $\varrho_m = z_m$ for SiLU. Fix a unit $j$ and let $\xi = c_j^\top/\norm{c_j}$. Call a unit row $\xi'$ \emph{admissible} if $\xi'c_j = 0$ and $\xi'c_k \neq 0$ for every unit $k \neq j$. The admissible rows are dense among the unit rows orthogonal to $c_j$, because no $c_k$ with $k \neq j$ is collinear with $c_j$. On the great circle of $\xi$ and an admissible $\xi'$,
\[
n(\vartheta)\,c_j \;=\; \sqrt{\dm}\,\norm{c_j}\cos\vartheta .
\]
For $m \in \mathbb Z$ let $\chi_m \in \R$ be the solution of $\sqrt{\dm}\,\norm{c_j}\sinh\chi_m = y_m$, and let $\vartheta_m = \tfrac{\pi}{2} - \mathrm i\chi_m$. Then $\cos\vartheta_m = \mathrm i\sinh\chi_m$ and $\sin\vartheta_m = \cosh\chi_m$, so
\[
n(\vartheta_m) \;=\; \sqrt{\dm}\,\big(\mathrm i\sinh\chi_m\;\xi + \cosh\chi_m\;\xi'\big), \qquad n(\vartheta_m)\,c_j = z_m .
\]
For $k \neq j$ the number $n(\vartheta_m)\,c_k$ has the real part $\sqrt{\dm}\cosh\chi_m\,(\xi'c_k) \neq 0$, so it is not a pole of $\sigma$, whose poles are imaginary. The function $\zeta(\vartheta) = n(\vartheta)\,c_j - z_m$ has a simple zero at $\vartheta_m$, as $\zeta'(\vartheta_m) = -\sqrt{\dm}\,\norm{c_j}\cosh\chi_m \neq 0$. So by \eqref{eq:sphereA}, near $\vartheta_m$,
\begin{equation}\label{eq:spherepole}
\operatorname{pr}_{n(\vartheta)}\big(J(n(\vartheta))\big) \;=\; -\frac{\varrho_m}{\zeta(\vartheta)^2}\,\Theta'_j\big(n(\vartheta)\big) + \frac{\varrho_m}{\zeta(\vartheta)}\,\Theta_j\big(n(\vartheta)\big) + O(1).
\end{equation}
The left side is meromorphic in $\vartheta$ and vanishes for real $\vartheta$, so it vanishes identically, by (a). The coefficient of $(\vartheta - \vartheta_m)^{-2}$ in \eqref{eq:spherepole} is $-\varrho_m\,\Theta'_j(n(\vartheta_m))/\zeta'(\vartheta_m)^2$, so $\Theta'_j(n(\vartheta_m)) = 0$ for every $m \in \mathbb Z$. The numbers $e^{\mathrm i\vartheta_m} = \mathrm i\,e^{\chi_m}$ are pairwise different, so by (b) $\Theta'_j(n(\vartheta)) = 0$ for all $\vartheta$, in particular on the real great circle of $\xi$ and $\xi'$. Every $n \in \mathbb S$ that is not a multiple of $\xi$ lies on the great circle of $\xi$ and of the unit row along $n - \langle n, \xi\rangle\,\xi$. So the real great circles of $\xi$ and of the admissible $\xi'$ cover a dense subset of $\mathbb S$, and $\Theta'_j = 0$ on $\mathbb S$. Hence $\psi_{c_j} = 0$ and $\Theta'_j = 0$ as a polynomial map. The coefficient of $(\vartheta - \vartheta_m)^{-1}$ in \eqref{eq:spherepole} is then $\varrho_m\,\Theta_j(n(\vartheta_m))/\zeta'(\vartheta_m)$, so $\Theta_j(n(\vartheta_m)) = 0$ for every $m$, and in the same way $\Theta_j = 0$ on $\mathbb S$.

\emph{GELU and GeGLU units: the error function.} Here $\sigma(z) = \tfrac{z}{2}\big(1 + \operatorname{erf}\tfrac{z}{\sqrt2}\big)$ and $\sigma'(z) = \tfrac12\big(1 + \operatorname{erf}\tfrac{z}{\sqrt2}\big) + \tfrac{z}{\sqrt{2\pi}}\,e^{-z^2/2}$. On a great circle let $z_j(\vartheta) = n(\vartheta)\,c_j$. At $n = n(\vartheta)$, \eqref{eq:sphereA} is
\begin{multline}\label{eq:sphereerf}
\operatorname{pr}_n\big(J(n)\big) \;=\; \Big[\operatorname{pr}_n(n\psi_S) + \tfrac12\sum_j\big(z_j\Theta_j + \Theta'_j\big)\Big] \\ + \sum_j \tfrac12\big(z_j\Theta_j + \Theta'_j\big)\,\operatorname{erf}\frac{z_j}{\sqrt2} \;+\; \sum_j \frac{z_j}{\sqrt{2\pi}}\,\Theta'_j\;e^{-z_j^2/2},
\end{multline}
where the bracket, $z_j\Theta_j + \Theta'_j$ and $z_j\Theta'_j$ are trigonometric polynomials in $\vartheta$. We need the following two lemmas.
\begin{nested}
\begin{lemma}[Dawson's integral]\label{lem:dawson}
For $y > 0$ let $\operatorname{daw}(y) = e^{-y^2}\int_0^y e^{\varsigma^2}\,d\varsigma$. For every integer $m \ge 1$, as $y \to \infty$,
\[
\operatorname{daw}(y) \;=\; \sum_{l=0}^{m-1}\frac{(2l-1)!!}{2^{\,l+1}}\,\frac{1}{y^{2l+1}} \;+\; O\Big(\frac{1}{y^{2m+1}}\Big), \qquad\text{where } (-1)!! = 1 .
\]
\end{lemma}

\begin{proof}
For $l \ge 0$ let $\mathcal I_l(y) = \int_1^y e^{\varsigma^2}\varsigma^{-2l}\,d\varsigma$. An integration by parts, with $e^{\varsigma^2} = \tfrac{1}{2\varsigma}\tfrac{d}{d\varsigma}e^{\varsigma^2}$, gives
\[
\mathcal I_l(y) \;=\; \frac{e^{y^2}}{2\,y^{2l+1}} + \frac{2l+1}{2}\,\mathcal I_{l+1}(y) + O(1),
\]
and applying this for $l = 0, \dots, m-1$,
\[
\mathcal I_0(y) \;=\; \sum_{l=0}^{m-1}\frac{(2l-1)!!}{2^{\,l+1}}\,\frac{e^{y^2}}{y^{2l+1}} + \frac{(2m-1)!!}{2^{\,m}}\,\mathcal I_m(y) + O(1).
\]
For $y \ge 2$,
\[
\mathcal I_m(y) \;\le\; \int_1^{y/2} e^{\varsigma^2}\,d\varsigma + \Big(\frac{2}{y}\Big)^{2m+1}\int_{y/2}^{y}\varsigma\,e^{\varsigma^2}\,d\varsigma \;\le\; \frac{y}{2}\,e^{y^2/4} + \frac{2^{2m}}{y^{2m+1}}\,e^{y^2} \;=\; O\Big(\frac{e^{y^2}}{y^{2m+1}}\Big).
\]
As $\int_0^y e^{\varsigma^2}\,d\varsigma = \mathcal I_0(y) + O(1)$, multiplying by $e^{-y^2}$ gives the claim.
\end{proof}

\begin{lemma}[The error function on a circle]\label{lem:erfcircle}
Let $r_j > 0$ and $\delta_j \in \R$, for $j$ in a finite set, be such that the points $r_j(\cos\delta_j, \sin\delta_j)$ of $\R^2$ are pairwise different and no two of them are opposite. Let $z_j(\vartheta) = r_j\cos(\vartheta - \delta_j)$, and let $p_0$, $p_j$ and $q_j$ be trigonometric polynomials with values in $\mathbb C^{1\times\dm}$. If
\begin{equation}\label{eq:erfid}
p_0(\vartheta) + \sum_j\Big[p_j(\vartheta)\,\operatorname{erf}\frac{z_j(\vartheta)}{\sqrt2} + q_j(\vartheta)\,e^{-z_j(\vartheta)^2/2}\Big] \;=\; 0 \qquad\text{for all real } \vartheta,
\end{equation}
then $p_0 = 0$ and $p_j = q_j = 0$ for every $j$.
\end{lemma}

\begin{proof}
Both sides of \eqref{eq:erfid} are entire functions of $\vartheta$, so \eqref{eq:erfid} holds for all $\vartheta \in \mathbb C$. We use induction on the number of indices. Let $j$ be an index with the largest $r_j$, and take $\vartheta = \delta_j + \tfrac{\pi}{2} + \mathrm it$ with $t \to +\infty$. Let $l_0$ bound the degrees of the trigonometric polynomials, which are then $O(e^{l_0t})$ on this line.

\emph{The index $j$.} Here $z_j = -\mathrm i\,r_j\sinh t$. With $y = r_j\sinh t/\sqrt2 > 0$,
\[
e^{-z_j^2/2} = e^{y^2}, \qquad \operatorname{erf}\frac{z_j}{\sqrt2} = \operatorname{erf}(-\mathrm iy) = -\frac{2\mathrm i}{\sqrt\pi}\int_0^y e^{\varsigma^2}\,d\varsigma = -\frac{2\mathrm i}{\sqrt\pi}\;e^{y^2}\operatorname{daw}(y).
\]

\emph{The other indices.} For $k \neq j$ let $\delta_{jk} = \delta_j - \delta_k$. Then $z_k = -r_k\sin(\delta_{jk} + \mathrm it)$ and
\[
\operatorname{Re}\big(z_k^2\big) = r_k^2\big(\sin^2\delta_{jk} - \cos(2\delta_{jk})\sinh^2 t\big), \qquad\text{so}\qquad \big\lvert e^{-z_k^2/2}\big\rvert \;\le\; \exp\Big(\frac{r_k^2\cos(2\delta_{jk})}{r_j^2}\,y^2\Big).
\]
As $r_k \le r_j$, and $\cos(2\delta_{jk}) < 1$ if $r_k = r_j$, there is a number $\varkappa \in [0, 1)$ with $\lvert e^{-z_k^2/2}\rvert \le e^{\varkappa y^2}$ for all $k \neq j$. For every $z \in \mathbb C$, $\lvert\operatorname{erf}z\rvert \le \tfrac{2}{\sqrt\pi}\,\lvert z\rvert\max\big(1, \lvert e^{-z^2}\rvert\big)$, by integrating $e^{-\varsigma^2}$ along the segment from $0$ to $z$. So $\lvert\operatorname{erf}(z_k/\sqrt2)\rvert \le r_k\cosh t\;e^{\varkappa y^2}$.

\emph{Together.} Dividing \eqref{eq:erfid} by $e^{y^2}$ gives, as $t \to \infty$,
\begin{equation}\label{eq:erfsmall}
q_j - \frac{2\mathrm i}{\sqrt\pi}\,\operatorname{daw}(y)\,p_j \;=\; O\big(e^{(l_0+1)t}\,e^{-(1-\varkappa)y^2}\big),
\end{equation}
which tends to zero faster than every power of $e^{-t}$. On the line, every component of $p_j$ and $q_j$ is a finite sum of powers $e^{lt}$ with $\lvert l\rvert \le l_0$. Suppose that a component $p$ of $p_j$ is not zero, and let $q$ be the same component of $q_j$. Then $\lvert p\rvert \ge e^{-(l_0+1)t}$ for large $t$, and \eqref{eq:erfsmall} gives, for every integer $m$,
\[
\operatorname{daw}(y) \;=\; \frac{\sqrt\pi}{2\mathrm i}\,\frac{q}{p} + O\big(y^{-m}\big) \qquad (y \to \infty).
\]
Now $e^{t} = \sqrt2\,y/r_j + \sqrt{2y^2/r_j^2 + 1}$ is $y$ times a convergent power series in $1/y$. So $\tfrac{\sqrt\pi}{2\mathrm i}\,\tfrac{q}{p}$, a rational function of $e^{t}$, equals for large $y$ a convergent series in powers of $1/y$ with finitely many positive powers of $y$. This series is then an asymptotic expansion of $\operatorname{daw}(y)$, and such expansions are unique, so by \cref{lem:dawson} its nonzero coefficients are the numbers $(2l-1)!!/2^{\,l+1}$. They grow faster than every geometric sequence, so the series does not converge, a contradiction. Hence $p_j$ vanishes on the line, and a trigonometric polynomial that vanishes on a line is zero, so $p_j = 0$. Then by \eqref{eq:erfsmall} $q_j$, a finite sum of powers of $e^{t}$ on the line, tends to zero faster than every power of $e^{-t}$, so $q_j = 0$. The identity \eqref{eq:erfid} now holds for the other indices, and by induction all $p_k$ and $q_k$ vanish. Then $p_0 = 0$.
\end{proof}
\end{nested}
Call the great circle of $\xi$ and $\xi'$ \emph{admissible} if the points $(\xi c_j,\ \xi'c_j)$ of $\R^2$, one per hidden unit, are nonzero and no two of them are equal or opposite. On an admissible circle $z_j(\vartheta) = r_j\cos(\vartheta - \delta_j)$ with $r_j(\cos\delta_j, \sin\delta_j) = \sqrt{\dm}\,(\xi c_j, \xi'c_j)$, and the left side of \eqref{eq:sphereerf} vanishes for real $\vartheta$. By \cref{lem:erfcircle}, $z_j\Theta'_j = 0$ and $z_j\Theta_j + \Theta'_j = 0$ on the circle for every $j$, and as $z_j$ is not identically zero, $\Theta'_j(n(\vartheta)) = 0$ and $\Theta_j(n(\vartheta)) = 0$. Every $n \in \mathbb S$ lies on an admissible circle: take $\xi = n/\sqrt{\dm}$; the circle of $\xi$ and $\xi'$ is admissible unless $\xi$ and $\xi'$ are both orthogonal to one of the finitely many nonzero columns $c_j$ and $c_j \pm c_k$, which excludes no $\xi'$ if $\dm = 2$ and a nowhere dense set of $\xi'$ if $\dm \ge 3$. So $\Theta'_j$ and $\Theta_j$ vanish on $\mathbb S$.

\emph{Step B: the second token fixed.} By Step~A every hidden unit is as in the statement, and its terms vanish at every input. Take a sequence whose second token is $u \in \mathbb S$ and whose other tokens all equal $w \in \mathbb S$, and let $v = u - w$. Let $k_i = 1$ if head $i$ is causal and $k_i = T - 1$ otherwise, and
\[
\omega_i \;=\; \tfrac12\big(uB_iv^\top - \log k_i\big).
\]
By \KMeq{eq:twotok}, with $M = 0$ and with $E$ and $D$ replaced by the identity, the block gives the second token the output $uS - \tfrac12vS + \tfrac12\sum_i\tanh(\omega_i)\,vC_i$. Its derivative along $\eta$ is
\[
J_2(u, w) \;=\; \tfrac12(u + w)\,\psi_S + \tfrac12\sum_i\Big[\tanh(\omega_i)\,v\psi_{C_i} + \tfrac12\tanh'(\omega_i)\,\big(u\psi_{B_i}v^\top\big)\,vC_i\Big],
\]
and \eqref{eq:radial} at the second token is $\operatorname{pr}_u(J_2(u, w)) = 0$ for all $u, w \in \mathbb S$. We fix $u$ and vary $w$. With the polynomial maps of $w$
\[
\Theta_{u,i}(w) = \operatorname{pr}_u\big(v\psi_{C_i}\big), \qquad \Theta'_{u,i}(w) = \big(u\psi_{B_i}v^\top\big)\;\operatorname{pr}_u\big(vC_i\big), \qquad v = u - w,
\]
\begin{equation}\label{eq:sphereB}
\operatorname{pr}_u\big(J_2(u, w)\big) \;=\; \tfrac12\operatorname{pr}_u\big((u + w)\,\psi_S\big) + \tfrac12\sum_i\Big[\tanh(\omega_i)\,\Theta_{u,i}(w) + \tfrac12\tanh'(\omega_i)\,\Theta'_{u,i}(w)\Big].
\end{equation}
Let $\mathcal U \subset \mathbb S$ be the set of $u$ such that $uB_i \neq 0$ for every head, $uB_i$ and $uB_{i'}$ are linearly independent for all $i \neq i'$, and no $C_i$ has all its rows in $\R\,u$. It is open and dense in $\mathbb S$.

Fix $u \in \mathcal U$ and a head $i$, and let $\xi = uB_i/\norm{uB_i}$. Call a unit row $\xi'$ \emph{admissible} if $\xi'(uB_i)^\top = 0$ and $\xi'(uB_{i'})^\top \neq 0$ for every head $i' \neq i$; the admissible rows are dense among the unit rows orthogonal to $uB_i$. Let $w(\vartheta) = \sqrt{\dm}\,(\cos\vartheta\;\xi + \sin\vartheta\;\xi')$ be the great circle of $\xi$ and an admissible $\xi'$. On it
\[
\omega_i(\vartheta) \;=\; \tfrac12\big(uB_iu^\top - \log k_i - \sqrt{\dm}\,\norm{uB_i}\cos\vartheta\big).
\]
For $m \in \mathbb Z$ choose $\vartheta_m \in \mathbb C$ with $\omega_i(\vartheta_m) = \mathrm i\pi(m + \tfrac12)$, a pole of $\tanh$, that is, with
\[
\sqrt{\dm}\,\norm{uB_i}\cos\vartheta_m \;=\; uB_iu^\top - \log k_i - \mathrm i\pi(2m+1).
\]
As $\cos\vartheta_m$ is not real, $\sin\vartheta_m \neq 0$, so $\zeta(\vartheta) = \omega_i(\vartheta) - \omega_i(\vartheta_m)$ has a simple zero at $\vartheta_m$. The numbers $\cos\vartheta_m$ are pairwise different, and hence so are the $e^{\mathrm i\vartheta_m}$.

For large $\lvert m\rvert$ no other head has a pole at $\vartheta_m$. Indeed, a pole of a head $i' \neq i$ at $w$ requires $\operatorname{Re}\big(uB_{i'}w^\top\big) = uB_{i'}u^\top - \log k_{i'}$, and
\[
uB_{i'}w(\vartheta_m)^\top \;=\; \sqrt{\dm}\,\Big(\cos\vartheta_m\;uB_{i'}\xi^\top + \sin\vartheta_m\;uB_{i'}\xi'^\top\Big).
\]
Here the real part of $\cos\vartheta_m$ does not depend on $m$, $uB_{i'}\xi'^\top \neq 0$, and $\lvert\operatorname{Re}\sin\vartheta_m\rvert \to \infty$ as $\lvert m\rvert \to \infty$, because $\sin^2\vartheta_m = 1 - \cos^2\vartheta_m$ and the imaginary part of $\cos\vartheta_m$ tends to infinity in absolute value. So there is an $m_0$ such that no head $i' \neq i$ has a pole at $w(\vartheta_m)$ for $\lvert m\rvert \ge m_0$. For these $m$, $\tanh(\omega_i) = 1/\zeta + O(\zeta)$ and $\tanh'(\omega_i) = -1/\zeta^2 + O(1)$ near $\vartheta_m$, so by \eqref{eq:sphereB}
\begin{equation}\label{eq:spherepoleB}
\operatorname{pr}_u\big(J_2(u, w(\vartheta))\big) \;=\; -\frac{1}{4\,\zeta(\vartheta)^2}\,\Theta'_{u,i}\big(w(\vartheta)\big) + \frac{1}{2\,\zeta(\vartheta)}\,\Theta_{u,i}\big(w(\vartheta)\big) + O(1).
\end{equation}
The left side is meromorphic in $\vartheta$ and vanishes for real $\vartheta$, so it vanishes identically, by (a). The coefficient of $(\vartheta - \vartheta_m)^{-2}$ gives $\Theta'_{u,i}(w(\vartheta_m)) = 0$ for $\lvert m\rvert \ge m_0$, so by (b) $\Theta'_{u,i}(w(\vartheta)) = 0$ for all $\vartheta$. As in Step~A, over the admissible $\xi'$ this gives $\Theta'_{u,i} = 0$ on $\mathbb S$.

Now $\Theta'_{u,i}(w)$ is the product of the number $u\psi_{B_i}(u - w)^\top$ and the row $\operatorname{pr}_u\big((u - w)\,C_i\big)$, both affine in $w$. This row does not vanish at all $w \in \mathbb S$: otherwise, adding and subtracting its values at $w$ and $-w$, $\operatorname{pr}_u(uC_i) = 0$ and $\operatorname{pr}_u(wC_i) = 0$ for all $w \in \mathbb S$, hence for all rows $w$, so every row of $C_i$ is a multiple of $u$, which $u \in \mathcal U$ excludes. So its zeros form an affine subspace that does not contain $\mathbb S$, they are nowhere dense in $\mathbb S$, and $u\psi_{B_i}(u - w)^\top = 0$ for all $w \in \mathbb S$. Subtracting this at $w$ and at $-w$ gives $u\psi_{B_i}w^\top = 0$ for all $w \in \mathbb S$, so
\[
u\,\psi_{B_i} = 0, \qquad\text{and}\quad \Theta'_{u,i} = 0 \ \text{ as a polynomial map.}
\]
The coefficient of $(\vartheta - \vartheta_m)^{-1}$ in \eqref{eq:spherepoleB} then gives $\Theta_{u,i}(w(\vartheta_m)) = 0$ for $\lvert m\rvert \ge m_0$, and as before $\Theta_{u,i} = 0$ on $\mathbb S$. Subtracting its values at $w$ and at $-w$ gives $\operatorname{pr}_u(w\psi_{C_i}) = 0$ for all $w \in \mathbb S$, hence for all rows $w$: every row of $\psi_{C_i}$ is a multiple of $u$.

Both conclusions hold for every $u \in \mathcal U$. The rows of $\mathcal U$ span $\R^{1\times\dm}$, so $\psi_{B_i} = 0$. Two non-collinear rows of $\mathcal U$ show that every row of $\psi_{C_i}$ vanishes, so $\psi_{C_i} = 0$. As in Step~1 for attention blocks, with $E$ and $D$ replaced by the identity, $\psi_{B_i} = 0$ and $\psi_{C_i} = 0$ give $\eta_{Q,i} = W_{Q,i}\Gamma_i$, $\eta_{K,i} = -W_{K,i}\Gamma_i^\top$, $\eta_{V,i} = W_{V,i}\Gamma'_i$ and $\eta_{O,i} = -\Gamma'_iW_{O,i}$. So $\eta$ is a direction of the symmetries inside the block, and the derivative of $g$ along it vanishes at every input.
\end{proof}

\subsection{Small writes, and the proof of the theorem}\label{app:smallwrites0}

\begin{lemma}[Small writes, $\varepsilon = 0$]\label{lem:smallwrites0}
Let the networks be those of \cref{thm:epszero}. For generic parameters
\[
\theta \;=\; \Big(E,\ \big\{\Wread^{\ell},\ \Wwrite^{\ell}\big\}_{\ell=1}^{L},\ D\Big)
\]
and a number $\tau$, let
\[
\theta_\tau \;=\; \Big(E,\ \big\{\Wread^{\ell},\ \tau\,\Wwrite^{\ell}\big\}_{\ell=1}^{L},\ D\Big).
\]
For all small $\tau \neq 0$, the residual network has $\operatorname{fd}(\theta_\tau) = P - \dm(\dm-1)/2 - 1 - s$, where the batches hold the inputs whose tokens $x_t$ all have $x_tE \neq 0$.
\end{lemma}

\begin{proof}
This is \KM{Lemma}{lem:smallwrites} in one more case, and we keep the notation of its proof: $J_\tau$ is the Jacobian of the outputs on a finite batch at $\theta_\tau$, $S_\tau$ is the diagonal matrix with the entry $1/\tau$ at the coordinates of the reads and $1$ elsewhere, $J_0 = \lim_{\tau\to0}J_\tau S_\tau$, and a direction $\eta \in \R^P$ is \emph{flat} if $(J_0\eta)(x) = 0$ at every input $x$. The wide block $g$ holds the hidden units and the heads of all blocks at $\theta$, and $\Der_\theta g(n)\,\eta$ is its derivative at the input $n$ along the entries of $\eta$ at its reads and writes. In that proof, under pre-RMSNorm with $\varepsilon > 0$, with $h = xE$ and $n = N_\varepsilon(h)$,
\[
(J_0\eta)(x) \;=\; \Der N_\varepsilon(h)\,\big(x\eta_E + \Der_\theta g(n)\,\eta\big)\,D + N_\varepsilon(h)\,\eta_D .
\]
We determine the flat directions for $\varepsilon = 0$.

With $\varepsilon = 0$, $N_\varepsilon(h) = \sqrt{\dm}\,h/\norm{h}$ on every nonzero row, and the inputs are those whose tokens $x_t$ all have $x_tE \neq 0$. On a finite batch of such inputs, every hidden state at $\theta_\tau$ has nonzero rows for small $\tau$, and the expansion of $y$ and the formula for $(J_0\eta)(x)$ hold as for $\varepsilon > 0$, with
\[
\Der N_\varepsilon(h)\,k \;=\; \frac{\sqrt{\dm}}{\norm{h}}\Big(k - \frac{\langle h, k\rangle}{\norm{h}^2}\,h\Big)
\]
for rows $h \neq 0$ and $k$. This derivative vanishes exactly when $k$ is a multiple of $h$. Let $\eta$ be flat. As in Step~1 of the proof of \KMp{Theorem}{thm:onelayer}{ii}, split $\eta_D = -YD + Z$; then $Z = 0$, because the rows $N_\varepsilon(h)$ span $\R^{\dm}$, and at every token
\[
\Der N_\varepsilon(h)\,\big(x\eta_E + \Der_\theta g(n)\,\eta\big) \;=\; N_\varepsilon(h)\,Y .
\]
The left side is orthogonal to $h$, so $hYh^\top = 0$ for all rows $h$, and $Y$ is antisymmetric. Then $\Der N_\varepsilon(h)\,(hY) = N_\varepsilon(h)\,Y$, so at every token $t$
\[
x_t\eta_E + \big(\Der_\theta g(n)\,\eta\big)_t - x_tEY \;=\; \mu_t(x)\,x_tE \qquad\text{for a number } \mu_t(x).
\]
Replace $x$ by $\lambda x$ with $\lambda > 0$. Then $n$ and $\Der_\theta g(n)\,\eta$ do not change, and
\[
\lambda\,x_t\,(\eta_E - EY) + \big(\Der_\theta g(n)\,\eta\big)_t \;=\; \lambda\,\mu_t(\lambda x)\,x_tE \qquad\text{for all } \lambda > 0 .
\]
With $\lambda = 1$ and $\lambda = 2$, both rows on the left are multiples of $x_tE$:
\[
x_t\,(\eta_E - EY) \in \R\,x_tE, \qquad \big(\Der_\theta g(n)\,\eta\big)_t \in \R\,n_t .
\]
The first relation holds for every row $x_t$ with $x_tE \neq 0$. If $x_0E = 0$, then $x_t\,(\eta_E - EY)$ and $(x_t + x_0)\,(\eta_E - EY)$ are both multiples of $x_tE$, so $x_0\,(\eta_E - EY)$ is a multiple of every nonzero row of $\R^{\dm}$, hence zero. So $\eta_E - EY = E\Lambda$ for a matrix $\Lambda \in \R^{\dm\times\dm}$, and every nonzero row $v$ of $\R^{\dm}$ has $v\Lambda \in \R\,v$. Hence $\Lambda = cI$ for a number $c$, and $\eta_E = E\,(Y + cI)$. The second relation holds for every sequence $n = (n_1, \dots, n_T)$ of rows of the sphere $\mathbb S$ of \cref{lem:sphere}, since $E$ has rank $\dm$. By that lemma, applied to the wide block $g$, the entries of $\eta$ at the reads and the writes are a direction of the symmetries inside the blocks. Conversely, a direction with the entries $E\,(Y + cI)$ at $E$ and $-YD$ at $D$, with $Y$ antisymmetric, is flat, since $\Der N_\varepsilon(h)\,\big(h\,(Y + cI)\big) = N_\varepsilon(h)\,Y$, and so are the directions of the symmetries inside the blocks. So the flat directions form a space of dimension $\dm(\dm-1)/2 + 1 + s$.

Let $m = \dm(\dm-1)/2 + 1 + s$. As in the last step of the proof of \KM{Lemma}{lem:smallwrites}, the kernel of $J_0$ decreases as the batch grows, and its intersection over all batches is the space of flat directions, so on some finite batch $\operatorname{rank}J_0 = P - m$. The rank is lower semicontinuous and $J_\tau S_\tau \to J_0$, so on this batch $\operatorname{rank}J_\tau = \operatorname{rank}(J_\tau S_\tau) \ge P - m$ for all small $\tau \neq 0$. Conversely $\operatorname{rank}J_\tau \le P - m$ on every batch, since the changes of basis in $\CO(\dm)$ of the whole network and the symmetries inside the blocks move $\theta_\tau$ in $m$ independent directions.
\end{proof}

\begin{proof}[Proof of \cref{thm:epszero}]
Let $r = P - \dim\CO(\dm) - s$. By \cref{lem:smallwrites0}, at generic $\theta$ and for small $\tau \neq 0$, the Jacobian at $\theta_\tau$ has rank $r$ on a finite batch $X$ whose tokens $x_t$ all have $x_tE \neq 0$. With $\varepsilon = 0$ the outputs $F_X$ on this batch are defined, and real-analytic in the parameters, where no row of a hidden state vanishes. We need the following lemma.
\begin{nested}
\begin{lemma}[The domain]\label{lem:domain}
Let $X$ be a finite batch of inputs with nonzero tokens, and let $\mathcal O_X \subset \R^P$ be the set of parameters at which every row of every hidden state $h_0, \dots, h_L$ is nonzero, for every $x \in X$. Then $\mathcal O_X$ is open and connected, and its complement has measure zero.
\end{lemma}

\begin{proof}
Let $\mathcal O^0$ be the set of parameters with $x_tE \neq 0$ for all tokens of all $x \in X$, and for $\ell \ge 1$ let $\mathcal O^\ell$ be the set of parameters in $\mathcal O^{\ell-1}$ at which every row of $h_\ell$ is nonzero for every $x \in X$. Then $\mathcal O_X = \mathcal O^L$, and $h_\ell$ is real-analytic on $\mathcal O^{\ell-1}$. The complement of $\mathcal O^0$ is a finite union of linear subspaces of codimension $\dm \ge 2$, so $\mathcal O^0$ is open and connected, and its complement has measure zero. Let $\ell \ge 1$, fix $x \in X$ and a token $t$, and let $f_{\ell,t}$ be row $t$ of $f_\ell\big(N_\varepsilon(h_{\ell-1})\Wread^{\ell}\big)$, so that row $t$ of $h_\ell$ is
\[
h_{\ell,t} \;=\; h_{\ell-1,t} + f_{\ell,t}\,\Wwrite^{\ell} .
\]
Neither $h_{\ell-1,t}$ nor $f_{\ell,t}$ depends on $\Wwrite^{\ell}$. At a point of $\mathcal O^{\ell-1}$ with $h_{\ell,t} = 0$ we have $f_{\ell,t}\,\Wwrite^{\ell} = -h_{\ell-1,t} \neq 0$, so $f_{\ell,t} \neq 0$, and the derivative of $h_{\ell,t}$ with respect to $\Wwrite^{\ell}$, the map $\eta_{\mathrm{write}} \mapsto f_{\ell,t}\,\eta_{\mathrm{write}}$, is onto $\R^{\dm}$. So the parameters in $\mathcal O^{\ell-1}$ with $h_{\ell,t} = 0$ form a closed submanifold of $\mathcal O^{\ell-1}$ of codimension $\dm \ge 2$. Removing finitely many closed submanifolds of codimension at least two from a connected open set leaves a connected open set, and removes a set of measure zero. So every $\mathcal O^\ell$ is open and connected, and its complement has measure zero.
\end{proof}
\end{nested}
For small $\tau$, every hidden state at $\theta_\tau$ is $xE + O(\tau)$, so $\theta_\tau \in \mathcal O_X$. The sum of the squares of the minors of size $r$ of the Jacobian of $F_X$ is real-analytic on the connected open set $\mathcal O_X$ and does not vanish at $\theta_\tau$, so its zero set has measure zero. Hence the rank on $X$ is at least $r$ at almost every parameter. At generic parameters it is at most $r$ on every batch: the changes of basis in $\CO(\dm)$ \KMp{Theorem}{thm:gaugefull}{iii} and the symmetries inside the blocks move the parameters in $\dim\CO(\dm) + s$ independent directions, as in the proof of \KM{Lemma}{lem:smallwrites}. So the functional dimension is $r$ at generic parameters.
\end{proof}

\section{Wide GELU blocks without normalization, and the theorem of Fefferman}\label{app:widegelu}

We recall the proposition.

\propWideGelu*

\begin{proof}
Without normalization $\dim\G = \dm^2$ and $s = 0$, so the bound \eqref{eq:e0upper} is $P - (L+1)\dm^2$, and it remains to find parameters and a batch on which the Jacobian has this rank.

Let block $\ell$ have $q_\ell \ge 2\dm$ hidden units, and take a residual network with blocks of $q_\ell - \dm$ hidden units at the parameters $\theta_\tau$ of \KM{Lemma}{lem:smallwrites}, in which every $\Wwrite^{\ell}$ is replaced by $\tau\,\Wwrite^{\ell}$. Since $\sigma(z) - \sigma(-z) = z$, each of its blocks equals a block without skip with $\dm$ more hidden units \citep{mijoski2026worth}: for $\dm$ hidden units $i$ whose reads $w_i$ form an invertible matrix, add the row $\gamma_i$ of its inverse to the write of unit $i$, and add a hidden unit with the read $-w_i$ and the write $-\gamma_i$. This construction is smooth in the residual parameters near these, where the chosen reads stay invertible, and gives a skipless network that computes the residual one. By the chain rule, the columns of the skipless Jacobian therefore span those of the residual Jacobian, which, with the read columns divided by $\tau$, tend to the limit $J_0$ of the proof of \KM{Lemma}{lem:smallwrites}, of rank $P_R - \dm^2$ on some finite batch, where $P_R = P - 2L\dm^2$ counts the parameters of the residual network. As $\sigma(z) - \sigma(-z) = z$ gives $\sigma'(-z) = 1 - \sigma'(z)$, and every block is the identity at $\tau = 0$, the reads of the added units give the further columns
\[
\big(\sigma'(xEw_i) - 1\big)(xE\delta_i)\,\gamma_iD,\qquad \delta_i \in \R^{\dm},
\]
in the limit. Suppose that $J_0\,\eta$ plus a sum of these vanishes for all $x$. Each added term has the Gaussian part of the term of $\psi_{c_i}$ in \KM{Lemma}{lem:unitid}, with $E\delta_i$ in place of $\psi_{c_i}$ and the row $\gamma_iD$ in place of $a_i$. In the proof of that lemma for GELU, on the rays $x = tn$, the coefficient of $t^2$ for the slowest unit $j$ becomes $u_j\big((n\psi_{c_j})\,a_j + (nE\delta_j)\,\gamma_jD\big)$. At generic parameters $\gamma_jD$ is not parallel to $a_j$, so $nE\delta_j = 0$ for all $n$ in an open set, and $\delta_j = 0$ as $E$ has rank $\dm$; the proof then goes on as before. So, on some finite batch, the added reads raise the rank in the limit by $L\dm^2$, to $P_R - \dm^2 + L\dm^2 = P - (L+1)\dm^2$. By lower semicontinuity the skipless network has this rank for small $\tau$, which is the bound \eqref{eq:e0upper}, and an analytic network keeps its largest rank at generic parameters.
\end{proof}

\begin{remark}[Skipless $\tanh$ networks without normalization, from the theorem of Fefferman]\label{rem:fefferman}
Without normalization, the equality condition for skipless $\tanh$ networks also follows from Theorem~I.B.9 of \citet{fefferman1994}, with or without biases. With read biases $b_\ell$ and write biases $c_\ell$, which may vanish, the skipless network computes
\[
x^0 = x,\qquad x^\ell = \tanh\big(x^{\ell-1}A_\ell + \beta_\ell\big)\ \ (1 \le \ell \le L),\qquad y = x^LA_{L+1} + \beta_{L+1},
\]
with $A_1 = EW_1$, $A_\ell = V_{\ell-1}W_\ell$ and $A_{L+1} = V_LD$, where $W_\ell$ and $V_\ell$ are the read and the write of block $\ell$, and $\beta_1 = b_EW_1 + b_1$, $\beta_\ell = c_{\ell-1}W_\ell + b_\ell$ and $\beta_{L+1} = c_LD + b_D$. Two such networks that agree on an open set of inputs agree on all inputs, by analyticity. As $\sigma(z) = \tanh(z/2)$ is injective, two such networks have the same output exactly when the nets of that paper with the hidden layers $x^\ell = \sigma(2x^{\ell-1}A_\ell + 2\beta_\ell)$ and the output layer $\sigma(x^LA_{L+1} + \beta_{L+1})$ do, and then so do, for every input coordinate and every output, the nets with one input and one output that keep only these. At generic parameters every entry of every $A_\ell$ is nonzero, and no ratio of the weights that one unit sends to two units of a layer of width $\dhid$ is a fraction $p/q$ with $1 \le q \le 100\,\dhid^2$. These conditions are open, and they are those of Theorem~I.B.9 (pages 517 and 518 there), which is stated for nets with one input and one output and puts no condition on the thresholds. So near generic parameters, two networks with the same output give, for every input coordinate and every output, nets with one input and one output that agree up to permutations of the units of each hidden layer and sign flips. These are the identity: at the first layer where a unit is moved or flipped, the weights it receives would move by a fixed amount, since none vanishes and the weights that one unit sends to two units differ in absolute value. So the two networks have the same $A_\ell$ and $\beta_\ell$. When $E$, $D$ and every read and write have rank $\dm$, the parameters with the same $A_\ell$ and $\beta_\ell$ are those reached by the changes of basis at the $L+1$ hidden states, and with biases by the translations there.

By analyticity, the functional dimension at generic parameters is the largest rank $r$ that the Jacobian attains on any batch at any parameters. Take generic parameters and a batch on which the rank is $r$. Near these parameters the rank on this batch is at least $r$, by lower semicontinuity, and so equal to $r$. Its kernel is then the space of directions that are flat for every input, and the parameters with the same outputs on the batch form a manifold of dimension $P - r$ along which the function does not change. So $P - r$ is at most the dimension of the symmetries, and the symmetries give the reverse inequality (\KMp{Theorem}{thm:count}{ii} and \KM{Lemma}{lem:translations}). The proof of \cref{thm:skiplessbound} does not use this theorem.
\end{remark}

\end{document}
